% Intended LaTeX compiler: pdflatex
\documentclass[11pt]{article}
\usepackage[utf8]{inputenc}
\usepackage[T1]{fontenc}
\usepackage{graphicx}
\usepackage{longtable}
\usepackage{wrapfig}
\usepackage{rotating}
\usepackage[normalem]{ulem}
\usepackage{amsmath}
\usepackage{amssymb}
\usepackage{capt-of}
\usepackage{hyperref}
\author{root}
\date{\today}
\title{}
\hypersetup{
 pdfauthor={root},
 pdftitle={},
 pdfkeywords={},
 pdfsubject={},
 pdfcreator={Emacs 29.3 (Org mode 9.6.15)}, 
 pdflang={English}}
\begin{document}

\tableofcontents

\section{org-pwa}
\label{sec:org6509236}
A single-file, offline-capable, mobile-first outliner for editing \texttt{.org} files in the browser. No server, no account required to use it locally — just a static site you can install as a PWA and point at a file.

This document describes what org-pwa actually does today, how to use it, and — since that's usually the more useful question for anyone coming from Emacs — \textbf{where it deliberately or incidentally diverges from real org-mode.}

\noindent\rule{\textwidth}{0.5pt}

\subsection{Contents}
\label{sec:org17d8aa6}
\begin{itemize}
\item \hyperref[sec:org8f2a9cf]{What org-pwa is (and isn't)}
\item \hyperref[sec:org1037811]{Getting started}
\item \hyperref[sec:org22c076e]{Editing your outline}
\item \hyperref[sec:org9a920ce]{Undo/redo and History}
\item \hyperref[sec:orgbe2db61]{Lists, tables, and body text}
\item \hyperref[sec:org1a88f72]{Inline text markup}
\item \hyperref[sec:orgf0231a2]{Searching}
\item \hyperref[sec:org01b1fbc]{Links and images}
\item \hyperref[sec:org4f4f89f]{Configuration}
\item \hyperref[sec:org8f91341]{Progress logging}
\item \hyperref[sec:orgf0de015]{Archiving}
\item \hyperref[sec:org6433005]{Refile}
\item \hyperref[sec:org7b27204]{Attachments}
\item \hyperref[sec:orgb9e46b3]{Clocking}
\item \hyperref[sec:orga9cae38]{The plain-text editor}
\item \hyperref[sec:org105cb3a]{Agenda}
\item \hyperref[sec:orgf6942c0]{The TODO view}
\item \hyperref[sec:org83b5238]{Capture Templates}
\item \hyperref[sec:org08ac3df]{Extras menu}
\item \hyperref[sec:orgfca7ef3]{Menu customization}
\item \hyperref[sec:org5584f85]{File management}
\item \hyperref[sec:orgda9f5ef]{Tabs}
\item \hyperref[sec:org053e4f4]{Export}
\item \hyperref[sec:orgb4d67ea]{Help}
\item \hyperref[sec:orgbb54b74]{Settings}
\item \hyperref[sec:org57d263c]{Offline behavior and sync}
\item \hyperref[sec:orgc3c9c80]{Platform support}
\item \hyperref[sec:org74c50ec]{Security and privacy}
\item \hyperref[sec:org996c4cc]{Keybindings}
\item \hyperref[sec:org2d2d41c]{Differences from Emacs org-mode}
\item \hyperref[sec:org52c5ce1]{Known limitations / not built yet}
\item \hyperref[sec:org45678a6]{Development}
\end{itemize}

\noindent\rule{\textwidth}{0.5pt}

\subsection{What org-pwa is (and isn't)}
\label{sec:org8f2a9cf}
org-pwa is a \textbf{touch-first outline editor}, not a port of Emacs. It reads and writes real \texttt{.org} files — round-trip fidelity is a hard requirement, so content it doesn't specifically understand is preserved verbatim rather than mangled or dropped. But the \emph{editing model} is built from scratch around tapping and gestures, not Emacs keybindings, and it deliberately implements a useful subset of org-mode rather than all of it.

If you're looking for org-mode-in-the-browser with full parity — babel, \LaTeX{} export, real Emacs Lisp — this isn't that (see \hyperref[sec:org52c5ce1]{Known limitations}). If you want to open a real org file on your phone, reorganize your outline, check off tasks, edit tables, and get it saved back to disk, GitHub, or a WebDAV server, that's what this is for.

\noindent\rule{\textwidth}{0.5pt}

\subsection{Getting started}
\label{sec:org1037811}
\begin{enumerate}
\item Open the app. If no file is open, tap \textbf{File → Open}.
\item Choose where to open from:

\begin{itemize}
\item \textbf{Local file} — uses the browser's native file picker and keeps a live handle, so Save writes straight back to the same file. Requires a browser with File System Access support (Chrome/Edge on desktop and Android).
\item \textbf{GitHub} — browse folders and \texttt{.org} files in a repo you've configured in Settings, tap one to open it. A "Type a path instead…" option is always available too, for a path that doesn't show up in the listing (e.g. permissions on a specific subfolder) or if listing fails outright.
\item \textbf{WebDAV} — same idea, browsing a WebDAV server you've configured in Settings via the standard \texttt{PROPFIND} directory-listing method. Same "Type a path instead…" fallback available here too.
\item \textbf{Import file} (shown instead of "Local file" on platforms without File System Access, i.e. everything on iOS) — picks a file once and reads it; there's no live link back to the original, so Save downloads a new copy you place back manually. See \hyperref[sec:orgc3c9c80]{Platform support}.
\end{itemize}

\item Edit. Every edit is applied instantly to an in-memory copy and cached locally in the background — nothing is lost if you close the tab, even before you hit Save.
\item Tap \textbf{Save} when you want it written back to wherever it came from. \textbf{Save As} lets you choose a different destination or backend.
\end{enumerate}

\textbf{Menus and popups}: the File, View, More and Extras menus and every popup open as a bottom sheet on a phone or tablet. On a wide screen with a mouse the menus drop down under their button, and \textbf{↑}, \textbf{↓}, \textbf{Home}, \textbf{End} and \textbf{Enter} move through them. Popups become centered dialogs from 600 px wide, and \textbf{Esc} closes either. Questions such as a file path or a replacement string are asked in the same kind of dialog; \textbf{Enter} accepts.

The modeline (see \hyperref[sec:org6cade45]{Modeline and minibuffer}) shows the current filename, which backend it came from, and whether there are unsaved edits — the same information the tab bar (see \hyperref[sec:orgda9f5ef]{Tabs}) also carries once a second tab is open, via a solid dot on any tab with unsaved changes.

\noindent\rule{\textwidth}{0.5pt}

\subsection{Editing your outline}
\label{sec:org22c076e}
\begin{itemize}
\item \textbf{Headings}: tap the title to reveal twelve actions: edit title, edit the heading's body text, edit details (the general editor, below), add a sub-heading, TODO, move up and down, archive, attach (see \hyperref[sec:org7b27204]{Attachments}), delete, promote and demote. Most actions close the row after running; move up and down and promote and demote leave it open so you can repeat them. Tap elsewhere on the row to close it. \textbf{Delete} always confirms, listing what would be lost (dates, tags, properties, state history). Every delete or discard confirmation in the app (heading, table, note, list item, row or column, attachment, the recent-files list) has \textbf{Cancel} as the default, focused when the dialog opens, so a stray Enter, an interrupted tap or a tap on the backdrop never deletes. The buttons arrange themselves to the width: four across in three rows on a narrow screen (with the arrows last, as left, right, up, down), six across in two rows, or one row of twelve on a wide one. With a keyboard, the arrow keys move between them and \textbf{Esc} closes the row; hovering shows a button's name, and \textbf{Delete} is red.
\item \textbf{TODO state}: tap the TODO badge to cycle the \texttt{\#+TODO:} sequence (default \texttt{TODO}, \texttt{DONE}) one step at a time. The action menu's \textbf{TODO} button (\texttt{C-c C-t}) does the same, except that for a heading with no keyword in a file with more than one parallel \texttt{\#+TODO:} workflow it opens a picker of every state of every workflow (see \hyperref[sec:org8f91341]{Progress logging}).
\item \textbf{Fold and unfold}: tap the chevron to toggle a heading's subtree, or swipe left on a heading to cycle collapsed → children only → fully expanded.
\item \textbf{Move}: drag the grip at the right of a heading (Settings → Quick Settings → \textbf{Drag handle on headings} puts it on the left or turns it off) to move it with its subtree. The blue line snaps to the gap between two headings, and moving sideways steps through the levels that fit there, one step per level, moving right being deeper. From a right-hand grip you start at the deepest level and move left to come out; from a left-hand grip you start at the outermost and move right to go in. A tag names the result (\textbf{into} a heading, \textbf{same level as} one, or \textbf{top level}). Letting go over the top bar, or over the mode line and minibuffer (they turn red and say \textbf{Release to cancel}), or pressing Escape, cancels; the very top and bottom of the outline are still drops. The Up, Down, Promote and Demote buttons in the action menu do the same in single steps.
\item \textbf{Complete or archive}: swipe right. With an active TODO keyword it jumps straight to the sequence's first DONE-type keyword. Without one, including an already-DONE heading, it opens the Archive, Refile or Cancel prompt. An already-archived heading unarchives directly.
\item \textbf{Archive} (\textbf{Unarchive} for an archived heading) moves the subtree to its \texttt{org-archive-location}, or back to where it came from, from a panel that shows the destination first (see \hyperref[sec:orgf0de015]{Archiving}). The \texttt{:ARCHIVE:} tag and \texttt{ARCHIVE\_*} properties are ordinary once the move lands and can be edited or removed in the general editor.
\item \textbf{Moving headings} (↑ ↓ ← →): ↑ and ↓ reorder a heading among its siblings, taking its subtree with it. ← promotes it one level, as a sibling right after its former parent. → demotes it one level, as the last child of the heading before it. At a boundary (first or last, top level, nothing to demote under) a brief status message says so.
\item \textbf{The general editor} (📋 \textbf{Edit details}) is a popup with five tabs, so nothing is a long scroll: \textbf{Dates}, \textbf{Tags}, \textbf{Props}, \textbf{Files} and \textbf{More} (priority and tables). A tab you have edited shows a dot, and a note beside the buttons names them ("Changed: Dates, Tags"). Each tab shows how many items it holds.
\begin{itemize}
\item \textbf{SCHEDULED, DEADLINE and a plain timestamp}: date and time pickers, a repeat selector (every N hours, days, weeks, months or years) and an optional warning delay (\texttt{-3d} shows a deadline three days early). The three are independent, each with its own \textbf{Clear}. See \hyperref[sec:org105cb3a]{Agenda} for how a plain timestamp differs (it doesn't carry forward, and lives in the title).
\item \textbf{Tags}: removable chips (tap to remove) and an input with an \textbf{Add} button (Enter also adds). Suggestions appear as you type: every tag used in the open document or a loaded agenda file (most used first) plus tags declared on a \texttt{\#+TAGS:} line, narrowed to what you typed and excluding tags the heading already has. One tap adds one.
\item \textbf{Priority}: buttons A, B, C and None. Priorities outside A to C aren't reachable here.
\item \textbf{Properties}: an editable key/value row for each, with a remove button and an \textbf{Add property} row. The Key field completes from every key used in the open document or a loaded agenda file (most used first) plus ones the app reads or writes (\texttt{CUSTOM\_ID}, \texttt{EFFORT}, \texttt{Effort\_ALL}, \texttt{ID}, the \texttt{ARCHIVE\_*} set); other keys still work. A row with a blank key is dropped on commit. This is also how you set a \texttt{CUSTOM\_ID} for \texttt{[[\#id]]}-style links (see \hyperref[sec:org01b1fbc]{Links and images}).
\end{itemize}

Three buttons govern every tab: \textbf{Cancel} discards everything and closes, \textbf{Reset} restores every field to how it was when the editor opened, and \textbf{OK} commits every tab and closes. Nothing is written until \textbf{OK}. Above them are three immediate actions. \textbf{Add table} inserts a table into the body and closes the editor. \textbf{Delete table} removes the heading's most recently added table after a confirmation, and is disabled when there is none. \textbf{Attach} closes the editor and opens the attachment picker (see \hyperref[sec:org7b27204]{Attachments}). If any field has changed, these three first ask \textbf{Commit and continue}, \textbf{Discard and continue} or \textbf{Keep editing}; with no changes they act at once.

Existing properties also show as a small read-only line under the heading while drawers are visible for it (\texttt{\#+STARTUP: showeverything}, or manually cycled to fully expanded; see \hyperref[sec:orga77e6fc]{Folding and \#+STARTUP}). Tapping it opens the action menu.

\item \textbf{Block content} (\texttt{\#+BEGIN\_SRC}, \texttt{\#+BEGIN\_QUOTE}, \texttt{\#+BEGIN\_EXAMPLE} and so on) is shown read-only, under the same drawer-visibility rule as properties. Folded, it shows a "collapsed block" label: tap it to reveal the content, and tap it again to re-fold. Edit blocks in \hyperref[sec:orga9cae38]{Text view}. Any \texttt{\#+BEGIN\_NAME} and \texttt{\#+END\_NAME} pair works, and three names are displayed specially:
\begin{itemize}
\item \texttt{QUOTE}: inline markup is interpreted, and the block is indented on both margins with a left border.
\item \texttt{VERSE}: inline markup is interpreted, and every source line stays its own line.
\item \texttt{CENTER}: inline markup is interpreted, and its paragraphs are centered.
\end{itemize}

Everything else (\texttt{SRC}, \texttt{EXAMPLE}, \texttt{COMMENT}, custom names) is shown literally.

On export (see \hyperref[sec:org053e4f4]{Export}), \texttt{SRC} becomes a fenced code block with its language kept (\texttt{\#+BEGIN\_SRC python} gives ```python), \texttt{QUOTE} becomes a blockquote, and \texttt{COMMENT} is dropped. Every other name exports as a plain fenced block, without \texttt{VERSE}'s line preserving or \texttt{CENTER}'s centering. A \texttt{SRC} language is kept only as that export hint: there is no syntax highlighting, and only \texttt{js} blocks can run (see \hyperref[sec:org41ed666]{Source blocks}).
\end{itemize}

\subsubsection{A heading's "text"}
\label{sec:orgcb56bee}
Org has no separate description field: everything after a heading, up to the next heading, belongs to it, and can include several paragraphs, lists and tables in any order. \textbf{Edit text} opens all of it as one block of raw org syntax in a popup, like Settings' \textbf{Menu labels} editor (see \hyperref[sec:orgfca7ef3]{Menu customization}): a large textarea with \textbf{Cancel} (discard and close), \textbf{Reset} (reload the original text and keep the popup open) and \textbf{OK} (save and close). The heading's rendered body is hidden from the outline while the popup is open. Typing \texttt{- [ ]} at the start of a line becomes a checklist item on save, as in the file.

\noindent\rule{\textwidth}{0.5pt}

\subsection{Undo/redo and History}
\label{sec:org9a920ce}
Every edit can be undone and redone: heading, paragraph, list-item and table-cell text, TODO state, priority, tags, properties, moving, promoting, archiving, capturing, all of it. The command palette's \textbf{Undo history} opens a panel with quick \textbf{‹ Undo} and \textbf{Redo ›} buttons for one step at a time, a scrollable list of every step since the file was opened (oldest to newest, current position marked), and a \textbf{diff} toggle on each entry showing what that step changed, as added and removed lines diffed against the step before it.

Tap \textbf{any} entry to jump straight there, forward or back. Entries past your position (available to redo) are dimmer. A genuinely new edit from an undone point discards what was ahead of it, as in most editors, but browsing the list, including jumping back and forward again, discards nothing.

Each step is a full snapshot of the document's text, one per completed edit action (saving a title, toggling a checkbox, archiving a heading) rather than one per keystroke. An edit that changes nothing, such as toggling something and toggling it back, gets no step.

History is unlimited while a file stays open. It also survives closing and reopening the file, capped at the 200 most recent steps, but only when the text on reopen exactly matches what the app last had: if the file changed in the meantime (edited elsewhere, on another device, or directly on GitHub or WebDAV), the old history is discarded. Undo and redo reapply the file's \texttt{\#+STARTUP} visibility instead of preserving fold state across the jump.

Landing on a step whose text is identical to the last confirmed saved version (what the app last read or wrote) clears the filename's modified indicator. That baseline updates on every successful Save during the session.

\textbf{If the filename shows modified (red) and nothing you did explains it}, check the first history entry. If it reads "Opened (resumed unsaved local version)" instead of "Opened", the file had an unsaved edit from an earlier session, and reopening resumed it instead of starting from what is stored. Save to make that version the real one, or jump back in the history list if it wasn't what you wanted.

\noindent\rule{\textwidth}{0.5pt}

\subsection{Lists, tables, and body text}
\label{sec:orgbe2db61}
\begin{itemize}
\item \textbf{Lists}: ordered, unordered, checkboxes, description lists (\texttt{term \textbackslash{}:: description}) and nested sub-lists. Tap an item's text for edit, add-below and delete actions.
\item \textbf{Checkboxes} cycle unchecked → in-progress → done. A progress cookie in a heading's title (\texttt{[3/8]}, \texttt{[40\%]}, or the blank \texttt{[/]} and \texttt{[/8]}) is recalculated whenever a checkbox in the heading's subtree changes, counting the heading's own list and every descendant heading's, as org does. A heading's \texttt{:COOKIE\_DATA:} property changes what is counted: \texttt{checkbox} or \texttt{todo} picks which kind when a heading has both, both words count both kinds, and \texttt{recursive} with \texttt{todo} extends TODO counting to the whole subtree. Without it, only checkboxes are counted, recursively.
\item \textbf{Tables}: tap a cell to edit it in place. A \textbf{⎌} button beside the field discards the edit and restores the cell's text, which is the way out on a phone, where there is no Escape key (\texttt{Escape} also works with a hardware keyboard). Row and column add and remove controls sit below the table. \texttt{\#+TBLFM:} formulas are evaluated (see \hyperref[sec:org1458240]{Table formulas}). \textbf{Header rows} are the rows before the table's first horizontal rule (\texttt{|-{}-{}-+-{}-{}-|}) and render bold; a table with no rule has no header, and a later rule is only a group separator.
\item \textbf{Paragraphs}: tap the text for edit, add-below and delete actions.
\item \textbf{Captions, names and image sizes}: \texttt{\#+CAPTION:}, \texttt{\#+NAME:}, \texttt{\#+ATTR\_HTML:} (or \texttt{\#+ATTR\_ORG:}), \texttt{\#+HEADER:} and \texttt{\#+RESULTS:} lines directly above a paragraph, image, table or block describe it. \texttt{\#+CAPTION:} shows as an italic caption, with its inline markup rendered. \texttt{\#+ATTR\_HTML: :width 300 :align center} sizes and centers the images in that paragraph (a plain number is pixels; \texttt{\%}, \texttt{em} and similar lengths work). The other lines stay visible as small muted text, and every line is kept exactly as written.
\end{itemize}

Every delete asks for confirmation, and Undo/redo covers every edit (see \hyperref[sec:org9a920ce]{Undo/redo and History}).

\subsubsection{Table formulas}
\label{sec:org1458240}
\texttt{\#+TBLFM:} formulas use a subset of org's formula language. Org hands anything beyond simple arithmetic to Emacs Calc, which is not reproduced here.

\begin{center}
\begin{tabular}{ll}
Name & What it does\\[0pt]
\hline
\texttt{+ - * /} & Arithmetic. Division by zero gives \texttt{\#ERROR}\\[0pt]
\texttt{\textasciicircum{}} & Exponent, right-associative (\texttt{2\textasciicircum{}3\textasciicircum{}2} is 512)\\[0pt]
\texttt{sqrt(x)} & Square root; a negative input gives a complex number, \texttt{sqrt(-4)} is \texttt{(0, 2)}\\[0pt]
\texttt{floor(x)}, \texttt{ceil(x)}, \texttt{round(x)}, \texttt{trunc(x)} & Round toward −∞, toward +∞, to nearest (an exact .5 goes away from zero), toward zero. An optional second argument keeps that many decimals: \texttt{round(\$1, 2)}\\[0pt]
\texttt{sin(x)}, \texttt{cos(x)}, \texttt{tan(x)} & Degrees by default (see the \texttt{D} and \texttt{R} flags)\\[0pt]
\texttt{arcsin(x)}, \texttt{arccos(x)}, \texttt{arctan(x)} & Inverse trig, under Calc's names (not \texttt{asin}); outside −1 to 1 the result is complex\\[0pt]
\texttt{ln(x)}, \texttt{log10(x)}, \texttt{log(x, base)} & Logarithms; \texttt{log(x)} is natural. \texttt{x} of 0, or a \texttt{base} of exactly 1, gives \texttt{\#ERROR}; a negative number gives a complex result\\[0pt]
\texttt{exp(x)}, \texttt{exp10(x)} & e\textsuperscript{x} and 10\textsuperscript{x}\\[0pt]
\texttt{pi}, \texttt{inf} & Constants, written bare; \texttt{-inf} works\\[0pt]
\texttt{\$name} & A constant from a \texttt{\#+CONSTANTS:} line\\[0pt]
\texttt{vsum}, \texttt{vmean}, \texttt{vcount}, \texttt{vmin}, \texttt{vmax}, \texttt{vmedian} (\texttt{range}) & Sum, mean, count of non-blank cells, minimum, maximum, median\\[0pt]
\texttt{vvar}, \texttt{vpvar}, \texttt{vsdev}, \texttt{vpsdev} (\texttt{range}) & Sample and population variance and standard deviation. As in Calc, there is no plain \texttt{variance} or \texttt{stddev}\\[0pt]
\texttt{if(c,t,f)} & \texttt{t} if \texttt{c} is true (non-zero or non-empty), else \texttt{f}; only the chosen branch is evaluated\\[0pt]
\texttt{and(a,b)}, \texttt{or(a,b)}, \texttt{not(a)} & Function forms of the logical operators; they don't short-circuit\\[0pt]
\texttt{string(ref)} & A cell's raw text as a string\\[0pt]
\texttt{"text"} & A string literal (\texttt{\textbackslash{}"} for a quote inside it)\\[0pt]
\texttt{date(ref)}, \texttt{date(Y,M,D)}, \texttt{now()} & Date values (see below)\\[0pt]
\texttt{deg(x)} & Turns an HMS duration back into a plain number of days\\[0pt]
\texttt{date-to-time(ref)} & A date value from a wider range of date strings\\[0pt]
\texttt{format-time-string(fmt, time, utc)} & Formats a time with \texttt{\%} codes (see \hyperref[sec:org83b5238]{Capture templates})\\[0pt]
\texttt{uconv(value, "from", "to")} & Unit conversion (see below)\\[0pt]
\end{tabular}
\end{center}

\textbf{Logic}: comparisons (\texttt{==} or \texttt{=}, \texttt{!=}, \texttt{<}, \texttt{>}, \texttt{<=}, \texttt{>=}) give \texttt{1} or \texttt{0}. They bind tighter than \texttt{\&\&} and \texttt{||}, which bind looser than \texttt{+ - * /}, so \texttt{\$1 >= 50 \&\& \$2 >= 50} means \texttt{(\$1 >= 50) \&\& (\$2 >= 50)}. \texttt{!x} is logical NOT. \texttt{\&\&}, \texttt{||} and \texttt{if} short-circuit, so \texttt{if(\$1 != 0, 100 / \$1, 0)} never divides by zero. Nested \texttt{if} calls make an else-if chain: \texttt{if(\$1 >= 90, "A", if(\$1 >= 80, "B", "C"))}. A string literal is text and \texttt{string(ref)} reads a cell as text, so \texttt{string(\$1) == "Sales"} compares text, where a bare \texttt{\$1 == "Sales"} would try to read the cell as a number.

\textbf{Named constants}: a \texttt{\#+CONSTANTS:} line anywhere in the file defines constants, used as \texttt{\$name}: \texttt{\#+CONSTANTS: c=299792458 g=9.81}. Names are space-separated, several lines combine, and a later value wins. An undefined \texttt{\$name} gives \texttt{\#ERROR}. There is no built-in physical-constants library (org's optional \texttt{constants.el}); only \texttt{pi} and \texttt{inf} are built in.

\textbf{Dates and times}: a cell holding an org timestamp (\texttt{[YYYY-MM-DD]} or \texttt{<YYYY-MM-DD>}, with an optional time and day name) is read as a date as soon as it is referenced, so \texttt{\$3 = \$1 - \$2} subtracts two timestamp cells. Every component is range-checked (month 1 to 12, a day within that month's length with leap years, hour 0 to 23, minute 0 to 59), so \texttt{<2026-06-34>} gives \texttt{\#ERROR}. \texttt{date(ref)} differs from a bare reference only for a cell that isn't a timestamp: bare, it reads as an ordinary number; wrapped, it gives \texttt{\#ERROR}, which is useful when a formula should fail loudly. \texttt{date(Y, M, D)} builds a date from three numbers (\texttt{date(2028, 2, 29)} is valid, \texttt{date(2026, 2, 29)} is not).

\begin{itemize}
\item Subtracting two dates gives days as an integer, or, if either has a time of day, a duration in Calc's HMS form (\texttt{8@ 30' 0"} for 8.5 hours).
\item Adding or subtracting a number shifts a date by that many days and writes it back as a timestamp: \texttt{date(\$1) + 14} gives \texttt{<2026-09-08>}.
\item \texttt{deg()} turns an HMS duration back into days for further arithmetic: \texttt{deg(date(\$2) - date(\$1)) * 24;\%.2f} gives \texttt{8.50} hours.
\item \texttt{now()} is the current date and time.
\item \texttt{date-to-time(ref)}, spelled as in Emacs, accepts every timestamp \texttt{date()} does, plus ISO 8601 (\texttt{2026-01-15T10:30:00Z}), RFC 822 (\texttt{Mon, 15 Jan 2026 10:30:00 GMT}) and US (\texttt{01/15/2026}) strings, parsed by the browser. An unparseable string gives \texttt{\#ERROR}.
\item \texttt{format-time-string(fmt, time, utc)} uses the same \texttt{\%} codes as capture templates: \texttt{format-time-string("\%Y-\%m-\%d", \$1)} gives \texttt{2026-01-15}. \texttt{time} defaults to now and accepts a date value or a number of seconds since the Unix epoch. A non-zero \texttt{utc} formats in UTC (\texttt{\%Z} and \texttt{\%z} then say UTC and +0000); otherwise the device's time zone is used.
\end{itemize}

\textbf{References}: \texttt{\$N} is column N of the current row, \texttt{@N\$M} is row N, column M, \texttt{@-N} and \texttt{@+N} are rows relative to the target, \texttt{@<}, \texttt{@>}, \texttt{\$<} and \texttt{\$>} are the first and last row or column, and a range such as \texttt{\$1..\$3} or \texttt{@1\$1..@3\$2} is an aggregate function's argument. A range may omit the column (\texttt{@2..@6}) to mean those rows in the formula's own column. Rows count every actual row, including the header; a horizontal rule has no row number. The column shorthand \texttt{\$M=} (no \texttt{@}, applied to every row) skips rows before the first horizontal rule, while an explicit \texttt{@1\$2=} still changes that row. A table with no rule has no header, so nothing is skipped.

\texttt{@I}, \texttt{@II}, \texttt{@III} refer to horizontal rules by repeated "I" characters, not Roman numerals (\texttt{@IIII} is the fourth). An offset such as \texttt{@III+2} is the data row that many rows after (or, with \texttt{-}, before) the rule. A range between two rules, \texttt{vsum(@I..@II)}, covers the data rows strictly between them. A count past the last rule clamps to the end of the table.

\textbf{Targets}: a formula's target is one cell (\texttt{@N\$M=}), a range (\texttt{@N\$M..@N\$M=}, applied to each cell; each end needs an explicit row) or a whole column (\texttt{\$M=}). Several formulas in one \texttt{\#+TBLFM:} line are separated by \texttt{::} and run in order, each seeing the earlier results. Only the first \texttt{\#+TBLFM:} line directly below a table counts.

\textbf{Recalculating} is explicit. A \textbf{🔢 Calc} button below a table with a \texttt{\#+TBLFM:} line recalculates that table. \textbf{Recalculate all tables} in the command palette (\texttt{org-table-recalculate-buffer-tables}, also usable in an Extras entry) recalculates every table as one undo step. A malformed formula is reported, and only its own table is skipped.

\textbf{Blank and non-numeric cells}: by default a blank cell in a direct reference (\texttt{\$1+\$2}) reads as 0 but is left out of a range, so \texttt{vmean} averages only the non-blank cells and \texttt{vcount} skips it. Text that isn't a number reads as \texttt{nan} and propagates, so a stray label shows a problem instead of a plausible wrong number (the \texttt{N} flag below reads it as 0).

\textbf{Mode flags}: a trailing \texttt{;flags} string after the expression:

\begin{center}
\begin{tabular}{ll}
Flag & Effect\\[0pt]
\hline
\texttt{\%.Nf}, \texttt{fN} & Fixed, N decimals (\texttt{\%0.Nf} also works): \texttt{4.85;f2} gives \texttt{4.85}\\[0pt]
\texttt{\%d} & Truncates toward zero: \texttt{4.85} gives \texttt{4}\\[0pt]
\texttt{nN} & Normal notation, N significant figures: \texttt{3.14159265;n5} gives \texttt{3.1416}\\[0pt]
\texttt{sN} & Scientific notation, N figures: \texttt{123456.789;s3} gives \texttt{1.23e5}\\[0pt]
\texttt{eN} & Engineering notation (exponent a multiple of 3): \texttt{123456.789;e3} gives \texttt{123e3}\\[0pt]
\texttt{pN} & Computes at N significant digits, in arbitrary precision (see \hyperref[sec:org45678a6]{Development}). Degree-mode trig is limited to about double precision, and sqrt of a negative still errors\\[0pt]
\texttt{E} & Keeps a blank cell in a range (as \texttt{nan}) and reads a blank direct reference as \texttt{nan}\\[0pt]
\texttt{N} & Non-numeric text (with \texttt{E}, blanks too) reads as 0 instead of \texttt{nan}\\[0pt]
\texttt{D}, \texttt{R} & Angle mode for the trig functions: degrees or radians. Degrees is the default, so \texttt{D} changes nothing\\[0pt]
\texttt{T}, \texttt{U}, \texttt{t} & Duration mode. Cells matching \texttt{HH:MM[:SS]} (or a bare integer, in seconds) are read as durations, and the result is \texttt{HH:MM:SS} (\texttt{T}), \texttt{HH:MM} (\texttt{U}) or decimal hours to two places (\texttt{t}). Without a flag, \texttt{HH:MM} text isn't a duration. Negative durations work, and \texttt{org-table-duration-hour-zero-padding} (see \hyperref[sec:org75c5aa1]{Local Variables}, default on) pads the hours\\[0pt]
\texttt{F} & Fraction mode: exact rational results (\texttt{1/3;F} is exactly \texttt{1/3}, and \texttt{\$1/\$2;F} on 3 and 6 is \texttt{1/2}; a whole result shows without \texttt{/1}). A function with a rational result stays exact (\texttt{sqrt(4);F} is \texttt{2}); otherwise a short approximation (denominator under 1,000,000) is shown\\[0pt]
\texttt{FS} & Fraction mode as a mixed number: \texttt{14/4;FS} is \texttt{3 + 1/2}. \texttt{S} alone does nothing, since general Symbolic mode isn't implemented\\[0pt]
\end{tabular}
\end{center}

An unrecognized flag is ignored. Examples:

\begin{verbatim}
|  Task 1 |   Task 2 |    Total |
|---------+----------+----------|
|    2:12 |     1:47 | 03:59:00 |
|    2:12 |     1:47 |    03:59 |
| 3:02:20 | -2:07:00 |     0.92 |
#+TBLFM: @2$3=$1+$2;T::@3$3=$1+$2;U::@4$3=$1+$2;t
\end{verbatim}

\begin{verbatim}
| Top | Bottom | Fraction | Mixed number |
|-----+--------+----------+--------------|
|   3 |      6 | 1/2      | 1/2          |
|  12 |      9 | 4/3      | 1 + 1/3      |
|  14 |      4 | 7/2      | 3 + 1/2      |
#+TBLFM: $3=$1/$2;F::$4=$1/$2;FS
\end{verbatim}

\textbf{Errors}: a formula that fails when evaluated (an out-of-range reference, a range where one value is expected) writes \texttt{\#ERROR} into that one cell, and the rest of the table still updates. A formula that can't be parsed (bad syntax, an unknown function) is reported for the whole recalculation.

\textbf{Raw Elisp}: a formula starting with a quote, like org's \texttt{'(function-name arg ...)}, bypasses the grammar above. One function is supported, \texttt{orgtbl-ascii-draw(value, min, max, width)}, which draws a bar in the cell using org's 11-character gradient. A \texttt{value} outside \texttt{[min, max]} shows \texttt{too small} or \texttt{too large}, and \texttt{width} defaults to 12. Each argument may be a number, a \texttt{\$N} reference or a parenthesized expression such as \texttt{(\$2 * 2)}.

\textbf{Unit conversion}: \texttt{uconv(value, "from", "to")}, such as \texttt{uconv(\$1, "mi/hr", "m/s")}, uses the mathjs unit library (see Development). It handles compound units (\texttt{kg/m\textasciicircum{}3}) and converts temperature as an absolute reading (\texttt{0 degC} is \texttt{32 degF}, like Calc's \texttt{u t} command rather than the relative change of \texttt{u c}). Unlike Calc's two-argument form, the source unit is explicit. An incompatible pair (\texttt{m} to \texttt{kg}) or an unknown unit gives \texttt{\#ERROR}. Common units: length (\texttt{m}, \texttt{ft}, \texttt{km}, \texttt{mi}, \texttt{in}, \texttt{cm}), mass (\texttt{kg}, \texttt{lb}, \texttt{g}, \texttt{oz}), temperature (\texttt{degC}, \texttt{degF}, \texttt{K}), volume (\texttt{L}, \texttt{gal}, \texttt{mL}), time (\texttt{s}, \texttt{min}, \texttt{hr}, \texttt{day}), speed (\texttt{mi/hr}, \texttt{m/s}), data (\texttt{MB}, \texttt{GB}), pressure (\texttt{bar}, \texttt{psi}, \texttt{atm}), energy (\texttt{J}, \texttt{kWh}) and area (\texttt{m\textasciicircum{}2}, \texttt{acre}). Names mostly match Calc's; \texttt{tsp} and \texttt{tbsp} are aliases for mathjs's \texttt{teaspoon} and \texttt{tablespoon}. The full list is at \url{https://mathjs.org/docs/datatypes/units.html}.

\textbf{Complex numbers}: \texttt{(re, im)} or \texttt{(r; θ)} in a cell or formula, and \texttt{i} for (0, 1). \texttt{+ - * / \textasciicircum{}}, \texttt{sqrt}, \texttt{exp}, \texttt{ln}, \texttt{log10}, the trig functions and \texttt{vsum}, \texttt{vmean}, \texttt{vprod} accept them, along with \texttt{re}, \texttt{im}, \texttt{conj}, \texttt{abs}, \texttt{arg}, \texttt{polar} and \texttt{rect}. Whole numbers stay exact and other parts show as floats (\texttt{2.}), a zero imaginary part drops to a plain number, and \texttt{pN} sets the working digits. Unlike Calc, \texttt{i} is the number, so \texttt{2+3*i} is \texttt{(2, 3)}, not \texttt{3 i + 2}.

\begin{verbatim}
| (1, 2) | (3, 4) | (-5, 10) |
#+TBLFM: $3=$1*$2
\end{verbatim}

\textbf{Vectors and matrices}: \texttt{[1, 2, 3]} and \texttt{[[1, 2], [3, 4]]} in a cell or formula; a cell range such as \texttt{\$1..\$3} is a vector of its non-blank cells. \texttt{+ -} and \texttt{* /} by a number work elementwise; \texttt{*} is a dot product for two vectors, a matrix product for matrices, and \texttt{\textasciicircum{}} takes a matrix to an integer power. Functions: \texttt{det}, \texttt{inv}, \texttt{trn}, \texttt{tr}, \texttt{cross}, \texttt{vlen}, \texttt{abs} (the length), and the \texttt{v} aggregates, which flatten.

Out of scope: Emacs Calc beyond the above (symbolic algebra, the rest of Calc's vector and matrix library), named-field references, and recalculation on every keystroke.

\subsubsection{Plots}
\label{sec:org0e13f05}
A \texttt{\#+PLOT:} line directly above a table (several stacked lines combine their options) graphs the table's data. A \textbf{📊 Plot} button below the table, beside \textbf{🔢 Calc} when present, shows the graph inline, and tapping it again hides it. This is a subset of org-plot: it reads org-plot's own \texttt{\#+PLOT:} options, so an unmodified line works, but it renders inline SVG instead of calling gnuplot.

\texttt{type:2d} (the default) and \texttt{type:radar} are supported; \texttt{type:3d} and \texttt{type:grid} are refused with an error. For 2d, \texttt{with:} picks \texttt{lines} (the default), \texttt{points}, \texttt{boxes} (\texttt{histograms} renders the same) or \texttt{impulses}.

\begin{center}
\begin{tabular}{ll}
Option & Meaning\\[0pt]
\hline
\texttt{title} & A quoted string\\[0pt]
\texttt{ind} & The 1-based column giving the x position (2d) or the axis label (radar); default 1\\[0pt]
\texttt{deps} & A parenthesized list of 1-based columns to plot; default every column except \texttt{ind}\\[0pt]
\texttt{labels} & A parenthesized list of quoted series names; default the header row\\[0pt]
\texttt{transpose}, \texttt{trans} & \texttt{yes}, \texttt{y} or \texttt{t} swaps rows and columns first, which is how a radar chart gets its axes from columns\\[0pt]
\texttt{min}, \texttt{max} & The y range (2d), or the shared range of every radar axis\\[0pt]
\texttt{xmin}, \texttt{xmax} & The x range (2d, when \texttt{ind} is numeric)\\[0pt]
\texttt{ticks} & The number of y gridlines; default 5\\[0pt]
\end{tabular}
\end{center}

org-plot's gnuplot scripting options (\texttt{set}, \texttt{line}, \texttt{script}, \texttt{file}, \texttt{map}, \texttt{timefmt}) are accepted but have no effect, so a copied line never errors.

The graph is cached for each \texttt{\#+PLOT:} line and its table data, so showing it again is instant, and editing the table re-renders it the next time it is shown. It follows the heading's fold state: collapsing the heading hides the graph, and after expanding it the \textbf{Plot} button needs a fresh tap, as with contact photos (see \hyperref[sec:org75c5aa1]{\texttt{org-xx-startup-with-show-photos}}).

\noindent\rule{\textwidth}{0.5pt}

\subsection{Inline text markup}
\label{sec:org1a88f72}
All six of org's emphasis markers work in headings, paragraphs and list items:

\begin{center}
\begin{tabular}{lll}
Markup & Syntax & Renders as\\[0pt]
\hline
Bold & \texttt{*bold text*} & \textbf{bold text}\\[0pt]
Italic & \texttt{/italic text/} & \emph{italic text}\\[0pt]
Underline & \texttt{\_underline text\_} & underlined\\[0pt]
Strikethrough & \texttt{+strikethrough text+} & \sout{strikethrough text}\\[0pt]
Verbatim & \texttt{=verbatim text=} & \texttt{verbatim text} (literal, not further parsed)\\[0pt]
Code & \texttt{\textasciitilde{}code text\textasciitilde{}} & \texttt{code text} (literal, not further parsed)\\[0pt]
\end{tabular}
\end{center}

A marker opens where it follows the start of a line, whitespace or one of \texttt{-(\{'"} and isn't followed by whitespace. It closes where it isn't preceded by whitespace and is followed by the end of the line, whitespace or closing punctuation. These are org's border rules for the common cases, not its full border-character tables. Bold, italic, underline and strikethrough nest (\texttt{*bold /italic/ bold*}); verbatim and code never do, and their content is literal.

\textbf{Horizontal rule}: a line of only dashes, at least five (\texttt{-{}-{}-{}-{}-}), is a horizontal rule. Whitespace around the dashes is fine; any other character on the line makes it a paragraph.

\textbf{Hard line break}: a line ending in two backslashes (\texttt{\textbackslash{}\textbackslash{}}) is org's explicit line break. Source lines in an ordinary paragraph or list item otherwise reflow together with a space (org's \texttt{org-export-preserve-breaks} is \texttt{nil} by default). The marker forces a break exactly where it is written, in paragraphs, \hyperref[sec:org22c076e]{QUOTE or CENTER blocks} and all four document export formats (HTML, Markdown, ASCII, ODT).

\textbf{Subscript and superscript}: \texttt{\_} and \texttt{\textasciicircum{}} mark them. \texttt{x\_i} and \texttt{x\textasciicircum{}2} work without braces, and braces (\texttt{x\_\{alpha\}}, \texttt{10\textasciicircum{}\{100\}}) are always accepted, mainly for readability. \texttt{org-use-sub-superscripts} (see \hyperref[sec:org75c5aa1]{Local Variables}) controls them:

\begin{center}
\begin{tabular}{ll}
Value & Behavior\\[0pt]
\hline
\texttt{t} (default) & Both \texttt{x\_i} and \texttt{x\_\{i\}} are interpreted\\[0pt]
\texttt{\{\}} & Only the braced form \texttt{x\_\{i\}} is interpreted; a bare \texttt{x\_i} stays literal text\\[0pt]
\texttt{nil} & Disabled entirely; \texttt{\_} / \texttt{\textasciicircum{}} are always literal characters\\[0pt]
\end{tabular}
\end{center}

\texttt{\_} also marks underline, as in org. The two are told apart by what precedes it: underline needs whitespace, a line start or certain punctuation before it (\texttt{\_underlined\_}), and subscript needs a non-space character (\texttt{a\_b}). Script content is literal: \texttt{x\_\{*bold*\}} doesn't interpret the \texttt{*bold*}.

\textbf{Footnotes}: \texttt{[fn:label]} is a reference, shown as a small tappable superscript that jumps to, and briefly highlights, the definition anywhere in the document, even under a folded heading. A definition can be inline at the reference, \texttt{[fn:label:the note text]}, or anonymous, \texttt{[fn::an anonymous note]}, and inline text renders in place, smaller and in italics. A separate line starting \texttt{[fn:label]} (like \texttt{[fn:1] The actual definition.}) is also a definition, shown with a small label. Markdown export uses GFM's \texttt{[\textasciicircum{}label]} and \texttt{[\textasciicircum{}label]: text}, and HTML export uses a jump-linked superscript with a back-linked Footnotes section.

Table cells render inline markup too. Tap a cell to see and edit its raw text, and tap away or tap another cell to commit (see \hyperref[sec:org22c076e]{Editing your outline}).

\subsubsection{Footnote commands}
\label{sec:orga27d7d4}
In the outline, open a paragraph, a list item or a heading's text with \textbf{Edit text}, put the cursor after the words the note belongs to and tap \textbf{Footnote}. A sheet (a dialog on a large screen) asks for the label, already the next free number, and the note, with a switch to write it inline; \textbf{Insert} adds the reference and the note under \texttt{* Footnotes}, and Undo takes it back. Tap a reference to see its note and \textbf{Edit note}, \textbf{Go to the note} or \textbf{Delete footnote}; a reference with no note offers \textbf{Create the note}. With a paragraph selected from the keyboard, the \textbf{Footnote} command opens that paragraph's sheet.

In Text view, \textbf{Footnote} (\texttt{org-footnote-action}, \texttt{C-c C-x f}, palette) does the right thing at the cursor: on a reference it jumps to the definition (offering to create one if there is none), on a definition's label it jumps back to the reference, and anywhere else it adds a new footnote. Where a footnote can't go (start of a line, source blocks, comments, links, inside another reference) it offers the one-key menu: \texttt{s\textasciitilde{}ort, \textasciitilde{}r\textasciitilde{}enumber, \textasciitilde{}S} both, \textasciitilde{}n\textasciitilde{}ormalize, \textasciitilde{}d\textasciitilde{}elete.

\begin{itemize}
\item \textbf{Insert footnote} (\texttt{org-footnote-new}): a reference at the cursor and the start of its definition under \texttt{* Footnotes} at the end of the file, with the cursor there to write it. The label is the next unused number; an existing label adds another reference to it, and an empty one makes an anonymous \texttt{[fn::]}.
\item \textbf{Normalize footnotes} (\texttt{org-footnote-normalize}): every footnote becomes a numbered one in reading order, inline and anonymous notes become definitions, and all definitions move under \texttt{* Footnotes} (\texttt{DEFINITION NOT FOUND.} for a reference nobody defined).
\item \textbf{Sort footnotes} (\texttt{org-footnote-sort}) puts the definitions in the order of their references. \textbf{Renumber footnotes} (\texttt{org-footnote-renumber-fn:N}) numbers \texttt{fn:1}, \textasciitilde{}fn:2\textasciitilde{}… by first reference.
\item \textbf{Delete footnote} (\texttt{org-footnote-delete}): the one at the cursor, or the label you give, with every definition and reference.
\end{itemize}

Insert, sort, renumber, normalize and delete work from any view, and Undo reverses them. Settings: \texttt{org-footnote-section} (the heading; \texttt{nil} keeps each definition at the end of its own section), \texttt{org-footnote-auto-label}, \texttt{org-footnote-define-inline} and \texttt{org-footnote-auto-adjust} (renumber and/or sort after adding or deleting). \texttt{\#+STARTUP:} accepts \texttt{fninline}, \texttt{fnlocal}, \texttt{fnauto}, \texttt{fnprompt}, \texttt{fnconfirm}, \texttt{fnplain}, \texttt{fnadjust} and \texttt{nofnadjust}.

\subsubsection{\LaTeX{} math fragments}
\label{sec:org0b05dda}
Math in a heading, paragraph or list item renders as typeset math with \href{https://katex.org}{KaTeX}, vendored locally (no CDN). It is display-only: the \texttt{.org} file's text is untouched, and editing shows the raw source.

\begin{center}
\begin{tabular}{lll}
Delimiter & Kind & Restrictions\\[0pt]
\hline
\texttt{\textbackslash{}(...\textbackslash{})} & Inline & None — org's own recommended form\\[0pt]
\texttt{\textbackslash{}[...\textbackslash{}]} & Display & None\\[0pt]
\texttt{\$\$...\$\$} & Display & None\\[0pt]
\texttt{\$...\$} & Inline & No whitespace directly inside the delimiters; the closing \texttt{\$} must be followed by whitespace or punctuation, but never a dash (avoiding \texttt{\$5-10} reading as math)\\[0pt]
\texttt{\textbackslash{}begin\{env\}...\textbackslash{}end\{env\}} & Display & \texttt{\textbackslash{}begin} must start a line, preceded by nothing but whitespace\\[0pt]
\end{tabular}
\end{center}

\texttt{\$...\$} avoids misreading currency: \texttt{costs \$5 and \$10 total} stays plain text. The opening \texttt{\$} may follow any character except another \texttt{\$} (or the start of the line). The first content character can't be a space, tab, newline, comma, semicolon or period, and the character before the closing \texttt{\$} can't be a space, tab, newline, comma or period (a semicolon is allowed). It needs at least two characters of content, so \texttt{\$x\$} doesn't match; use \texttt{\textbackslash{}(x\textbackslash{})} for a single character.

\texttt{\textbackslash{}(...\textbackslash{})}, \texttt{\textbackslash{}[...\textbackslash{}]}, \texttt{\$\$...\$\$} and \texttt{\textbackslash{}begin\{env\}...\textbackslash{}end\{env\}} can span any number of lines, even across a blank line (an \texttt{align} environment with a blank line between rows). \texttt{\$...\$} allows at most two embedded newlines. KaTeX doesn't treat \texttt{\%} as a comment, so text after it renders literally. A fragment KaTeX can't render (an unsupported command, a syntax error) shows its raw source with a dashed border, without breaking the paragraph. KaTeX supports a large subset of \LaTeX{}, not all of it.

Math renders wherever inline markup does: the outline, the read-only Docs viewer and HTML export (KaTeX's CSS is inlined only into an export that contains math). Markdown, ASCII and ODT export pass the raw \LaTeX{} through as text.

\textbf{Examples}: the source \texttt{The quadratic formula is \textbackslash{}(x = \textbackslash{}frac\{-b \textbackslash{}pm \textbackslash{}sqrt\{b\textasciicircum{}2-4ac\}\}\{2a\}\textbackslash{}).} renders inline with a proper fraction bar and a square-root radical. This \texttt{align} environment:

\begin{verbatim}
\begin{align}
E &= mc^2 \\
F &= ma \\
a^2 + b^2 &= c^2
\end{align}
\end{verbatim}

renders three right-aligned equations, lined up at the \texttt{\&} marks and automatically numbered (1) to (3), which is KaTeX's own behavior for \texttt{align}. The \texttt{\textbackslash{}\textbackslash{}} ending a line there ends an equation; it is unrelated to org's hard line break above.

\noindent\rule{\textwidth}{0.5pt}

\subsection{Searching}
\label{sec:orgf0231a2}
Open the command palette's \textbf{Search} (or god-mode's \texttt{s}, see \hyperref[sec:org996c4cc]{Keybindings}) to search the open document, case-insensitively: headings, tags, TODO keyword, priority, properties (key or value), SCHEDULED and DEADLINE timestamps, paragraphs, list items, table cells and block content. When any are configured it also searches every file in \texttt{org-agenda-files}.

Results appear at once from what is already loaded and fill in as each agenda file finishes loading; a file that fails to fetch is left out. The query box and its options live in the minibuffer at the bottom of the screen (see \hyperref[sec:org6cade45]{Modeline and minibuffer}). Results update as you type, grouped under the heading each belongs to, with a snippet that highlights the matched keyword. A result from another file shows that file's name.

Search covers the whole document whatever the fold state: a match inside a collapsed heading, hidden body text, a folded drawer or a block still shows up. Tapping a result expands what is needed to reveal it (the heading's body, and its drawers for a property or block match), scrolls to it with a brief highlight, and switches to that file first when the result came from another one.

\subsubsection{Filter tokens}
\label{sec:orgbbfbd4f}

Tokens in the query box combine as an implicit AND:

\begin{itemize}
\item \texttt{word} or \texttt{+word}: the entry (title, tags, TODO state, priority, properties, planning and body text) must contain it. \texttt{-word} excludes entries that contain it. \texttt{+computer +wifi -ethernet}, or just \texttt{computer wifi -ethernet}, finds entries with "computer" and "wifi" but not "ethernet", as in org's \texttt{org-search-view}. So a query like \texttt{-5} means "exclude entries containing 5", not the literal text; use Regex mode for that.
\item \texttt{tag:urgent}: the heading must have this tag, inherited tags included. \texttt{-tag:urgent} excludes it.
\item \texttt{todo:KEYWORD}: the TODO state must equal this (case-insensitive), as in \texttt{todo:WAITING}. \texttt{-todo:KEYWORD} excludes it; a heading with no keyword passes an exclude.
\item \texttt{priority:X}: the priority must equal this, as in \texttt{priority:A}. \texttt{-priority:X} excludes it.
\item \texttt{key:value}: any other \texttt{key:value} matches a property \texttt{key} whose value contains \texttt{value} (substring, case-insensitive), such as \texttt{spouse:Mary}. \texttt{-key:value} excludes it.
\item \texttt{|} inside \texttt{tag:}, \texttt{todo:}, \texttt{priority:} or \texttt{key:value} means "any of these": \texttt{tag:urgent|blocked}, \texttt{todo:TODO|WAITING}, \texttt{assignee:Bob|Alice}. Negated, it means "none of these". It combines alternatives only within one filter. There is no OR between filters and no grouping; Match mode has both.
\end{itemize}

A word starting with \texttt{+} or \texttt{-} (a negative number, a hyphenated word) or shaped \texttt{word:word} (such as a \texttt{mailto:} link) can be read as a filter instead of text. \texttt{http://} and \texttt{https://} URLs are not. Tag inheritance is on by default and property inheritance off, as in org (see \hyperref[sec:org2d2d41c]{Differences from Emacs org-mode} and the \texttt{org-use-tag-inheritance} and \texttt{org-use-property-inheritance} variables). A query with only excludes or filters and no keyword (\texttt{tag:work -someday}) returns the matching headings directly. The box's placeholder, \texttt{Search, +word -word, or [tag|todo|priority|key]:value…}, is a reminder of the pattern.

Example: \texttt{+work todo:WAITING budget} finds entries tagged \texttt{work}, in WAITING state, containing "budget".

\subsubsection{Regex mode}
\label{sec:org9edddd8}

The \textbf{Regex} toggle in the minibuffer switches from plain substring matching to a regular expression, case-insensitive either way, in Emacs's syntax rather than JavaScript's. The biggest difference: an unescaped \texttt{(}, \texttt{)}, \texttt{\{}, \texttt{\}} or \texttt{|} is literal, and grouping needs \texttt{\textbackslash{}(} and \texttt{\textbackslash{})}, alternation \texttt{\textbackslash{}|}, and an interval \texttt{\textbackslash{}\{n,m\textbackslash{}\}}. With Regex off, a \texttt{.} or \texttt{(} in the query simply matches that character.

Supported: backreferences (\texttt{\textbackslash{}1} to \texttt{\textbackslash{}9}), word and symbol boundaries (\texttt{\textbackslash{}<}, \texttt{\textbackslash{}>}, \texttt{\textbackslash{}\_<}, \texttt{\textbackslash{}\_>}, \texttt{\textbackslash{}b}, \texttt{\textbackslash{}B}), POSIX classes in brackets (\texttt{[[:digit:]]}, \texttt{[[:alpha:]]}) and whole-string anchors (\texttt{\textbackslash{}`} and \texttt{\textbackslash{}'}, distinct from the line anchors \texttt{\textasciicircum{}} and \texttt{\$}). Constructs that depend on a buffer's syntax table (\texttt{\textbackslash{}s<}, \texttt{\textbackslash{}s>}, \texttt{\textbackslash{}sc}, \texttt{\textbackslash{}S-} and similar) aren't supported and show an error naming the construct. An invalid pattern shows a clear error, so it is never mistaken for "no matches".

Example: \texttt{\textasciicircum{}\textbackslash{}(Buy\textbackslash{}|Fix\textbackslash{})} finds headings whose title starts with "Buy" or "Fix".

\subsubsection{Match mode}
\label{sec:orgcbe7580}

\textbf{Match}, beside \textbf{Regex} in the options popover (the \texttt{⋯} button at the search bar's right edge, which shows a check mark by the active one and a dot when either is on), uses org's tag and property match grammar, the one behind \texttt{C-c / m}. Match and Regex are mutually exclusive.

\texttt{|} is the lowest-precedence operator and splits the query into OR'd groups. Within a group, \texttt{+family} (or a bare \texttt{family}) and \texttt{-work} combine as AND, so \texttt{family+work|urgent} means "(family AND work) OR urgent". A condition can also compare a property, \texttt{PROPERTY="value"} with \texttt{=}, \texttt{<>}, \texttt{<}, \texttt{>}, \texttt{<=} or \texttt{>=}, combined freely with tags (\texttt{work+PRIORITY\_TEST="B"}). \texttt{TODO}, \texttt{LEVEL}, \texttt{PRIORITY} and \texttt{CATEGORY} work as pseudo-properties, as in org. An empty query, an empty group beside a stray \texttt{|}, or an unparseable condition gives a clear error. Replace is disabled in this mode.

Example: \texttt{work+urgent|blocked} finds headings tagged both \texttt{work} and \texttt{urgent}, or tagged \texttt{blocked}.

\subsubsection{Search extra files}
\label{sec:orgdf17868}
\texttt{org-agenda-text-search-extra-files} lists more files to search, in addition to the file you have open and your \hyperref[sec:org2e9fd13]{Agenda Files}. It is the same \texttt{scheme:path} syntax (\texttt{github:}, \texttt{webdav:}, \texttt{local:}), semicolon-separated:

\begin{verbatim}
org-agenda-text-search-extra-files: github:notes/reference.org;webdav:archive/2024.org
\end{verbatim}

Set it in \textbf{Settings → Quick Settings → Search extra files}, or per document in Local Variables. The files are searched, and nothing else: they don't appear in the Agenda, TODO or tag views, which makes this the place for notes that should be findable but not scheduled. A file listed here and as an agenda file is searched once. Tapping a result opens its file. A file that is still loading, or can't be loaded, is named above the results, so it never just quietly finds nothing. \textbf{Replace} works on the open file and your agenda files only.

\subsubsection{Search and replace}
\label{sec:orgce67b07}

The options button (\texttt{⋯}) also shows \textbf{Replace}, which starts an Emacs-style \texttt{query-replace} (or \texttt{query-replace-regexp}, following the Regex toggle) with the search text as the pattern. It walks every \hyperref[sec:org2e9fd13]{Agenda Files} file that has a match, starting with the open one, and a file it enters opens as its own tab. It asks for the replacement text, then offers each match:

\begin{center}
\begin{tabular}{ll}
Key & Action\\[0pt]
\hline
\texttt{y} or Space & Replace this one and go on\\[0pt]
\texttt{n} or Delete & Skip it\\[0pt]
\texttt{q} or Return & Stop the walk (replacements already made stay; the rest is untouched)\\[0pt]
\texttt{.} & Replace this one and stop\\[0pt]
\texttt{!} & Replace all remaining matches in the current file, then go on to the next file with a match\\[0pt]
\end{tabular}
\end{center}

Buttons for all five sit beside the prompt for devices without a keyboard. \texttt{\$1} to \texttt{\$9} in the replacement text are filled from the match's capture groups. As in Emacs, nothing is saved automatically: each touched file is changed but unsaved, shown by its tab's modified indicator. A configured agenda file that can't be reached (a local file not already open) is left out; only open files and GitHub or WebDAV files can be walked.

Replace works on heading titles, paragraphs, list items and table cells. Tags, TODO state, priority, properties, planning timestamps and block content (\texttt{\#+BEGIN\_VERSE}, \texttt{\#+BEGIN\_QUOTE} and the like) are never replaced, because that could corrupt the document's syntax. Search still finds text inside blocks, and when all of a query's matches are inside block content, Replace says so ("N matches found only inside block content\ldots{} not editable here; use Text view instead"). \hyperref[sec:orga9cae38]{Text view}'s own search and replace has no such restriction and stays within that one document.

\subsubsection{Narrow}
\label{sec:org5954f28}

When at least one match is in the open document, the results show \textbf{Narrow} with a count of those matches (cross-file results aren't counted). It works with any search mode. Tapping it restricts the outline to the matching headings plus their ancestors, so each match keeps its context, and a banner shows the count with a \textbf{Widen} button. It composes with \hyperref[sec:org77ba459]{subtree narrowing}: narrow to a subtree, then to matches within it. Like subtree narrowing it survives a reload, and a heading renamed or removed in the meantime simply drops out.

Org's own search is several separate tools: sparse trees (\texttt{C-c /}), a match grammar, and \texttt{org-search-view}. Here one box covers tags, TODO state, priority, properties and free text together, with Match mode for the fuller boolean grammar. Not included: sparse-tree-style in-place refolding; Narrow is the nearest equivalent.

\noindent\rule{\textwidth}{0.5pt}

\subsection{Modeline and minibuffer}
\label{sec:org6cade45}
A fixed two-row bar at the bottom of the screen, laid out like Emacs's frame: a dense modeline above a minibuffer.

\textbf{The modeline} is read-only, monospace and terse:

\begin{itemize}
\item \textbf{Modified indicator}: \texttt{-{}-{}} unmodified, \texttt{**} modified, \texttt{\%\%} read-only, \texttt{\%*} read-only and modified, Emacs's four states. Emacs's coding-system and line-ending segment is omitted, since the app always writes UTF-8 with Unix line endings.
\item \textbf{Buffer name}: the document's name plus a suffix, \texttt{(GitHub)}, \texttt{(WebDAV)}, \texttt{(Local)} or \texttt{(Imported)}.
\item \textbf{Position}: \texttt{Top}, \texttt{Bot}, \texttt{All} or a percentage for how far down the visible area is, plus \texttt{L} and the line number when a heading is focused, and \texttt{C} and the column only while a text field is being edited. The app has no persistent cursor between edits, so otherwise there is no real column to show.
\item \textbf{Date and time}, when \texttt{display-time-mode} is on (the default): formatted by \texttt{display-time-format} (default \texttt{"\%H:\%M"}), using the same \texttt{\%<FORMAT>} subset as \hyperref[sec:org83b5238]{Capture templates}. It updates at least every 30 seconds.
\item \textbf{Clock}, when a clock is running anywhere in the document: the elapsed time of this clock-in in brackets, then the heading's title, like \texttt{[⏱ 0:35] Write the quarterly report}. It shows the current session, not the task's lifetime total.
\end{itemize}

\textbf{The minibuffer}, below it, is the one strip for the status line, god-mode's chord echo and the search box (see \hyperref[sec:orgf0231a2]{Searching}). Only one shows at a time: the search box while search is open, then the god-mode echo while a chord is in progress or has just completed (the chord and its command's name; see \hyperref[sec:orgb59a778]{God-mode}), and otherwise the last status message. A status message clears at your next keypress or tap, as Emacs's echo area clears at each new command.

Two variables control it, set in Global or Local Variables or in Settings → Quick Settings → \textbf{Modeline}:

\begin{center}
\begin{tabular}{lll}
Variable & Default & Effect\\[0pt]
\hline
\texttt{display-time-mode} & \texttt{1} (on) & Emacs's minor-mode toggle, with its \texttt{1} and \texttt{0} convention instead of \texttt{t} and \texttt{nil}. \texttt{0} hides the date and time.\\[0pt]
\texttt{display-time-format} & \texttt{"\%H:\%M"} & Emacs's variable and default. Any \texttt{\%<FORMAT>} token that \hyperref[sec:org83b5238]{Capture templates} support works.\\[0pt]
\end{tabular}
\end{center}

\noindent\rule{\textwidth}{0.5pt}

\subsection{Links and images}
\label{sec:org01b1fbc}
\begin{itemize}
\item \texttt{[[target][description]]} links render as tappable text.
\item \textbf{Bare URLs} link too: \texttt{https://example.com} typed in a sentence, or wrapped as \texttt{<https://example.com>}. The recognized schemes are \texttt{http}, \texttt{https}, \texttt{ftp}, \texttt{mailto}, \texttt{doi}, \texttt{file}, \texttt{github}, \texttt{webdav} and \texttt{local}, so a time like \texttt{10:30} or a ratio like \texttt{3:4} is never read as a link. Trailing sentence punctuation (\texttt{.}, \texttt{,}, \texttt{)}) is left out: \texttt{see https://example.com/page.} links the URL without the period. Wrap a URL in \texttt{<...>} or \texttt{[[...]]} if it must end in one of those characters.
\item \texttt{doi:10.1145/1327452.1327492} resolves to \texttt{https://doi.org/...} and opens like any external link, bracketed or bare.
\item \textbf{Internal links} resolve by heading title (\texttt{[[*Some Heading]]}), by \texttt{CUSTOM\_ID} (\texttt{[[\#my-id]]}), or by a title search for bare text. Tapping one expands the ancestors and scrolls to the target. Set a \texttt{CUSTOM\_ID} in the general editor (see \hyperref[sec:org22c076e]{Editing your outline}).
\item \textbf{ID links}: \texttt{[[id:3f2a-…][Meeting notes]]} go to the heading with that \texttt{:ID:} property, wherever it now lives. Tapping one looks in the open file, then the other open tabs, then your \hyperref[sec:org2e9fd13]{Agenda Files} (the files org-id scans), then the GitHub, WebDAV and local files you have opened or saved here before (a remembered table, like \texttt{org-id-locations}), and opens whichever holds it; the ⏎ button comes back. If nothing has the ID the status line says so, and if several headings share it the first is used. \textbf{Copy link to heading} (\texttt{org-store-link}, \texttt{C-c l}, \texttt{c SPC l}) puts \texttt{[[id:…][Title]]} on the clipboard, first giving the heading an \texttt{:ID:} if it has none. \textbf{Add ID to heading} (\texttt{org-id-get-create}) only adds the \texttt{:ID:}. \textbf{Find links to this heading} searches this file and your Agenda Files for \texttt{[[id:…}.
\item \textbf{Targets}: \texttt{<{}<{}name>{}>{}} is a spot a link can go to, and the same words in triple angle brackets make a \textbf{radio target}: every plain use of the same words in the file (any case, any spacing, whole words only) becomes a link to it. A link whose text is the target's name goes to the heading that holds it, after trying a heading with that title. Both are shown dimmed as written. Code blocks are not read, so a noweb \texttt{<{}<{}ref>{}>{}} there is left alone. HTML export turns them into anchors and links.
\item \textbf{The ⏎ button}: after you follow a link, a floating ⏎ button appears at bottom-left to go back, so a document with a table of contents (Docs' own included) is navigable on a touch device without a browser Back. Chained links unwind one hop at a time, and the button disappears when there is nowhere left to go back to. It is low-opacity at rest and uses the theme's text colours.
\item \textbf{Image links}: \texttt{http://} and \texttt{https://} image links render inline when \texttt{\#+STARTUP: inlineimages} is set (off by default).
\item \textbf{Local and relative image paths} on GitHub and WebDAV render too: \texttt{[[photo.png]]}, \texttt{[[images/photo.png]]}, \texttt{[[file:photo.png]]} and \texttt{[[github:images/photo.png]]} resolve relative to the document's path (a bare filename lands beside it), under the same \texttt{\#+STARTUP: inlineimages} gate, and are cached for the session. Local files (File System Access) and iOS imports show an \texttt{[image: path]} placeholder instead, because reading a sibling file needs browser permission that can't be granted mid-render.
\item \textbf{\texttt{file:}, \texttt{github:}, \texttt{webdav:} and \texttt{local:} links} parse into a scheme, a path and an optional in-file target, using org's \texttt{::} convention: an absolute path (\texttt{file:/home/user/documents/notes.org}), a relative one (\texttt{file:projects/todo.org}), a headline (\texttt{file:\textasciitilde{}/notes.org::*Project Alpha}) or a text search (\texttt{file:\textasciitilde{}/notes.org::exact phrase}). Tapping one opens the document, as File → Open does (with its unsaved-changes handling), and jumps to the headline or match. A path with no scheme uses the backend the current document came from; an explicit \texttt{github:}, \texttt{webdav:} or \texttt{local:} always uses that backend (\texttt{local:} is a file on this device: see \hyperref[sec:org3d41625]{Android app}).

The back button follows these links across documents: it returns to the heading the link was in, with its ancestors and body expanded, found by the heading's title path (a scroll offset wouldn't survive a re-parse). It falls back to the remembered scroll position when the link had no heading (it was in the preamble) or the heading has since been renamed or removed. This works only when both documents are on GitHub or WebDAV. A local file or iOS import can't be reopened from a link tap or the back button, because opening a path needs a file-picker gesture the browser requires for each file, so a clear message is shown instead.
\item \textbf{\texttt{attachment:} links} are org-attach's link type, produced by \textbf{Attach} (see \hyperref[sec:org7b27204]{Attachments}). They are bracketed only (\texttt{[[attachment:photo.jpg]]}); a bare \texttt{attachment:photo.jpg} in text is not a link.
\end{itemize}

\noindent\rule{\textwidth}{0.5pt}

\subsection{Configuration}
\label{sec:org4f4f89f}
Three layers set variables (Emacs Lisp variables or org directives), resolved in org's own precedence: \textbf{Local Variables beats \#+STARTUP beats Global Variables}. A variable set at a higher layer wins for that file, and one left unset falls through to the next layer down.

\begin{enumerate}
\item \textbf{Global Variables} (\hyperref[sec:orga7a5637]{below}): an app-wide baseline, set in Settings.
\item \textbf{\texttt{\#+STARTUP:}} (\hyperref[sec:orga77e6fc]{below}): a file's own directive. Only three variables have a \texttt{\#+STARTUP:} shorthand, as in real org:
\end{enumerate}

\begin{center}
\begin{tabular}{ll}
Variable & \texttt{\#+STARTUP:} keyword(s)\\[0pt]
\hline
\texttt{org-log-done} & \texttt{logdone} (→ \texttt{time}), \texttt{lognotedone} (→ \texttt{note})\\[0pt]
\texttt{org-startup-folded} & \texttt{overview}, \texttt{content}, \texttt{showall}, \texttt{showeverything}, \texttt{show2levels}, \texttt{show3levels}, \texttt{show4levels}, \texttt{show5levels}\\[0pt]
\texttt{org-startup-with-inline-images} & \texttt{inlineimages} (→ \texttt{t}), \texttt{noinlineimages} (→ \texttt{nil})\\[0pt]
\end{tabular}
\end{center}

\begin{enumerate}
\item \textbf{Local Variables} (\hyperref[sec:org75c5aa1]{below}): a file's own \texttt{\# Local Variables:} block, a general Emacs mechanism processed last.
\end{enumerate}

\noindent\rule{\textwidth}{0.5pt}

\subsubsection{Global Variables}
\label{sec:orga7a5637}
Set in \textbf{Settings → Global Variables}. This is the lowest layer, and the only one not stored in a \texttt{.org} file. The format is one \texttt{name: value} per line, like a \hyperref[sec:org75c5aa1]{Local Variables} block without its wrapper:

\begin{verbatim}
org-log-done: 'time
org-agenda-start-on-weekday: 1
\end{verbatim}

Every Local Variables variable works here too, at app-wide precedence. Changes take effect on \textbf{Save}, with no reload, and they are part of \textbf{Export Settings} and \textbf{Import Settings…} (see \hyperref[sec:org5584f85]{File management}).

\textbf{Smart quotes are normalized}: iOS and many Android keyboards replace a typed straight quote with a curly one, but quote-delimited values here (\texttt{org-xx-extra-menu}, \texttt{org-xx-menu-aliases} and the \hyperref[sec:org83b5238]{Capture Templates} JSON) need the ASCII character. Quick Settings, this text box and Capture Templates convert curly quotes back to straight ones on save. Quotes in ordinary prose elsewhere are untouched.

\noindent\rule{\textwidth}{0.5pt}

\subsubsection{Folding and \texttt{\#+STARTUP}}
\label{sec:orga77e6fc}
org-pwa applies \texttt{\#+STARTUP:} directives, typically on the first line. They can share a line or be spread across several, and the last matching keyword wins.

\begin{center}
\begin{tabular}{lll}
Category & Keywords & Default\\[0pt]
\hline
Heading visibility & \texttt{overview}, \texttt{content}, \texttt{showall}, \texttt{showeverything}, \texttt{show2levels}, \texttt{show3levels}, \texttt{show4levels}, \texttt{show5levels} & \texttt{showeverything}\\[0pt]
Inline images & \texttt{inlineimages}, \texttt{noinlineimages} & \texttt{noinlineimages}\\[0pt]
Progress logging on DONE & \texttt{logdone}, \texttt{lognotedone} & neither (no logging)\\[0pt]
\end{tabular}
\end{center}

\begin{itemize}
\item \texttt{overview}: only top-level headings, everything else folded.
\item \texttt{show2levels} to \texttt{show5levels}: like \texttt{overview} with that cutoff depth. Headings are shown through that level with no body text at any visible level, and deeper levels are folded.
\item \texttt{content}: every heading line shown, but body text (paragraphs, lists, tables) and properties hidden until you open a heading.
\item \texttt{showall}: body text, lists and tables shown. Property drawers (\texttt{:PROPERTIES:}) and blocks (\texttt{\#+BEGIN\_SRC}, \texttt{\#+BEGIN\_QUOTE} and the like) stay folded, one line each.
\item \texttt{showeverything}: \texttt{showall} plus drawers and blocks fully expanded.
\item \texttt{logdone}, \texttt{lognotedone}: set \texttt{org-log-done} to \texttt{time} or \texttt{note} (see \hyperref[sec:org8f91341]{Progress logging}).
\end{itemize}

Cycling a heading to fully expanded (slide-left, or tapping through the fold states) always reveals its properties and blocks, whatever \texttt{\#+STARTUP:} says, and collapsing it hides them again. The defaults match org's own: with no \texttt{\#+STARTUP} line, everything is shown. Archived-heading cycling is a separate mechanism (see \hyperref[sec:org75c5aa1]{Local Variables}).

\noindent\rule{\textwidth}{0.5pt}

\subsubsection{Local Variables}
\label{sec:org75c5aa1}
Two forms, both general Emacs mechanisms: a first-line \texttt{-*- var: val; -*-} cookie, and a \texttt{\# Local Variables: ... \# End:} block, conventionally at the end of the file. Both can appear in one file, and the block wins a conflict. Most variables below also have a labeled control in Settings → \textbf{Quick Settings} (see \hyperref[sec:orgbb54b74]{Settings}) for the app-wide baseline.

\begin{verbatim}
# -*- buffer-read-only: t; -*-
\end{verbatim}

\begin{verbatim}
# Local Variables:
# org-agenda-start-on-weekday: 1
# org-cycle-open-archived-trees: nil
# org-agenda-skip-comment-trees: t
# org-agenda-skip-archived-trees: t
# End:
\end{verbatim}

The recognized variables, alphabetically. Those marked "(own)" are this app's extensions, not real org variables. Booleans are Lisp \texttt{t} or \texttt{nil}, not JavaScript truthiness: the string \texttt{"true"} is not \texttt{t}.

\begin{center}
\begin{tabular}{lll}
Variable & Default & What it does\\[0pt]
\hline
\texttt{buffer-read-only} & \texttt{nil} & Whether the buffer is read-only (see \hyperref[sec:org996c4cc]{Keybindings}: \texttt{C-x C-q}, or \textbf{Toggle read-only} in the command palette)\\[0pt]
\texttt{calendar-latitude} & \texttt{35.994} & Latitude for the solar forms and \texttt{org-weather} (Durham, NC)\\[0pt]
\texttt{calendar-longitude} & \texttt{-78.8986} & Longitude for the same (Durham, NC)\\[0pt]
\texttt{org-agenda-files} & unset (current file only) & Semicolon-separated \texttt{scheme:path} entries: more files for the Agenda and TODO views (see \hyperref[sec:org2e9fd13]{Agenda Files})\\[0pt]
\texttt{org-agenda-skip-archived-trees} & \texttt{t} & Exclude archived headings (tagged \texttt{:ARCHIVE:}) from the agenda; \texttt{nil} includes them\\[0pt]
\texttt{org-cite-global-bibliography} & unset & Bibliography files used by every document, as a Lisp list of paths (see \hyperref[sec:org261744f]{Citations})\\[0pt]
\texttt{org-cite-basic-sorting-field} & \texttt{author} & Field the reference list is sorted by\\[0pt]
\texttt{org-cite-basic-author-year-separator} & \textasciitilde{}, \textasciitilde{} & Text between author and year in a citation\\[0pt]
\texttt{org-footnote-section} & \texttt{Footnotes} & Heading that holds footnote definitions; \texttt{nil} puts each at the end of its own section (see \hyperref[sec:orga27d7d4]{Footnote commands})\\[0pt]
\texttt{org-footnote-auto-label} & \texttt{t} & Label for a new footnote: \texttt{t} next number, \texttt{confirm} next number you can edit, \texttt{nil} ask, \texttt{random}\\[0pt]
\texttt{org-footnote-define-inline} & \texttt{nil} & Write new footnotes as \texttt{[fn:label:text]} at the reference\\[0pt]
\texttt{org-footnote-auto-adjust} & \texttt{nil} & After adding or deleting a footnote: \texttt{renumber} fn:N, \texttt{sort} the definitions, both (\texttt{t}) or neither\\[0pt]
\texttt{org-agenda-skip-comment-trees} & \texttt{t} & Exclude commented headings (\texttt{** COMMENT draft}) from the agenda. They stay visible in the outline; \texttt{nil} includes them\\[0pt]
\texttt{org-xx-agenda-hide-title-stamps} & \texttt{t} & Leave active timestamps out of a heading's title in the agenda (\texttt{Do task <2026-10-11 Sun +14d>} shows as \texttt{Do task}; each day of a range shows \texttt{(day 2/3)}). \texttt{nil} shows the title as written\\[0pt]
\texttt{org-agenda-start-on-weekday} & \texttt{1} (Monday) & The Week view's first day: \texttt{0} is Sunday through \texttt{6} Saturday. An out-of-range value falls back to Monday\\[0pt]
\texttt{org-ascii-text-width} & \texttt{72} & Wrap width, in characters, for ASCII export\\[0pt]
\texttt{org-ascii-charset} & \texttt{ascii} & Characters used by ASCII export: \texttt{ascii}, \texttt{latin1} or \texttt{utf-8}\\[0pt]
\texttt{org-closed-keep-when-no-todo} & \texttt{nil} & \texttt{t} keeps a CLOSED timestamp when a DONE heading is cycled back to no keyword. Cycling to another keyword always removes it (see \hyperref[sec:org8f91341]{Progress logging})\\[0pt]
\texttt{org-contacts-birthday-property} & \texttt{BIRTHDAY} & The property \texttt{org-contacts-anniversaries} reads, as a plain string (\texttt{\# org-contacts-birthday-property: EVENT})\\[0pt]
\texttt{org-contacts-files} & unset (current file only) & More files for \hyperref[sec:org053e4f4]{Contacts (.vcf)} export (see \hyperref[sec:orgba4514f]{Contacts Files})\\[0pt]
\texttt{org-agenda-text-search-extra-files} & unset & More files for text search only, in the same syntax as \texttt{org-agenda-files} (see \hyperref[sec:orgdf17868]{Search extra files})\\[0pt]
\texttt{org-cycle-open-archived-trees} & \texttt{nil} & \texttt{nil}: an archived heading starts folded whatever \texttt{\#+STARTUP:} says, and the swipe gesture refuses to open it (the chevron still does, as org's force-open commands do). \texttt{t} makes archived headings behave like any other\\[0pt]
\texttt{org-deadline-warning-days} & \texttt{14} & How many days before its date a DEADLINE with no warning suffix starts showing in the agenda\\[0pt]
\texttt{org-global-properties} & unset & Properties that apply everywhere. Only \texttt{Effort\_ALL} is used, as in \texttt{Effort\_ALL: 0:15 0:30 1:00 2:00} (see \hyperref[sec:orgf7daf41]{Effort\textsubscript{ALL}})\\[0pt]
\texttt{org-log-done} & \texttt{nil} (no logging) & \texttt{'time} inserts a \texttt{CLOSED:} timestamp when a heading becomes DONE; \texttt{'note} also asks for a note, kept in a \texttt{:LOGBOOK:} drawer. Its \texttt{\#+STARTUP:} shorthands are \texttt{logdone} and \texttt{lognotedone} (see \hyperref[sec:org8f91341]{Progress logging} for the per-keyword \texttt{!} and \texttt{@} forms)\\[0pt]
\texttt{org-refile-targets} & current file, level 1 & Which headings Refile offers, such as \texttt{current maxlevel=3; notes.org tag=work level=2; agenda-files todo=NEXT} (see \hyperref[sec:org6433005]{Refile})\\[0pt]
\texttt{org-startup-folded} & \texttt{showeverything} & The initial visibility: \texttt{'overview}, \texttt{'content}, \texttt{'showall}, \texttt{'showeverything} or \texttt{'show2levels} to \texttt{'show5levels} (the quote is optional); see \hyperref[sec:orga77e6fc]{Folding}\\[0pt]
\texttt{org-startup-with-inline-images} & \texttt{nil} & Whether images render inline when a file opens\\[0pt]
\texttt{org-table-duration-hour-zero-padding} & \texttt{t} & Whether the \texttt{T} and \texttt{U} duration flags zero-pad the hours (see \hyperref[sec:org1458240]{Table formulas})\\[0pt]
\texttt{org-use-property-inheritance} & \texttt{nil} & \texttt{t} lets a \texttt{key:value} search filter also match a value inherited from an ancestor (a heading's own value wins). A \texttt{t} or \texttt{nil} simplification of org's fuller value space\\[0pt]
\texttt{org-use-sub-superscripts} & \texttt{t} & \texttt{t}, \texttt{\{\}} or \texttt{nil} (three values, as in org); see \hyperref[sec:org1a88f72]{Inline text markup}. An unrecognized value falls back to \texttt{t}\\[0pt]
\texttt{org-use-tag-inheritance} & \texttt{t} & \texttt{t} lets a \texttt{tag:} search filter match a tag inherited from an ancestor; \texttt{nil} limits it to the heading's own tags; a list such as \texttt{("work" "home")} or a quoted regexp such as \texttt{"\textasciicircum{}proj\_"} inherits only those tags. Applies to search, the Agenda tag filters and the tags shown on a heading\\[0pt]
\texttt{org-tags-exclude-from-inheritance} & (none) & Tags that never inherit, as a list (\texttt{("private")}) or just words (\texttt{private archive})\\[0pt]
\texttt{org-agenda-show-inherited-tags} & \texttt{t} & Whether Agenda and TODO list rows show inherited tags, dimmed after the title\\[0pt]
\texttt{org-xx-show-inherited-tags} & \texttt{t} & (own) Whether the outline shows inherited tags, dimmed with a dashed border, after a heading's own\\[0pt]
\texttt{org-weather-format} & \texttt{"Weather: \%desc, \%tcur(\%tmin-\%tmax)\%tu, \%p\%pu, \%h\%hu, \%s\%su, \%a\%au"} & (own) The template \texttt{\%\%(org-weather)} fills in (see \hyperref[sec:org58e8c79]{Diary-sexp forms})\\[0pt]
\texttt{org-weather-refresh-interval} & \texttt{never} & (own) \texttt{never}, or an automatic refresh every \texttt{1}, \texttt{2}, \texttt{4}, \texttt{6}, \texttt{12} or \texttt{24} hours\\[0pt]
\texttt{org-weather-speed-unit} & \texttt{mph} & (own) \texttt{km/h}, \texttt{m/s}, \texttt{mph} or \texttt{Knots}\\[0pt]
\texttt{org-weather-temperature-unit} & \texttt{°F} & (own) \texttt{°C} or \texttt{°F}\\[0pt]
\texttt{org-html-head} & unset & A quoted string added to the \texttt{<head>} of every HTML export; a file's own \texttt{\#+HTML\_HEAD:} is added after it. Settings → Export sets it\\[0pt]
\texttt{org-html-preamble} & unset (none) & A quoted format string shown above the content of every HTML export (\texttt{\%t} title, \texttt{\%s} subtitle, \texttt{\%a} author, \texttt{\%d} date, \texttt{\%e} email). A file's own \texttt{\#+OPTIONS: html-preamble:} wins\\[0pt]
\texttt{org-html-postamble} & \texttt{auto} & A quoted format string for the foot of every HTML export, or \texttt{nil} for none. A file's own \texttt{\#+OPTIONS: html-postamble:} wins\\[0pt]
\texttt{org-xx-extra-menu} & unset (button hidden) & (own) The floating extras (☰) button's menu (see \hyperref[sec:org08ac3df]{Extras menu})\\[0pt]
\texttt{org-babel-load-languages} & \texttt{((emacs-lisp . t))} & \texttt{(js . t)} in the list lets \texttt{js} source blocks run (see \hyperref[sec:org41ed666]{Source blocks}); read from Global Variables only, never from a file; also Settings → Quick Settings → \textbf{Editing}\\[0pt]
\texttt{org-xx-extensions} & \texttt{off} & (own) \texttt{on} lets the approved init script run (see \hyperref[sec:org34d11ec]{Extensions}); read from Global Variables only, never from a file; also Settings → Quick Settings → \textbf{Editing}\\[0pt]
\texttt{org-xx-drag-handle} & \texttt{right} & (own) Where the drag grip sits on a heading row: \texttt{right}, \texttt{left} or \texttt{off}; also Settings → Quick Settings → \textbf{Keyboard}\\[0pt]
\texttt{org-xx-god-mode-button} & \texttt{t} (shown) & (own) Whether the floating \textbf{[g]} god-mode button is shown (see \hyperref[sec:org7b13dfa]{God-mode without a hardware keyboard}); also Settings → Quick Settings → \textbf{Keyboard}\\[0pt]
\texttt{org-xx-menu-aliases} & unset (no relabeling) & (own) Relabel or hide buttons across the File, More, Export, Import and View menus (see \hyperref[sec:orgfca7ef3]{Menu customization})\\[0pt]
\texttt{org-xx-startup-with-show-photos} & \texttt{nil} & (own) \texttt{t} shows a contact's photo as a thumbnail beside its field when a file opens: next to \texttt{:PHOTO:} in Flat style, in the "Photo" sub-heading's drawer in Tree style. It follows that field's visibility, so a collapsed contact or folded drawer shows no photo. It has no \texttt{\#+STARTUP:} keyword (see \hyperref[sec:org053e4f4]{Contacts (.vcf) export})\\[0pt]
\texttt{solar-ampm} & \texttt{nil} (24-hour) & (own) \texttt{t} uses 12-hour am/pm for the solar forms\\[0pt]
\texttt{solar-hide-label} & \texttt{nil} (label shown) & (own) \texttt{t} drops the trailing label ("Sunrise", "daylight") and shows the bare time or duration\\[0pt]
\end{tabular}
\end{center}

Any other \texttt{\# key: value} line is parsed and ignored, so new variables can be added later without changing the format.

\noindent\rule{\textwidth}{0.5pt}

\subsection{Progress logging}
\label{sec:org8f91341}
Two related mechanisms write into the same place: a heading's \texttt{:LOGBOOK:} drawer, org's running history of state changes and notes, kept apart from the body.

\subsubsection{\texttt{org-log-done}}
\label{sec:orgb667252}

Set by \texttt{\#+STARTUP: logdone} or \texttt{\#+STARTUP: lognotedone} (see \hyperref[sec:orga77e6fc]{Folding and \#+STARTUP}), it controls what happens when a heading enters a DONE-type state (tapping the TODO badge to cycle there, or the action menu toggle). The two modes are mutually exclusive:

\begin{itemize}
\item \texttt{logdone} (\texttt{time}) inserts a \texttt{CLOSED: [timestamp]} planning line when a heading becomes DONE. Cycling it to a non-done keyword removes the timestamp, and cycling it back to no keyword also removes it unless \texttt{org-closed-keep-when-no-todo} is set (see \hyperref[sec:org75c5aa1]{Local Variables}). Re-marking a done heading with a \textbf{different} done keyword (DONE to KILL) doesn't add or disturb a timestamp.
\item \texttt{lognotedone} (\texttt{note}) prompts for a note in a small panel when a heading becomes DONE. The badge flips at once and never waits for the note, so \textbf{Skip} only declines to add one. A saved note becomes a timestamped \texttt{:LOGBOOK:} entry.
\end{itemize}

They apply to a plain DONE-type keyword with no logging spec of its own, and a file has one or the other as its default. With neither, which is the default as in org, entering DONE inserts nothing extra.

\subsubsection{The per-keyword logging spec}
\label{sec:orgb424b11}

The more general mechanism applies to any keyword, not only done ones, and overrides \texttt{org-log-done} for that keyword. It uses the parenthesized fast-key syntax in \texttt{\#+TODO:}:

\begin{verbatim}
#+TODO: TODO(t) WAIT(w@/!) | DONE(d!) KILL(k@)
\end{verbatim}

Four markers, each optional, as in org:

\begin{center}
\begin{tabular}{ll}
Marker & Meaning\\[0pt]
\hline
\texttt{!} & Log a timestamp when \textbf{entering} this state\\[0pt]
\texttt{@} & Log a note (prompted for, same panel as \texttt{org-log-done}'s own) when \textbf{entering} this state\\[0pt]
\texttt{@/!} & Log a note when entering, \emph{and} a timestamp when \textbf{leaving} — conditionally, see below\\[0pt]
\texttt{/!} & Log a timestamp only when \textbf{leaving} this state — no logging on entering it at all\\[0pt]
\end{tabular}
\end{center}

A leaving-log (\texttt{@/!} or \texttt{/!}) fires only if the target state doesn't already log something on entry. In the example, WAIT to DONE fires nothing extra from WAIT's \texttt{/!}, because DONE already logs a timestamp on entry (\texttt{d!}). WAIT to TODO does fire it, since TODO logs nothing of its own and the move would otherwise go unrecorded.

An explicit per-keyword spec always beats \texttt{org-log-done}, which only fills in for a done-type keyword with no spec. Both can fire for one transition (a \texttt{CLOSED} timestamp and a \texttt{:LOGBOOK:} entry), which can look redundant but is org's behavior.

\textbf{Fast-access keys} (the letter before any \texttt{@} or \texttt{!}, as in \texttt{TODO(t)}) are parsed, so a \texttt{\#+TODO:} line using them doesn't break TODO recognition. A \texttt{\#+TODO:} line with no \texttt{|} treats its last keyword as the only done state, as in org.

\textbf{Parallel workflows}: several \texttt{\#+TODO:} lines in one file (\texttt{\#+TODO: TODO(t) | DONE(d)} then \texttt{\#+TODO: REPORT(r) BUG(b) KNOWNCAUSE(k) | FIXED(f)}) define coexisting workflows, as in org. Tapping the badge cycles within the sequence a heading's keyword already belongs to. For a heading with no TODO state, the badge, the action menu's \textbf{TODO} button and \texttt{C-c C-t} all behave alike: with one workflow they cycle as usual, and with more than one they open a picker showing every state of every workflow, one row per sequence, so you can jump to any state. Each sequence's first state (where a fresh heading starts) is green and its done-type states are red, since a workflow can wrap onto several lines and the last button isn't always its end. \textbf{Clear} resets the heading to blank, and \textbf{Cancel} dismisses the picker unchanged. \texttt{\#+SEQ\_TODO:} and \texttt{\#+TYP\_TODO:} are read the same way as \texttt{\#+TODO:}.

\subsubsection{The \texttt{:LOGBOOK:} drawer}
\label{sec:orgbf80fd9}

Shown under a heading like \texttt{:PROPERTIES:}, with the same visibility rule (hidden by default unless \texttt{\#+STARTUP: showeverything} or the heading is swiped fully open) and the same read-only, tap-for-the-action-menu treatment. A state change shows the transition and timestamp (\texttt{WAIT → DONE   [timestamp]}), and a note appears indented in italics beneath it. Entries follow org's default log format, including leaving out the "from" clause when there was no previous state and continuing multi-line notes with a trailing backslash. \texttt{CLOCK:} lines display the same way (see \hyperref[sec:orgb9e46b3]{Clocking}): \texttt{Clock: start–end ⇒ duration} once closed, and \texttt{Clock: start (running)} while open.

\subsubsection{Agenda: Log}
\label{sec:org6e300bb}

The Agenda's date-range line has a \textbf{Log} checkbox, off by default, like org's \texttt{org-agenda-log-mode}. On, every \texttt{:LOGBOOK:} entry across the aggregated files appears as its own dated item, interleaved by timestamp with SCHEDULED, DEADLINE and anniversaries, with a 📝 icon and any note beneath. It shows a DONE heading's history even though the Agenda normally excludes DONE headings, since seeing what got done is the point of Log mode.

\noindent\rule{\textwidth}{0.5pt}

\subsection{Archiving}
\label{sec:orgf0de015}
Tap a heading's title, then \textbf{Archive}, to move its subtree to the archive location (real org's \texttt{org-archive-subtree}, \texttt{C-c C-x C-s} or \texttt{C-c \$}). A panel opens in two parts, each scrolling on its own: \textbf{Archive \& details} names the destination and what will be added, and \textbf{Refile \& details} lists places to move it instead (see \hyperref[sec:org6433005]{Refile}). \textbf{Archive}, \textbf{Refile} and \textbf{Cancel} stay at the bottom; \textbf{Cancel} has focus, so a stray tap does nothing. On a screen 600 px wide or more the panel, like the other popups, is a centered dialog instead of a bottom sheet; \textbf{Esc} closes it, \textbf{Tab} stays inside it, and focus returns to where you were.

The moved heading gets an \texttt{:ARCHIVE:} tag and properties recording where it came from: \texttt{ARCHIVE\_TIME}, \texttt{ARCHIVE\_FILE}, \texttt{ARCHIVE\_OLPATH}, \texttt{ARCHIVE\_CATEGORY}, \texttt{ARCHIVE\_TODO} and \texttt{ARCHIVE\_ITAGS}, as in org's \texttt{org-archive-save-context-info}. The original is removed only after the copy is safely written, so a failure leaves it where it was.

The location is \texttt{org-archive-location}: the heading's own \texttt{ARCHIVE} property, else the file's \texttt{\#+ARCHIVE:} keyword, else the default \texttt{\%s\_archive::}. It has two parts joined by \texttt{::}: a file (\texttt{\%s} is the current file's name; empty means the current file) and an optional heading to file entries under, where the stars give its level (\texttt{* Archived Tasks}).

\begin{center}
\begin{tabular}{ll}
Location & Behavior\\[0pt]
\hline
\texttt{\%s\_archive::} (default) & Sibling file \texttt{notes.org\_archive}, top-level entries\\[0pt]
\texttt{::* Archived Tasks} & Same file, filed under a top-level \texttt{* Archived Tasks} heading (created if it doesn't exist yet)\\[0pt]
\texttt{\textasciitilde{}/org/archive.org::} & One central file, top-level entries\\[0pt]
\texttt{\textasciitilde{}/org/archive.org::* \%s} & One central file, filed under a heading named after the source file\\[0pt]
\end{tabular}
\end{center}

Archiving to a different file works on GitHub and WebDAV only; a local or iOS-imported file needs a file picker that the browser won't show mid-archive, and the app says so. The heading goes into that file's buffer, which is saved like any other (see \hyperref[sec:org5584f85]{File management}). Archiving within the current file works everywhere. Not supported: \texttt{datetree/} targets and multi-level heading paths.

\subsubsection{Restoring (unarchive)}
\label{sec:org1a89165}
Open the archive file (\textbf{File → Open} lists \texttt{\_archive} files), tap the archived heading, then \textbf{Unarchive}. The same panel opens with \textbf{Restore} in place of \textbf{Archive}, naming the destination; the heading then moves back to its original file and place, loses \texttt{:ARCHIVE:} and the \texttt{ARCHIVE\_*} properties, and gets its TODO state back from \texttt{ARCHIVE\_TODO}. As with archiving, it leaves the archive only after the other file's buffer has taken it.

A heading tagged \texttt{:ARCHIVE:} by hand, with no recorded \texttt{ARCHIVE\_FILE}, just has the tag and properties stripped in place. In the command palette it is \textbf{Unarchive (restore)} (\texttt{org-unarchive-subtree}). There is no separate list for browsing the archive.

\subsubsection{Refile}
\label{sec:org6433005}
Tap a heading's title, then \textbf{Archive}, pick a target in the lower part of the panel and tap \textbf{Refile}, to move the subtree elsewhere in the outline or into another file (real org's \texttt{org-refile}, also in the command palette). The list shows each target by its full outline path, with a search box above it (every word you type must match), with your last five destinations under \textbf{RECENT} and the rest under \textbf{ALL TARGETS}. The subtree becomes the target's last child, with its levels adjusted. A subtree is never offered as its own destination.

Targets come from \texttt{org-refile-targets} (see \hyperref[sec:org75c5aa1]{Local Variables} and \hyperref[sec:org4f4f89f]{Configuration}). The default is the current file, top-level headings only. Entries are separated by \texttt{;}, each a file followed by criteria:

\begin{verbatim}
org-refile-targets: current maxlevel=3; notes.org tag=work level=2; agenda-files todo=NEXT
\end{verbatim}

\begin{itemize}
\item \textbf{File}: \texttt{current}, \texttt{agenda-files}, or a file name (a bare name is a sibling of the current file), optionally with a \texttt{github:}, \texttt{webdav:} or \texttt{local:} prefix, like \texttt{local:inbox.org level=1}.
\item \textbf{Criteria}: \texttt{level=N}, \texttt{maxlevel=N}, \texttt{tag=NAME} (the heading's own tag; inheritance is not considered), \texttt{todo=KEYWORD}, and \texttt{regexp=PATTERN} (matched against the whole heading line, so it must come last; write a literal \texttt{;} as \texttt{\textbackslash{};}). Several criteria on one entry must all match, which goes beyond real org.
\end{itemize}

An entry the app can't understand is skipped and listed with the reason at the top of the picker. Refiling to another file works on GitHub and WebDAV, and the subtree goes into that file's buffer; a local file must have been opened once, except in the \hyperref[sec:org3d41625]{Android app}, where a file in the org-pwa folder works. Not built: creating a new heading as the destination.

\noindent\rule{\textwidth}{0.5pt}


\subsection{Extensions}
\label{sec:org34d11ec}
An init script can add your own diary functions. Set \textbf{Run extension scripts} to \texttt{On} in Quick Settings, paste the script under Settings → \textbf{Extensions}, and tap \textbf{Approve and run}. The script stays on this device (it is not synced, exported or read from any file), and editing it withdraws the approval.

\texttt{org.sexp(name, fn)}: makes \texttt{\%\%(name args)} work in an agenda file. \texttt{fn} is called once per day with \texttt{ctx} (\texttt{date}, \texttt{year}, \texttt{month}, \texttt{day}, \texttt{weekday} with 0 for Sunday, \texttt{today}, \texttt{calendarLatitude}, \texttt{calendarLongitude}) and then your arguments as plain values. Return \texttt{true} for an entry with the line's own text, a string for an entry with that text, anything else for none.

\begin{verbatim}
org.sexp("every-other-friday", (ctx, start) => {
  const days = (Date.parse(ctx.date) - Date.parse(start)) / 86400000;
  return days >= 0 && days % 14 === 0;
});
\end{verbatim}

\begin{verbatim}
* Pay day
%%(every-other-friday "2026-01-02") Pay day
\end{verbatim}

\texttt{org.command(id, label, fn)}: adds \textbf{label} to the command palette. \texttt{fn(ctx)} gets \texttt{ctx.heading}, the heading the command was run on (or null), and can read the whole file through \texttt{org.document}: its \texttt{name} and its \texttt{headings}, each with \texttt{index}, \texttt{level}, \texttt{title}, \texttt{todo}, \texttt{priority}, \texttt{tags}, \texttt{properties}, \texttt{scheduled}, \texttt{deadline}, \texttt{closed}, \texttt{body} and \texttt{path}. It changes the file with \texttt{org.edit}: \texttt{setTitle}, \texttt{setTodo}, \texttt{setTags}, \texttt{setProperty}, \texttt{deleteProperty} and \texttt{appendBody}, each taking a heading's \texttt{index}. The edits are checked, then applied together as one undo step; none is applied if one is invalid, the buffer is read-only or the file changed while the command ran. \texttt{org.notify(text)}, or a string returned from \texttt{fn}, shows a message in the status line.

\begin{verbatim}
org.command("tag-urgent", "Tag as urgent", (ctx) => {
  if (!ctx.heading) return "Put the cursor on a heading first";
  org.edit.setTags(ctx.heading.index, [...ctx.heading.tags, "urgent"]);
});
\end{verbatim}

\texttt{org.on(event, fn)}: runs \texttt{fn(ctx)} when something happens. The events are \texttt{open} (a file finished loading), \texttt{save} (a save succeeded), \texttt{todo-change} (a heading's TODO state changed; \texttt{ctx.heading}, \texttt{ctx.from} and \texttt{ctx.to}) and \texttt{capture} (a capture was filed; \texttt{ctx.description}, \texttt{ctx.file}). A hook can use \texttt{org.notify} in any event, and \texttt{org.edit} after \texttt{todo-change} and \texttt{capture}; its edits are one separate undo step, and do not set hooks off again. Edits after \texttt{save} or \texttt{open} are ignored, since they would mark the file modified the moment it was saved or opened. A hook that fails is logged and never blocks saving or opening.

\begin{verbatim}
org.on("todo-change", (ctx) => {
  if (ctx.to === "DONE") org.edit.setProperty(ctx.heading.index, "FINISHED", new Date().toISOString().slice(0, 10));
});
\end{verbatim}

A name is letters, digits and dashes, and \texttt{org.sexp} cannot take a standard form such as \texttt{diary-float}. The script runs in the same sandbox as \texttt{js} blocks. A diary function that takes longer than 3 seconds, a command longer than 5, or a hook longer than 2, stops the script. Emacs does not know your functions, so a file that uses one needs the same function defined in Emacs Lisp to work there. Not built: the rest of the planned API.

\subsubsection{Table functions, link types and exports}
\label{sec:org4c39a2c}
\begin{itemize}
\item \texttt{org.tableFunction(name, fn)}: use \texttt{name(...)} in a \texttt{\#+TBLFM:} formula. Arguments are numbers or text; return a number or text. Run with the table's Calc button or \emph{Recalculate all tables}.
\item \texttt{org.linkType(prefix, fn)}: makes \texttt{[[prefix:path][label]]} clickable. \texttt{fn(path)} returns an address (https, mailto, tel or geo) that is opened.
\item \texttt{org.exportBackend(id, label, fn)}: adds \emph{Export: label} to the palette. \texttt{fn(ctx)} gets \texttt{ctx.document} and returns text, or \texttt{\{ text, filename, mime \}} (up to 2 MB).
\end{itemize}

\begin{verbatim}
org.tableFunction("twice", (n) => n * 2);
org.linkType("weather", (place) => "https://wttr.in/" + encodeURIComponent(place));
org.exportBackend("titles", "Titles only", (ctx) => ctx.document.headings.map((h) => h.title).join("\n"));
\end{verbatim}

\subsubsection{Agenda sources and header lines}
\label{sec:org3553688}
\begin{itemize}
\item \texttt{org.agenda.source(name, fn)}: \texttt{fn(\{ from, to \})} returns items for the agenda, each \texttt{\{ title, start, end, repeat \}} with \texttt{start} as \texttt{"2026-12-25"} or \texttt{"2026-12-25 09:30"} and \texttt{repeat} such as \texttt{"+1w"}. Asked again after ten minutes.
\item \texttt{await org.ics.parse(text)}: the events of a calendar feed as \texttt{\{ summary, start, end, allDay, location, repeat, cancelled \}}. Complex recurrences are not expanded.
\item \texttt{org.ui.line(id, fn)}: \texttt{fn()} returns one line of text shown above the agenda.
\end{itemize}

\begin{verbatim}
org.agenda.source("holidays", async () => {
  const r = await org.fetch("https://example.org/us-holidays.ics");
  return org.ics.parse(await r.text());
});
org.ui.line("mood", () => "Have a good day");
\end{verbatim}

\subsubsection{Network, position and cache}
\label{sec:orgb9cd593}
The sandbox has no network, but a script or a \texttt{js} block can ask the page:
\begin{itemize}
\item \texttt{await org.fetch(url, \{ method, headers, body \})} returns \texttt{\{ status, ok, url, text(), json() \}}. Only https (or http to this device), GET or POST, the headers Accept, Content-Type, Authorization and X-Api-Key, a body up to 100 KB, no cookies, 8 seconds, up to 1 MB back, 30 requests per run. A redirect to a host you did not allow is refused.
\item \texttt{await org.location.get()} returns \texttt{\{ lat, lon, accuracy, time \}}, rounded to about 10 m.
\item \texttt{await org.cache.get(key, maxAgeMs)} returns a stored value or \texttt{null}; \texttt{await org.cache.set(key, value)} stores one (200 KB each, 1 MB and 100 keys per script or document) and returns it.
\end{itemize}
You are asked before the first use. A block lists what it needs: \texttt{\#+BEGIN\_SRC js :net api.open-meteo.com :location yes} (several hosts separated by spaces; \texttt{*.example.org} covers subdomains) and you are asked once per document. The init script is asked per host and for the position when it first uses them. The question warns that the code can send anything it can read, including the open file, to the site, so allow only code you trust.

\noindent\rule{\textwidth}{0.5pt}

\subsection{Source blocks}
\label{sec:org41ed666}
\texttt{js} blocks run in place, as in Org Babel. They are off until you set \textbf{Run JavaScript source blocks} to \texttt{On} in Settings → Quick Settings → \textbf{Editing}. The setting comes from your Global Variables, never from a file, so a document cannot turn it on for itself.

Each \texttt{js} block then has a \textbf{Run} button; \texttt{C-c C-c} runs the block under the cursor, and \textbf{Run source block} is in the command palette. The result is written under the block as \texttt{\#+RESULTS:} and replaced on the next run. One Undo removes it.

\begin{verbatim}
#+NAME: sales
| q  |  n |
|----+----|
| Q1 | 10 |
| Q2 | 20 |

#+BEGIN_SRC js :var data=sales
return data.slice(1).map(([q, n]) => [q, n * 2]);
#+END_SRC

#+RESULTS:
| Q1 | 20 |
| Q2 | 40 |
\end{verbatim}

\begin{itemize}
\item \texttt{:results}: \texttt{value} (the default: what the block \texttt{return\textasciitilde{}s) or \textasciitilde{}output} (the \texttt{console.log} lines), with \texttt{table}, \texttt{list}, \texttt{raw} or \texttt{silent}.
\item \texttt{:var name=value}: a number, a quoted string, or the \texttt{\#+NAME:} of a table, which arrives as an array of rows.
\item \texttt{:timeout N}: seconds, 5 by default and 60 at most. A block still running then is stopped.
\item \texttt{:wrap NAME}: wraps the result in a \texttt{\#+begin\_NAME} block.
\end{itemize}

The code runs in a sandbox with no access to the page, your files or your accounts, and with no network unless the header asks for it and you allow it (see Network, position and cache). \texttt{org.vars.get("calendar-latitude")} reads a setting. The first run in a saved document asks for confirmation; an unsaved one asks every time.


\subsubsection{Literate programming: header arguments, Noweb and Tangle}
\label{sec:org229bca3}
Header arguments are inherited as in Org: a \texttt{\#+PROPERTY: header-args :results output} line applies to the whole file, \texttt{header-args:js} to \texttt{js} blocks only, and a \texttt{:header-args:} property on a heading applies to its subtree (\texttt{header-args+} adds to what is inherited instead of replacing it). A block's own arguments, and its \texttt{\#+HEADER:} lines, win. This applies to every argument, including \texttt{:results}, \texttt{:var}, \texttt{:net}, \texttt{:timeout} and \texttt{:wrap}.

\textbf{Noweb}: with \texttt{:noweb yes}, \texttt{<{}<{}name>{}>{}} in a block is replaced by the block \texttt{\#+NAME: name}, by all blocks with \texttt{:noweb-ref name} (joined with \texttt{:noweb-sep}, a newline by default), or by the text of the heading whose \texttt{CUSTOM\_ID} is \texttt{name}. Text before the reference on its line is put before every line of the expansion, so \texttt{// <{}<{}name>{}>{}} comments out all of it. \texttt{:noweb} also takes \texttt{tangle}, \texttt{no-export} (expand to run and tangle, show the references on export), \texttt{strip-export} (references removed on export) and \texttt{strip-tangle}. \texttt{yes} expands when running, tangling and exporting. A reference that matches nothing, or that goes round in a circle, is an error rather than empty text.

\textbf{Tangle}: the palette command \textbf{Tangle this file} (\texttt{org-babel-tangle}) writes the blocks with \texttt{:tangle} to files: \texttt{:tangle yes} names the file after the document, \texttt{:tangle path/file.js} names it, \texttt{no} skips it. Blocks for the same file are joined in order, separated by a blank line (\texttt{:padline no} removes it). Also supported: \texttt{:shebang} (and the file mode that goes with it), \texttt{:tangle-mode}, \texttt{:comments link}, \texttt{org}, \texttt{both} and \texttt{noweb} (which also wraps each \texttt{<{}<{}name>{}>{}} chunk in comments naming the block it came from), \texttt{:prologue}, \texttt{:epilogue}, \texttt{:no-expand}, and \texttt{:var} for \texttt{js} blocks (written as \texttt{var x=3;}). \texttt{COMMENT} and \texttt{:ARCHIVE:} subtrees are skipped. One file is saved as it is; several, or one in a folder, come as a zip with the file modes set. Paths may not be absolute or climb out with \texttt{..}. The output was checked against Emacs's own.

\textbf{Cache}: \texttt{:cache yes} skips a block whose code, variables and \texttt{:results} shape are unchanged since its last run; the hash is kept in its results line (\texttt{\#+RESULTS[hash]:}). Delete the results to run it again. Do not cache a block that reads the time or a live feed.

\textbf{Export} shows what \texttt{:exports} asks: \texttt{code} (the default; stored results are not shown), \texttt{results}, \texttt{both} or \texttt{none`. Nothing is run on export; \textasciitilde{}results} shows what \texttt{\#+RESULTS:} already holds. This applies to ASCII, HTML, Markdown, ODT and Org export.

\textbf{The init script from an Org file}: in Settings → Extensions, \textbf{Use an Org file…} takes \texttt{github:path}, \texttt{webdav:path} or \texttt{local:name}, then \texttt{::target} (default \texttt{::init.js}). The blocks of that file that tangle to the target become the init script, Noweb included. Only the file you name is read; the approval covers the assembled text, so a change to the file asks again.

\noindent\rule{\textwidth}{0.5pt}

\subsection{Attachments}
\label{sec:org7b27204}
Files kept with a heading, as in Org's \texttt{org-attach}. Press \texttt{C-c C-a} on a heading, or tap \textbf{Attach} in 📋 Edit details, for a grid of buttons:

\begin{itemize}
\item \textbf{Attach} (\texttt{a}): on a phone or tablet, opens the system chooser, where you pick a file or a photo, or take a photo or a video with the camera; the file is copied into the heading's folder (in a browser the chooser takes several files at once). On ChromeOS and desktop, whose file chooser has no camera, it first asks \textbf{Choose a file}, \textbf{Take a photo} or \textbf{Record a video}; the last two open an in-app webcam panel (take or record, then Attach or Retake).
\item \textbf{Record audio} (\texttt{r}): record from the microphone, play it back, then save.
\item \textbf{Attach open document} (\texttt{b}, \texttt{org-attach-buffer}): save one of your open documents, as text, in the folder.
\item \textbf{New text file} (\texttt{n}, \texttt{org-attach-new}): give a name and type its contents.
\item \textbf{Open} (\texttt{o}): show a picture, recording, video or text file inside the app; anything else opens in the browser's viewer, or downloads when there is none.
\item \textbf{Folder} (\texttt{f}, \texttt{org-attach-reveal}): the folder as a compact list, one row per file: a checkbox to select, then \texttt{O} (open), \texttt{S} (save to disk) and \texttt{D} (delete), then the name. The strip above selects all, saves or deletes the selection (deleting asks once), and adds files.
\item \textbf{Delete one} (\texttt{d}) and \textbf{Delete all} (\texttt{D}): remove files after a confirmation.
\item \textbf{Set DIR} (\texttt{s}) and \textbf{Unset DIR} (\texttt{S}): keep a heading's files in a folder of your choosing instead. Unset DIR asks whether to move the files to the ID folder, delete them, or leave them.
\item \textbf{Sync} (\texttt{z}, \texttt{org-attach-sync}): set the \texttt{:ATTACH:} tag to match what is in the folder, for files added or removed by other means.
\end{itemize}

Each heading's files live in \texttt{data/<first 2 characters of its ID>/<rest of the ID>/} beside the document, and the heading gets an \texttt{ID} property and an \texttt{:ATTACH:} tag the first time. A \texttt{DIR} property names another folder, relative to the document (a leading \texttt{/} starts at the root of the storage); it wins over the ID. The attachments of a heading are what is in that folder, so a file added there by other means shows up too. With several files, Open, Preview, Save and Delete ask which one.

Example: on a heading, \texttt{C-c C-a}, then \texttt{a}, pick \texttt{receipt.pdf}. The heading now carries \texttt{:ATTACH:} and \texttt{receipt.pdf} is in its folder.

\textbf{Links in the body}: Settings → Attachments chooses whether attaching also adds an \texttt{[[attachment:name]]} link: for pictures, audio and video (the default), always, or never, as in Org. A bare link to a picture shows it inline; audio gets a play button. Tapping an \texttt{attachment:} link does the same as \textbf{Open}. A link finds its folder from the \texttt{ID} of its heading, else the nearest ancestor's.

\textbf{Requirements}: GitHub or WebDAV must be connected. The \hyperref[sec:org3d41625]{Android app} instead keeps a local file's attachments in the org-pwa folder you choose once.

\textbf{Large files}: over 25MB asks before attaching, since the whole file is held in memory and uploaded in one go; GitHub refuses over 95MB (its limit is 100MB), so use WebDAV for long videos.

\textbf{Filenames}: path-unsafe characters become \texttt{\_}; a name already in the folder gets \texttt{-1}, \texttt{-2} and so on before the extension, so repeated camera captures named \texttt{image.jpg} all survive.

\textbf{Not built}: Org's copy/move/link variants (\texttt{c}, \texttt{m}, \texttt{l}) and absolute \texttt{DIR} paths outside the storage.

\noindent\rule{\textwidth}{0.5pt}

\subsection{Clocking}
\label{sec:orgb9e46b3}
Every clocking action is a \hyperref[sec:orgc687081]{command palette} command: \textbf{Clock in}, \textbf{Clock out}, \textbf{Cancel clock}, \textbf{Continue last clock}, \textbf{Go to clocked task} and \textbf{Set effort estimate}, found by name or by their real Org name (\texttt{org-clock-in} and so on). The app runs one clock at a time, though a hand-edited file with several running \texttt{CLOCK:} lines round-trips fine.

\begin{itemize}
\item \textbf{org-clock-in} inserts a bare \texttt{CLOCK: [timestamp]} line on the heading whose action menu is open, else the keyboard-focused one (a status message says so when there is neither). On the heading already clocked it does nothing. On a different heading it clocks the running one out at the same moment, with no gap.
\item \textbf{org-clock-out} turns the line into \texttt{CLOCK: [start]-{}-{}[end] =>  H:MM}, with the duration computed. \textbf{org-clock-cancel} removes the running \texttt{CLOCK:} line entirely and records no duration, for a clock started by mistake. Both act on whichever clock is running anywhere in the document; with none running they do nothing and say so.
\item \textbf{org-clock-in-last} (\texttt{C-c C-x C-x}; the palette's \textbf{Continue last clock}, \texttt{c x x} in god-mode) resumes clocking on the most recently clocked heading, anywhere in the document. A status message covers "a clock is already running" and "nothing has been clocked yet". In an Extras entry it is \texttt{'org-clock-in-last} (see \hyperref[sec:org08ac3df]{Extras Menu}).
\item \textbf{org-clock-goto} (\texttt{C-c C-x C-j}; the palette's \textbf{Go to clocked task}, \texttt{c x j}) jumps to the clocked-in heading. There is one clock across all open tabs, so it works from anywhere: it switches to the clock's tab, unfolds whatever hides the task, and leaves the floating back button behind. With no clock running the palette shows it dimmed ("no clock is running"). With \texttt{C-u} (\texttt{u c x j} in god-mode, or the palette's \textbf{Go to a recently clocked task}, named \texttt{org-clock-select-task}) it lists recently clocked tasks from every open tab, most recent first, and jumps to the one you choose.
\end{itemize}

Clock lines live in a \texttt{:LOGBOOK:} drawer from the first one. Org itself keeps the first one bare (see \hyperref[sec:org8f91341]{Progress logging}).

\textbf{Showing clocked time}:

\begin{itemize}
\item \textbf{org-clock-display}: a \textbf{Clock} checkbox on \textbf{View → TODO} (like Agenda's \textbf{Log}; see \hyperref[sec:org105cb3a]{Agenda}) shows each item's total clocked time beside its title: all time ever clocked on the task and its descendants, including a running clock's elapsed time. It is off by default and is a session toggle, not a variable.
\item \textbf{org-clock-report}: a \texttt{\#+BEGIN: clocktable ... \#+END:} block, computed for a date range and shown read-only under \textbf{View → TODO}'s \textbf{Clock report} (see \hyperref[sec:orgf6942c0]{The TODO view}).
\item While a clock runs, the modeline (see \hyperref[sec:org6cade45]{Modeline and minibuffer}) shows \texttt{[⏱ H:MM] Heading title}, refreshed every 30 seconds.
\end{itemize}

\textbf{Effort estimates}: \textbf{Set effort estimate} in the command palette (\texttt{C-c C-x e}, \texttt{org-set-effort}) prompts for the \texttt{EFFORT} of the heading whose action menu is open, else the focused one. Any \texttt{org-duration} form is accepted and stored as typed: \texttt{1:30} or \texttt{1:23:45}; units like \texttt{2h}, \texttt{45min}, \texttt{1d 3h} or \texttt{1y 3d 3h 4min} (spaces optional); a mix like \texttt{3d 13:35}; decimals like \texttt{2.35h}; or a bare number, which is minutes. The units are Emacs's defaults: \texttt{min}, \texttt{h}, \texttt{d} (24 hours), \texttt{w} (7 days), \texttt{m} (a \textbf{month} of 30 days, not minutes: write \texttt{30min} for half an hour) and \texttt{y} (365.25 days). A blank value clears the estimate, and an invalid one is refused with a status message and the old value kept. You can also type it into \textbf{Edit details → Properties}.

An \texttt{EFFORT} shows beside the title as a \texttt{⏱} badge, as typed, and tapping the badge opens the same prompt. While that heading is clocked in, the modeline shows the running time against it, \texttt{[⏱ 0:45 / 2:00] Heading title} (the estimate as \texttt{H:MM}), with a trailing \texttt{!} once the time passes the estimate.

Not built: custom duration units (\texttt{org-duration-units}), the \texttt{Effort\_ALL+} append form, and a reminder tied to the estimate beyond the modeline's \texttt{!}. See \hyperref[sec:org2d2d41c]{Differences from Emacs org-mode}.

\noindent\rule{\textwidth}{0.5pt}

\subsubsection{Effort\textsubscript{ALL}: allowed effort values}
\label{sec:orgf7daf41}
Org's effort prompt completes against a list of allowed values, \texttt{Effort\_ALL}. It looks in this order, and the first source that gives at least one usable value wins:

\begin{enumerate}
\item an \texttt{Effort\_ALL} property on the heading itself, or on the nearest ancestor that has one;
\item the file's \texttt{\#+PROPERTY: Effort\_ALL 0 0:10 0:30 1:00 2:00} line;
\item the \textbf{Global properties} setting (\texttt{org-global-properties}, in Settings → Advanced or a file's Local Variables): semicolon-separated \texttt{NAME: value} entries, such as \texttt{Effort\_ALL: 0:15 0:30 1:00 2:00 4:00}. Only \texttt{Effort\_ALL} is used.
\end{enumerate}

Values are space-separated, and each must be a valid \texttt{org-duration} (see \hyperref[sec:orgb9e46b3]{Clocking}); anything else is skipped. When a list applies, \textbf{Set effort estimate} shows its values as buttons above the text box: one tap saves that value, and the box still takes any other duration. With no list the prompt is plain free text. The same list supplies the Agenda's effort filter (see \hyperref[sec:org105cb3a]{Agenda}); with none defined, the filter uses org's own fallback set (\texttt{0 0:10 0:30 1:00 2:00 3:00 4:00 5:00 6:00 7:00}).

Not supported: \texttt{Effort\_ALL+}, and quoted values containing spaces.


\subsection{Narrowing}
\label{sec:org77ba459}
A heading's action menu has a \textbf{🔽 Narrow} action that restricts the outline to that heading and its subtree: every sibling, ancestor and unrelated heading disappears from view. On the narrowed heading it becomes \textbf{🔼 Widen}, while every other heading still shows Narrow, and tapping it re-narrows to that heading instead. A banner at the top, \textbf{Narrowed to: TITLE — Widen}, stays while narrowing is active. Tapping \textbf{Widen} there, or on the heading's own action menu, shows the whole file again. These are org's \texttt{C-x n s} and \texttt{C-x n w}, also available in \hyperref[sec:orgb59a778]{god-mode} as \texttt{x SPC n s} and \texttt{x SPC n w}.

Narrowing only restricts what is displayed; the file is untouched. \textbf{Save, Export, Agenda and Search all still work on the whole document.} Narrowing follows the heading it points at: promote, demote, move or rename it and the narrowed view updates, and if it is deleted the view widens automatically. Narrowing to another heading replaces the restriction, with no stack of levels to pop, as in org. Folding is independent of narrowing.

\textbf{Narrowing persists} across switching tabs and across closing and reopening the app: a file you left narrowed reopens narrowed to the same heading.

\textbf{Text view while narrowed} shows only the narrowed subtree's raw text, so you can edit one section as plain org syntax (see \hyperref[sec:orga9cae38]{The plain-text editor}). You reach it only by being narrowed in the outline first; switching to Text without narrowing shows the whole file.

\noindent\rule{\textwidth}{0.5pt}

\subsection{The plain-text editor}
\label{sec:orga9cae38}
Tap \textbf{View → Text} to replace the outline with one full-width textarea holding the whole document as raw org text: everything, including \texttt{\#+STARTUP:}, \texttt{\#+TODO:}, a \hyperref[sec:org75c5aa1]{Local Variables} block, and syntax the structured view has no UI for (editing block content, changing the \texttt{\#+TODO:} sequence itself). Tap \textbf{View → Org} to switch back. The text is reparsed at that point, so changes (to \texttt{\#+STARTUP}, \texttt{\#+TODO} or Local Variables, say) take effect at once, mark the document modified, and persist like any other edit. This is the escape hatch for what the tap-driven UI doesn't cover.

\textbf{Position}: switching in lands near where you last navigated, not always at the top. After a search result, an internal link or an agenda item, the cursor and scroll position land at that heading's line. Manual scrolling isn't tracked, so scrolling around and then switching opens at the top. The scroll position is proportional to the line's place in the file, so it is near the line, not exact. While \hyperref[sec:org77ba459]{narrowed}, the textarea starts at the top of the narrowed text.

\textbf{Narrowed}: only the narrowed subtree's text is editable. Switching back to Org splices the box into its original place in the full document. The lines it replaces are fixed when you enter narrowed Text view, so retyping the heading's asterisks to promote or demote it, renaming it, or adding a new heading inside the fragment are ordinary edits that can't touch anything outside the region. Narrowing follows a level or title change made this way: rename or re-level the heading you're narrowed to, and the narrowed view updates to match.

\textbf{Select All} and \textbf{Copy All} buttons sit above the textarea, because selecting a long document with the phone's drag handles is painful. \textbf{Select All} selects the full text, and the device's own Copy still works from there. \textbf{Copy All} writes it straight to the clipboard, and if the clipboard API is unavailable or denied it selects the text instead. \texttt{Ctrl-A} or \texttt{Cmd-A} works natively inside the textarea, so the buttons are for devices without a keyboard. While narrowed, both act on just the narrowed text.

\textbf{Search works here without leaving Text view}, searching and replacing the textarea's live raw text. It is a separate implementation from Org view's structured \hyperref[sec:orgf0231a2]{Searching}, which reads a parsed document and so can never safely edit property drawers, tags, block content and the like. Text view has no such restriction: a value inside a \texttt{\#+BEGIN\_VERSE} or \texttt{\#+BEGIN\_QUOTE} block (an imported contact's note, say) can be found and replaced here. \textbf{Regex} and \textbf{Replace} work as in Org view, including the interactive \texttt{y}, \texttt{n}, \texttt{q}, \texttt{.} and \texttt{!} walk. \textbf{Match} mode isn't offered, since it needs a parsed document. A replacement stays in the textarea until you leave (switch views, or Save), like every edit made here.

\noindent\rule{\textwidth}{0.5pt}

\subsection{Agenda}
\label{sec:org105cb3a}
Tap \textbf{View → Agenda} for everything with a date, in Day, Week or Month view. One button cycles the three views. \textbf{‹} and \textbf{›} step back or forward by the current unit, \textbf{Today} jumps to today, and a horizontal swipe does the same as ‹ and ›. Tapping an item returns to the Org view at that item.

\textbf{Row actions}: press and hold an item for \textbf{Mark done}, \textbf{Set state}, \textbf{Reschedule}, \textbf{Priority}, \textbf{Refile} and \textbf{Open heading}. Swipe an item to the right to mark it done; \textbf{Undo} stays for a few seconds, and a repeating item moves to its next date. They change the open file only: an item from another agenda file offers \textbf{Open heading} alone.

\textbf{Filter chips}: \textbf{Tag}, \textbf{State} and \textbf{Category} under the controls keep the items that match; a chip shows its choice and an ✕ clears it. They match a heading's own tags, and last for the session.

\textbf{Titles}: a timestamp written in a heading's title is left out of its row, since the day and the time column already say it. Set \texttt{org-xx-agenda-hide-title-stamps} to \texttt{nil} to show titles as written.

\textbf{Empty days}: every day shows by default, even when nothing is on it. The \textbf{g} button toggles \texttt{org-agenda-show-all-dates} for the session; it is bold while empty days are hidden. The starting value comes from \texttt{org-agenda-show-all-dates} in \hyperref[sec:org75c5aa1]{Local Variables} or the \textbf{Show empty days in Agenda} Quick Setting.

\textbf{Weeks}: Week view always shows a full calendar week, Monday to Sunday by default (\texttt{org-agenda-start-on-weekday}), whichever day you opened it on. Day headers show the ISO week number on the first day of each week (always in Day view), and the range label shows the week numbers, like \texttt{(W36)} or \texttt{(W31-W36)}.

\textbf{Categories}: every item shows a category prefix: the file's \texttt{\#+CATEGORY:} keyword, else its file name without directory or \texttt{.org}. Put the keyword at the top of the file:

\begin{verbatim}
#+CATEGORY: Projects
* TODO Review weekly goals
  SCHEDULED: <2026-09-07 Mon>
\end{verbatim}

\textbf{Times}: a timestamp can carry a range, like \texttt{SCHEDULED: <2026-08-31 Mon 08:00-09:00>}. Week and Month views show \texttt{08:00-09:00}; a single time shows as \texttt{08:00 ……}; an untimed SCHEDULED or DEADLINE shows its \texttt{Scheduled:} or \texttt{Deadline:} label instead.

\textbf{Time grid}: Day view adds an hour-by-hour grid (org's \texttt{org-agenda-time-grid} default: dashed lines every two hours, 8am to 8pm). Untimed items come first; timed items sit between the lines at their time. Week and Month views have no grid.

\textbf{Effort}: an item with an \texttt{EFFORT} property (any \texttt{org-duration} form, see \hyperref[sec:orgb9e46b3]{Clocking}) shows it after its title, as typed: \texttt{⏱ 1:30}. Each day header shows the day's total on the right (\texttt{⏱ 3:30}), with a trailing \texttt{+2?} for TODOs that have no estimate yet. The \hyperref[sec:orgf6942c0]{TODO view} shows \texttt{⏱} too. How it counts:

\begin{itemize}
\item A heading shown twice in one day (scheduled and deadline) counts once. An overdue task carried forward counts toward today.
\item An appointment counts when it has an effort. One with neither a TODO keyword nor an effort is not counted in \texttt{+N?}.
\item An explicit \texttt{0:00} is an estimate. An invalid effort counts as missing.
\item Logged entries, diary lines and sunrise/sunset lines never count. A day with no estimates shows no total.
\end{itemize}

\textbf{Sorting, filtering and limiting by effort}: the \texttt{⏱} button in the toolbar opens a panel, the touch version of org's \texttt{effort-up} and \texttt{effort-down}, \texttt{org-agenda-filter-by-effort} and \texttt{org-agenda-max-effort}:

\begin{itemize}
\item \textbf{Sort}: least or most effort first within each day.
\item \textbf{Filter}: tasks with effort \textbf{at most}, \textbf{at least} or \textbf{exactly} a value from the \hyperref[sec:orgf7daf41]{Effort\textsubscript{ALL}} list (or org's default set).
\item \textbf{Max per day}: keep tasks while the running total stays within, say, \texttt{4:00}.
\item \textbf{Reset} clears all three.
\end{itemize}

The rules are org's own: \textbf{at most} and \textbf{at least} are inclusive; a task with no effort counts as high effort (it sorts last, is dropped by \textbf{at most} and \textbf{exactly}, and kept by \textbf{at least}); the limit keeps a prefix of the sorted day, so once the total passes the maximum everything after it is dropped, and an unestimated task ends the list; entries with no heading behind them (sunrise lines, anniversaries) are never filtered or counted; and the order is sort, limit, filter. These are session settings, not remembered after a reload.

\subsubsection{Dated entry sources}
\label{sec:org082ccb7}
Five kinds of entry show up. They differ in whether they carry forward, as in real org:

\begin{itemize}
\item \texttt{SCHEDULED:} shows on its date and, until the heading is done, on today as overdue ("3 days overdue"), as in Emacs. A delay suffix (\texttt{SCHEDULED: <2026-01-10 Sat -2d>}) hides it on its own date and shows it on today once the delay has passed.
\item \texttt{DEADLINE:} shows on its date, and on today when it is overdue or inside its warning period ("due in 3 days"). A warning suffix (\texttt{DEADLINE: <2026-01-10 Sat -3d>}) sets that period. With no suffix it is \texttt{org-deadline-warning-days} (default 14; see \hyperref[sec:org75c5aa1]{Local Variables}). An explicit suffix overrides that. Overdue and upcoming entries show on today only, not on the days between, and a repeating one counts from its first date.
\item \textbf{Plain active timestamps} anywhere in a heading's entry: its title, body, property values and drawers, such as \texttt{**** Jane Doe <2000-01-01 Sat +1y>}. They are listed beside the heading's own SCHEDULED or DEADLINE, and never carry forward. Inactive \texttt{[...]} timestamps, \texttt{\# \textasciitilde{} comment lines, \textasciitilde{}\#+} keyword lines and source, example, export and comment blocks are ignored.
\item A \textbf{general sexp timestamp} \texttt{<\%\%(sexp)>} in the title follows the same rules; see \hyperref[sec:org36f3daa]{General sexp timestamps}.
\item \texttt{org-contacts-anniversaries} reads yearly events (birthdays, anniversaries) from contact headings, as Emacs's org-contacts does. The contacts are the headings of your \hyperref[sec:orgba4514f]{Contacts Files} (\texttt{org-contacts-files}); with none set, the headings of the files you are looking at. A contact's date is a property, written \texttt{YYYY-MM-DD}, with an optional description after it.
\end{itemize}

\begin{verbatim}
* John Doe
:PROPERTIES:
:BIRTHDAY: 2000-01-01
:END:

* Mary Doe
:PROPERTIES:
:ANNIVERSARY: 2015-06-15
:SPOUSE: John Doe
:END:
\end{verbatim}

A line \texttt{\%\%(org-contacts-anniversaries FIELD FORMAT)} in any agenda file is a switch, not an event. \texttt{FIELD} is the property to read, as a quoted string (\texttt{BIRTHDAY} by default, or \texttt{org-contacts-birthday-property}), and \texttt{FORMAT} says how each entry reads; either may be left out or written \texttt{nil}. Each distinct line is its own scan, so a plain one and one with \texttt{"ANNIVERSARY"} give both, in one file or two, and the same line twice is one:

\begin{verbatim}
%%(org-contacts-anniversaries)
%%(org-contacts-anniversaries "ANNIVERSARY")
%%(org-contacts-anniversaries "ANNIVERSARY" "Anniversary: %h & %s (%Y)")
\end{verbatim}

Entries read \texttt{Birthday: John Doe (26th)}, or \texttt{Anniversary: Mary Doe \& John Doe (11th)} for a contact with a \texttt{:SPOUSE:}, with the age worked out for each year shown. In a \texttt{FORMAT}, \texttt{\%h} and \texttt{\%l} are the name, \texttt{\%y} the age, \texttt{\%Y} the age as an ordinal (\texttt{26th}), \texttt{\%s} the spouse, \texttt{\%n} the name with the spouse, and \texttt{\%\%} a percent sign (\texttt{\%s} and \texttt{\%n} are this app's). \texttt{YYYY} may be \texttt{nil} when the year is unknown, shown as \texttt{(??)}. Text after the date in the property replaces the label (\texttt{:BIRTHDAY: 2000-06-15 Wedding Anniversary} reads \texttt{Wedding Anniversary: Mary \& John Doe (26th)}). A value that can't be parsed is skipped. No \texttt{:EMAIL:} is required.

This works with \hyperref[sec:orgc97ce31]{vCard import}. In Flat style the properties are bare dates. In Tree style they are descendant headings (\texttt{:FIELDTYPE: birthday}, \texttt{anniversary} and \texttt{spouse}), found anywhere in the contact's subtree unless the contact has its own property.

\subsubsection{Diary-sexp forms}
\label{sec:org58e8c79}
These are self-contained lines anywhere in a heading's body, with no switch line needed. They show with a 🔁 icon (☀️ for the solar ones).

\begin{itemize}
\item \texttt{\%\%(org-anniversary YEAR MONTH DAY) text}: applies only to years after YEAR. In the text, \texttt{\%d} becomes the elapsed years and \texttt{\%s} its ordinal suffix, so \texttt{\%\%(org-anniversary 2001 6 12) \%d\%s wedding anniversary} shows "25th wedding anniversary" in 2026. February 29 counts as March 1 in non-leap years.
\item \texttt{\%\%(org-cyclic N YEAR MONTH DAY) text}: every N days from that date, never before it.
\item \texttt{\%\%(diary-float MONTH DAYNAME N [DAY] [YEAR]) text}: the Nth weekday of a month. MONTH is 1 to 12, a list like \texttt{(1 4 7 10)} (a quoted \texttt{'(1 4 7 10)} also works), or \texttt{t} for every month. DAYNAME is 0 (Sunday) to 6. N counts from the start of the month when positive and from the end when negative (\texttt{-1} is the last). \texttt{\%\%(diary-float 11 4 4) Thanksgiving} is the 4th Thursday of November. An optional DAY moves where scanning starts (\texttt{\%\%(diary-float 11 2 1 2) US Election Day} starts at the 2nd, skipping a Tuesday on the 1st), and an optional YEAR limits the entry to that year.
\item \texttt{\%\%(org-block Y1 M1 D1 Y2 M2 D2) text}: every day in an inclusive range, in either order.
\item Lines built from \texttt{and}, \texttt{or}, \texttt{not}, \texttt{when}, \texttt{org-class} or \texttt{org-date}, such as \texttt{\%\%(and (org-block 2026 9 1 2026 12 18) (org-class 2026 9 1 2026 12 18 1)) Seminar}, use the evaluator described under \hyperref[sec:org36f3daa]{General sexp timestamps}. The text after the closing parenthesis is the entry, and a string result replaces it.
\item \texttt{\%\%(diary-sunrise)}, \texttt{diary-sunset}, \texttt{diary-civil-dawn}, \texttt{diary-civil-dusk}, \texttt{diary-nautical-dawn}, \texttt{diary-nautical-dusk}, \texttt{diary-astronomical-dawn} and \texttt{diary-astronomical-dusk} are this app's own. They use \texttt{calendar-latitude} and \texttt{calendar-longitude} (default Durham, NC; set them in Local Variables or with the device-location button in Quick Settings) and show, for example, "06:30 Sunrise", "20:15 Sunset", "06:02 Dawn", "20:43 Dusk" in 24-hour time. Civil, nautical and astronomical twilight are the sun 6°, 12° and 18° below the horizon; \texttt{n/a} replaces the time on a day when it doesn't occur. \texttt{solar-ampm} (default \texttt{nil}) switches to 12-hour ("6:30am Sunrise") and \texttt{solar-hide-label} (default \texttt{nil}) drops the label ("06:30"); set them in Local Variables or Quick Settings (\hyperref[sec:orgbb54b74]{Settings}).
\item \texttt{\%\%(diary-day-length)} shows the day's daylight as "13:32 daylight". \texttt{solar-hide-label} applies and \texttt{solar-ampm} doesn't. Polar day and night show "24:00" and "00:00".
\item \texttt{\%\%(org-weather)} shows current conditions and today's forecast from \href{https://open-meteo.com/}{Open-Meteo} (free, no key), for the same coordinates. It is one entry, on today only, showing the last-fetched snapshot, and nothing until the first refresh.
\end{itemize}

\texttt{org-weather-format} controls the text. The default is \texttt{"Weather: \%desc, \%tcur(\%tmin-\%tmax)\%tu, \%p\%pu, \%h\%hu, \%s\%su, \%a\%au"}, which gives "Weather: Overcast, 86.3(75.5-91.8)°F, 1002.8hPa, 49\%, 5.1mph, 12m". An unrecognized \texttt{\%} sequence is left as is.

\begin{center}
\begin{tabular}{ll}
Placeholder & Meaning\\[0pt]
\hline
\texttt{\%desc}, \texttt{\%icon} & Description and emoji for the current weather code\\[0pt]
\texttt{\%tcur}, \texttt{\%tmin}, \texttt{\%tmax} & Current temperature, today's low and high\\[0pt]
\texttt{\%tamin}, \texttt{\%tamax} & Today's "feels like" low and high\\[0pt]
\texttt{\%tu} & Temperature unit (\texttt{org-weather-temperature-unit}, default \texttt{°F})\\[0pt]
\texttt{\%h}, \texttt{\%hu} & Humidity and its unit\\[0pt]
\texttt{\%p}, \texttt{\%pu} & Surface pressure and its unit (\texttt{hPa})\\[0pt]
\texttt{\%s}, \texttt{\%su} & Wind speed and unit (\texttt{org-weather-speed-unit}: \texttt{mph} by default, \texttt{km/h}, \texttt{m/s} or \texttt{Knots})\\[0pt]
\texttt{\%a}, \texttt{\%au} & Age of the snapshot, in \texttt{m}, \texttt{h} or \texttt{d}\\[0pt]
\end{tabular}
\end{center}

All three variables can be edited in Quick Settings, with a reset to default.

\textbf{Refreshing}: \textbf{Settings → Weather} shows the last refresh time and a \textbf{Refresh weather now} button, and says whether \texttt{\%\%(org-weather)} is used in the open document, only through \texttt{org-agenda-files}, or nowhere. \texttt{org-weather-refresh-interval} also refreshes automatically: \texttt{never} (the default) or every \texttt{1}, \texttt{2}, \texttt{4}, \texttt{6}, \texttt{12} or \texttt{24} hours, checked at open and periodically, and only while \texttt{\%\%(org-weather)} is in use. A failed refresh keeps the cached data, which persists across restarts.

\textbf{Not supported on purpose}: \texttt{diary-anniversary}, \texttt{diary-cyclic}, \texttt{diary-block} and \texttt{diary-date}. They read dates in an order set by \texttt{calendar-date-style}, while here a date is always year month day, as in their \texttt{org-*} twins. To use a line from an Emacs diary, rename it and put the year first: \texttt{\%\%(diary-anniversary 10 31 1948)} becomes \texttt{\%\%(org-anniversary 1948 10 31)}. A line using the diary names matches nothing.

\subsubsection{General \texttt{<\%\%(sexp)>} timestamps}
\label{sec:org36f3daa}
A \texttt{<\%\%(expression)>} goes where a plain timestamp would, in a heading's title. The expression is evaluated for each day shown, and every day it is non-nil is an occurrence. It is a matcher, not a repeater: marking the heading DONE has no effect, and it can't carry a repeater suffix. It is skipped when the heading has SCHEDULED or DEADLINE, and it never carries forward.

\begin{verbatim}
* Weekly Sunrise <%%(when (org-cyclic 7 2026 8 9) (diary-sunrise))>
* TODO Team Sync <%%(diary-float t 6 1)>
* Water the plants <%%(org-cyclic 3 2026 1 1)>
* Company Anniversary <%%(org-anniversary 2018 5 14)>
\end{verbatim}

The evaluator is small, not a Lisp interpreter:

\begin{itemize}
\item \texttt{(when CONDITION EXPR)}: the result of EXPR if CONDITION is non-nil, otherwise no match.
\item \texttt{(and EXPR...)}, \texttt{(or EXPR...)}, \texttt{(not EXPR)}: as in Lisp. \texttt{and} yields its last value, so a string survives; \texttt{or} yields the first that matches. \texttt{t} and \texttt{nil} work as values, and list arguments may be quoted (\texttt{'(10 11 12)}). \texttt{<\%\%(and (org-block 2026 9 1 2026 12 18) (diary-float t 1 1))>} is the first Monday of each month, within that range.
\item \texttt{(today-p)}: this app's shorthand for \texttt{(equal date (calendar-current-date))}. \texttt{<\%\%(when (today-p) (diary-sunrise))>} shows sunrise only today.
\item \texttt{(org-block Y M D Y M D)} is an inclusive range, in either order. \texttt{(org-date Y M D)} takes, for each field, an integer, a list or \texttt{t}: \texttt{(org-date t '(10 11 12) 22)} is the 22nd of October, November and December every year. Every date is year month day, and impossible dates match nothing.
\item \texttt{(org-class Y1 M1 D1 Y2 M2 D2 DAYNAME [SKIP-WEEK...])}: DAYNAME (0 Sunday to 6) every week between the two dates, except the listed ISO week numbers. Skipping holidays isn't supported; an expression that asks for it matches nothing.
\item \texttt{(org-cyclic ...)}, \texttt{(org-anniversary ...)} and \texttt{(diary-float ...)}, as above, work as sub-expressions. So do the solar functions, \texttt{(diary-day-length)} and \texttt{(org-weather)}, which return a formatted string; a string result becomes the displayed text ("Weekly Sunrise: 06:30 Sunrise"). \texttt{org-contacts-anniversaries} isn't recognized inside a sexp timestamp.
\item \texttt{(format FORMAT-STRING ARG...)} combines string-returning functions into one line: \texttt{<\%\%(when (today-p) (format "Rise \%s (\%s)" (diary-sunrise) (diary-civil-dawn)))>} shows "Rise 06:30 Sunrise (06:02 Dawn)", or "Rise 06:30 (06:02)" with \texttt{solar-hide-label}. Only \texttt{\%s} and \texttt{\%\%} are recognized, and each \texttt{\%s} evaluates its argument, including a nested \texttt{format}. Extra arguments are ignored; too few give no match. A title that is only the timestamp shows the result exactly as built, with no title prefix.
\item Anything else (an unknown function, a malformed expression, a bare value) is no match, never an error.
\end{itemize}

\textbf{Not included}: \texttt{let}, a \texttt{date} variable, arithmetic and comparison, and \texttt{calendar-*} helpers, so the Emacs manual's "pay check" example can't run. Also not implemented: \texttt{diary-remind}; the display-only \texttt{diary-day-of-year}, \texttt{diary-iso-date}, \texttt{diary-julian-date}, \texttt{diary-astro-day-number} and \texttt{diary-lunar-phases}; alternate calendars and Hebrew observances (\texttt{diary-yahrzeit}, \texttt{diary-rosh-hodesh} and the rest); and \texttt{calendar-date-style}.

\subsubsection{Other agenda behavior}
\label{sec:orgba6bdc3}
\begin{itemize}
\item \textbf{Completed items are excluded}, using the done keywords of the file's own \texttt{\#+TODO:} sequence.
\item \textbf{Repeating timestamps expand} for SCHEDULED, DEADLINE and title timestamps: \texttt{<2000-01-01 Sat +1y>} shows every year within the range displayed. A repeating SCHEDULED or DEADLINE shows only on its repeat dates, with no carry-forward or early warning between them.
\item \textbf{Marking a repeating task done shifts it}, as in org: its date moves forward, \texttt{LAST\_REPEAT} is stamped, and the TODO state returns to the first keyword. \texttt{+N} adds one interval to the previous date, even if that is still in the past. \texttt{++N} adds intervals until the date is later than now, keeping the weekday and time (a timed one can land later today). \texttt{.+N} counts from today, or from the current time for hours. The DONE transition is logged normally; the bounce-back is not, and the \texttt{CLOSED} timestamp is removed again.
\item \textbf{Archived and commented headings are excluded}. A commented heading starts with \texttt{COMMENT} (\texttt{** COMMENT draft, not ready yet}), as in Emacs. The settings are \texttt{org-agenda-skip-archived-trees} and \texttt{org-agenda-skip-comment-trees} in \hyperref[sec:org75c5aa1]{Local Variables}.
\end{itemize}

A trailing tag needs a space before it (\texttt{<2000-01-01 Sat +1y> :BDAY:}), as in org. Without the space the tag stays as title text, but the timestamp still shows.

\subsubsection{Habits}
\label{sec:org321c6c2}
A habit is an open TODO with \texttt{:STYLE: habit} and a repeating SCHEDULED date, as in Emacs's \texttt{org-habit}. The agenda row gets a strip of days: 21 before today, today (outlined, \texttt{!}) and 7 after. Blue: not yet due. Green: can be done. Yellow: last chance. Red: overdue. A \texttt{*} is a day it was done.

\begin{verbatim}
* TODO Shave
SCHEDULED: <2026-10-13 Tue .+2d/4d>
:PROPERTIES:
:STYLE: habit
:END:
\end{verbatim}

\texttt{.+2d/4d} means every 2 days, no later than every 4. Set the second number under Details > Dates > "No later than". Marking the habit done records a \texttt{State "DONE"} line in its \texttt{:LOGBOOK:}; the strip is drawn from those lines. A habit shows in today's agenda once it is due. Settings > Agenda has \texttt{org-habit-preceding-days}, \texttt{org-habit-following-days}, \texttt{org-habit-show-habits-only-for-today} and \texttt{org-habit-show-all-today}. The ⏱ panel hides habits for the session.

Tap a habit's strip for its history: how many times in a row it was done on time (each no later than the last-chance day after the one before), the best run, the last 12 weeks as a grid, and the latest completions. Details > Props has a \textbf{Habit} switch that sets \texttt{STYLE: habit}.

\subsubsection{Agenda Files}
\label{sec:org2e9fd13}
Add files in \textbf{Settings → Quick Settings → Agenda files}, or in a document's own \texttt{\# Local Variables:} block to override that for the one file (see \hyperref[sec:org75c5aa1]{Local Variables}). Entries are semicolon-separated \texttt{scheme:path} pairs, like \texttt{github:journal.org;webdav:notes/todo.org}, the same syntax as \texttt{org-refile-targets} (see \hyperref[sec:org6433005]{Refile}). The Agenda and TODO views then scan all of them. \texttt{github} and \texttt{webdav} files are fetched; \texttt{local:name.org} reads a file on this device: one you have opened with \textbf{File → Open → Local file} or, in the \hyperref[sec:org3d41625]{Android app}, one in the org-pwa folder. A malformed entry is skipped. If the open file is in the list, its live in-memory version is used.

Files are fetched once per session. Once any are configured, a \textbf{↻} button in the agenda controls refetches them, and a file that fails to load is reported above the list. A local file is read-only here: Agenda, TODO, Search and Export use it, but Refile, Replace and Capture never write into it. It loads only if the browser already allows access; otherwise the agenda says it "needs permission", and the \textbf{↻} button, being a tap, is what asks. Export's \textbf{This file + Agenda Files} scope always fetches fresh (see \hyperref[sec:org053e4f4]{Export}). Two entries with the same file name, such as \texttt{github:notes.org} and \texttt{local:notes.org}, can't both be used: the first is, and the agenda says so for the second.

\subsubsection{Contacts Files}
\label{sec:orgba4514f}
Set in \textbf{Settings → Quick Settings → Contacts files}, or per document in Local Variables, exactly like \hyperref[sec:org2e9fd13]{Agenda Files} (real org-contacts' \texttt{org-contacts-files}; \texttt{github}, \texttt{webdav} and \texttt{local}). \hyperref[sec:org053e4f4]{Contacts (.vcf)} export's \textbf{This file + Contacts Files} scope scans them. The list is independent of Agenda Files, because the two commonly differ.

\subsubsection{Calendar and contacts sync (CalDAV, CardDAV)}
\label{sec:org66ee6b6}
Mirror your agenda into a calendar on your own server, so its items appear, with alarms, in your phone's calendar app. Create a calendar on the server first (a Radicale or Nextcloud calendar), then set it in \textbf{Settings → Calendar (CalDAV) / Contacts (CardDAV)}: the calendar's address, and a username and password (blank means the WebDAV ones). Leave the address blank to turn it off.

It is one-way: the calendar follows the agenda, and edits made in a calendar app are overwritten, so give it a calendar of its own. The app only ever removes events it created, and an event is removed only when its own file was loaded, so opening a different file never removes another file's events.

What goes in is what \textbf{View → Agenda} shows, from the open file and your \hyperref[sec:org2e9fd13]{Agenda Files}: SCHEDULED, DEADLINE, plain timestamps, diary entries, and birthdays and anniversaries (from your Contacts Files). Each event has its title, its start and end time, the heading's tags as categories, its \texttt{:LOCATION:} and the start of its body (400 characters), and a deadline's warning delay (\texttt{-3d}) becomes an alarm. Every event with a time of day also gets an alarm \texttt{appt-message-warning-time} minutes before it (Settings → Agenda notifications); all-day events don't. A repeating item is one event with a repeat rule. A birthday or anniversary is one event for each year in the window, with that year's age. Completed items, and the agenda's overdue and warning-day entries, are not sent. Items are kept from a month back to six months ahead.

A \texttt{<\%\%(…)>} timestamp in a heading title becomes one all-day event for each day it matches, so \texttt{<\%\%(when (today-p) (org-weather))>} is a single event for today. It reads "Weather as of 10:45" in its description, from the last forecast fetched, and the calendar sync updates that one event whenever the weather refreshes (see \texttt{org-weather-refresh-interval}). Yesterday's is removed. Your location, the weather settings and any script-defined functions (see \hyperref[sec:org34d11ec]{Extensions}) are used just as in the Agenda; a script's functions are asked first and the export waits for the answers. Plain \texttt{\%\%(diary-sunrise)}, \texttt{\%\%(diary-sunset)} and \texttt{\%\%(org-weather)} lines in a body are Agenda-only and are not exported.

It syncs after a save, when you return to the app, and at launch. \textbf{Sync} in Settings does it now, and \textbf{Rebuild calendar} (or the palette's \textbf{Rebuild calendar from the agenda}) removes the app's events and sends them all again. After a wrong address or password, automatic syncing stops until you change the settings or sync by hand.

\textbf{Contacts} work the same way. Create an address book on the server, and set its address as the \textbf{Contacts address}. It uses the same username and password as the calendar, and either address can be left blank. Your \hyperref[sec:orgba4514f]{Contacts Files} are mirrored into it, one vCard per contact, in Flat or Tree style. The open file is not sent unless it is one of them, and a contact is removed only when its own file was loaded. \textbf{Sync} does the calendar and the contacts together, and says how each went. \textbf{Rebuild contacts} (or \textbf{Rebuild contacts from the contacts files}) removes the app's contacts and sends them all again. \textbf{Sync agenda to calendar} and \textbf{Sync contacts to the address book} do one each.

Examples: \texttt{https://dav.example.com/radicale/me/calendar/} and \texttt{https://dav.example.com/radicale/me/contacts/}. Like WebDAV, the server must allow requests from the app (CORS), including \texttt{PUT}, \texttt{DELETE} and \texttt{PROPFIND} with the \texttt{Authorization}, \texttt{Content-Type} and \texttt{Depth} headers.

\subsubsection{Agenda notifications}
\label{sec:org0153874}
Get a reminder shortly before an appointment, as Emacs's \texttt{appt} does. Turn it on in \textbf{Settings → Agenda notifications}; the settings use Emacs's names and belong to this device.

\begin{itemize}
\item \textbf{Agenda notifications} (\texttt{appt-activate}): on or off.
\item \textbf{Warn this many minutes before} (\texttt{appt-message-warning-time}): 12 by default, as in Emacs, adjustable in 1-minute steps.
\item \textbf{Remind again every} (\texttt{appt-display-interval}): 3 minutes by default. Once an appointment is within its warning time, it is announced again at this interval, as in Emacs: 10 minutes warned every 3 reads 10, 7, 4 and 1. A warning time that is a multiple of the interval also ends with a reminder at the start.
\item \textbf{Scan agenda files every} (\texttt{appt-agenda-scan-interval}): 5 minutes by default. The app also looks again at launch, when you return to it, and after a save.
\item \textbf{Countdown in the mode line} (\texttt{appt-display-mode-line}): shows \texttt{Appt: 10m} while an appointment is inside its warning time.
\item \textbf{Show reminders as} (\texttt{appt-display-format}): \textbf{window} is a system notification plus a banner in the app; \textbf{echo} is the status line only.
\end{itemize}

An appointment is an item with a time of day, from the open file and your \hyperref[sec:org2e9fd13]{Agenda Files}: SCHEDULED, DEADLINE or a plain timestamp, repeating ones included. Done items and items with no time of day are ignored. A reminder is titled by the appointment, with the time left and its time below: \texttt{Dentist}, \texttt{in 10 min · 09:30–10:15}. Under that comes the start of the heading's body (up to 400 characters), the same text the calendar sync sends; on Android it shows when the notification is expanded. The banner stays until you tap \textbf{Open day} or \textbf{Dismiss}, and the next reminder replaces it. A system notification stays until you swipe it away, and the repeats replace each other. Tapping a reminder opens that day in the Agenda.

The first time, tap \textbf{Allow notifications}. If they are blocked, reminders still show in the app. In the browser, reminders appear while the app is open. In the \hyperref[sec:org3d41625]{Android app} they are scheduled with the system, so they arrive when the app is closed; allowing \textbf{Exact timing} keeps them on the minute.

\noindent\rule{\textwidth}{0.5pt}

\subsection{The TODO view}
\label{sec:orgf6942c0}
Tap \textbf{View → TODO} for every active TODO-state heading in the file, independent of any date, like org's global TODO list (\texttt{C-c a t}, as distinct from the calendar-style agenda, \texttt{C-c a a}). A TODO with no SCHEDULED, DEADLINE or timestamp never shows in Agenda, so this is where it lives.

It is a flat list in document order, each item showing its TODO keyword, title, category and SCHEDULED or DEADLINE date (overdue in red). The filter chips work as in the Agenda, and \textbf{Sort} switches between document order and date order, with undated tasks last. Tap an item to jump to it in the outline; press and hold or swipe it as in the \hyperref[sec:org105cb3a]{Agenda}. It has the same exclusions as Agenda (completed items, by the file's \texttt{\#+TODO:} sequence, archived headings and commented headings, with the same \hyperref[sec:org75c5aa1]{Local Variables} overrides).

\textbf{Clock report}: a checkbox on the same row as \textbf{Clock} (a rule separates the clock controls from the task list) shows org's \texttt{org-clock-report}, a \texttt{\#+BEGIN: clocktable ... \#+END:} block, rendered read-only instead of written into the document. Checking it shows a start and end date picker (the last 7 days including today the first time, then whatever you pick) and a \texttt{maxlevel} selector (1 to 10, default 2, as in org). The table lists every heading with clocked time in the range, down to \texttt{maxlevel}, and each row's total includes all its descendants. A deeper heading still adds its time to the ancestor shown, without a row of its own, and a heading with no clocked time in its subtree is omitted. A clock that is still running is excluded, because its duration would change with every view. A \textbf{Copy} button, on the same row as the \textbf{Clock} and \textbf{Clock report} checkboxes, copies the rendered text as valid org dynamic-block syntax, ready to paste into a file as a permanent record.

\noindent\rule{\textwidth}{0.5pt}

\subsection{Capture Templates}
\label{sec:org83b5238}
The command palette's \textbf{Capture} command (or \texttt{C-c c}, see \hyperref[sec:orgb59a778]{God-mode}) opens a panel for org's \texttt{org-capture-templates}: filing a note, task or table row into a specific place in the outline without navigating there. Each template is listed by its description. Tap one; if it has \texttt{\%\textasciicircum{}\{Prompt\}} or \texttt{\%?} fields, a form asks for each, with a \textbf{live preview} of the expanded result above the fields and the first field focused with the keyboard up. A template with no prompts captures at once.

After a capture the form returns to the template list, so you can capture again straight away. If a capture fails (a network error, a sync conflict) the form stays as you left it, with your text intact. Every template is also a palette command, \textbf{Capture: <description>}, that captures with that template directly, and the Extras menu (see \hyperref[sec:org08ac3df]{Extras Menu}) can jump to one template by its key.

A capture into a different file goes into that file's buffer, as in Emacs: the entry is added to its tab (opened in the background if it has none), which shows as modified until you save it with \textbf{Save} or \textbf{Save Some Buffers} (see \hyperref[sec:org5584f85]{File management}). If the file has no tab and has unsaved local changes from an earlier session, the capture refuses: open the file, resolve its pending changes, then retry.

Four templates are the default: \texttt{b} (Bullet List), \texttt{c} (Check List), \texttt{m} (Meeting) and \texttt{t} (Table Row). Edit or replace them in \textbf{Settings → Capture Templates}, as a JSON array of template objects.

\textbf{Template shape:}

\begin{verbatim}
{
  "key": "b",
  "description": "Bullet List",
  "type": "item",
  "olp": ["heading 1", "heading n"],
  "template": "%? [The captured text or note]",
  "file": "journal.org",
  "prepend": true,
  "prependHeading": true
}
\end{verbatim}

\begin{itemize}
\item \texttt{key}: must be unique. It is the hotkey in a god-mode picker and shows beside the description on the capture form.
\item \texttt{type}: \texttt{item} (a bullet), \texttt{checkitem} (a checkbox item), \texttt{plain} (raw text parsed as an org fragment; the only type that can contain headings, drawers or TODO keywords) or \texttt{table-line} (a row added to the nearest table under the target, or a new table).
\item \texttt{olp}: the outline path to insert into, found or created heading by heading. A segment wrapped in \texttt{\%<FORMAT>} (like \texttt{"\%<\%Y-\%m>"}) expands dynamically, so a monthly table can target "2026-07" without creating that heading by hand. Other segments are literal.
\item \texttt{file} (optional): capture into a different file (org's \texttt{file+olp}). A bare name is a sibling of the open file, a path containing \texttt{/} is used as is, and an optional \texttt{github:}, \texttt{webdav:} or \texttt{local:} prefix names the backend; a prefix that doesn't match the open document's backend is refused. If it resolves to the open file, the capture lands in that live document. Otherwise the other file is read, changed and written back without being opened, and nothing changes if the write fails. Omit it to capture into the open file. GitHub and WebDAV targets work automatically; a local or iOS-imported target must have been opened or saved once first, except in the \hyperref[sec:org3d41625]{Android app}, where a file in the org-pwa folder works automatically and is created if it doesn't exist.
\item \texttt{template}: the text inserted, after \%-expansion. A JSON \texttt{\textbackslash{}n} is a newline; a blank line after the entry is \texttt{"...\%?\textbackslash{}n"}.
\item \texttt{prepend} (default \texttt{false}): put the new entry at the top instead of the bottom. A new list item goes before the first item, and a new table row goes right after the header row and rule. For \texttt{item}, \texttt{checkitem} and \texttt{plain} it also puts any auto-created \texttt{olp} heading first among its siblings.
\item \texttt{prependHeading} (default \texttt{false}; \texttt{table-line} only): put newly created \texttt{olp} headings first among their siblings, independently of \texttt{prepend}. That keeps the newest month's heading on top while its rows still append in order. It applies at every level that has to be created; existing segments are reused as they are.
\item \texttt{omitEmptyEntries} (default \texttt{false}; \texttt{plain} only): remove entries whose expanded value is empty because a \texttt{\%\textasciicircum{}\{Prompt\}} was left blank. That covers \texttt{:KEY:} drawer entries and body lines shaped \texttt{KEY:value}, such as raw vCard lines (for \texttt{N} and \texttt{ADR}, a value of only whitespace and semicolons counts as empty). It also removes an empty trailing tag (a blank \texttt{:\%\textasciicircum{}\{Tag\}:} would show as \texttt{::}), in every sub-heading the capture defines.
\item \texttt{preText}, \texttt{postText} (optional, default empty): template text, in the same \%-language, added around the output. For \texttt{item}, \texttt{checkitem} and \texttt{plain} it wraps every capture, with no automatic spacing. For \texttt{table-line} it is added once, when the target has no table yet: \texttt{preText}, the new table with this first row, then \texttt{postText}; later captures only add rows. No blank line separates them from the table, so a header row and rule in \texttt{preText} join the table, and a \texttt{\#+TBLFM:} or \texttt{\#+PLOT:} line is picked up (see \hyperref[sec:org1458240]{Table formulas} and \hyperref[sec:org0e13f05]{Plots}). A \texttt{\%\textasciicircum{}\{Prompt\}} in them still appears on the form even when the capture won't use it, unless the target is the open document, a loaded Agenda File or a local file, where the form can check first.
\end{itemize}

\textbf{Template text, the \%-escape language} (local time):

\begin{itemize}
\item \texttt{\%<FORMAT>} is a \texttt{format-time-string} subset. The examples use Wednesday, March 5, 2025, 14:30:07:
\end{itemize}

\begin{center}
\begin{tabular}{lll}
Specifier & Meaning & Example output\\[0pt]
\hline
\texttt{\%a} & Abbreviated weekday name & \texttt{Wed}\\[0pt]
\texttt{\%A} & Full weekday name & \texttt{Wednesday}\\[0pt]
\texttt{\%b} & Abbreviated month name & \texttt{Mar}\\[0pt]
\texttt{\%B} & Full month name & \texttt{March}\\[0pt]
\texttt{\%d} & Day of month, zero-padded (01–31) & \texttt{05}\\[0pt]
\texttt{\%e} & Day of month, space-padded ( 1–31) & \texttt{␣5}\\[0pt]
\texttt{\%F} & ISO date (\texttt{\%Y-\%m-\%d} combined) & \texttt{2025-03-05}\\[0pt]
\texttt{\%H} & Hour, 24-hour, zero-padded (00–23) & \texttt{14}\\[0pt]
\texttt{\%I} & Hour, 12-hour, zero-padded (01–12) & \texttt{02}\\[0pt]
\texttt{\%j} & Day of year, zero-padded to 3 digits (001–366) & \texttt{064}\\[0pt]
\texttt{\%m} & Month, zero-padded (01–12) & \texttt{03}\\[0pt]
\texttt{\%M} & Minute, zero-padded (00–59) & \texttt{30}\\[0pt]
\texttt{\%p} & \texttt{AM} or \texttt{PM} & \texttt{PM}\\[0pt]
\texttt{\%R} & Hour:minute (\texttt{\%H:\%M} combined) & \texttt{14:30}\\[0pt]
\texttt{\%S} & Second, zero-padded (00–59) & \texttt{07}\\[0pt]
\texttt{\%T} & Hour:minute:second (\texttt{\%H:\%M:\%S} combined) & \texttt{14:30:07}\\[0pt]
\texttt{\%y} & 2-digit year & \texttt{25}\\[0pt]
\texttt{\%Y} & 4-digit year & \texttt{2025}\\[0pt]
\texttt{\%z} & Numeric UTC offset (±HHMM); depends on the device's time zone (shown for US Eastern summer time) & \texttt{-0400}\\[0pt]
\texttt{\%Z} & Time zone abbreviation; device-dependent like \texttt{\%z} & \texttt{EDT}\\[0pt]
\texttt{\%\%} & A literal percent sign & \texttt{\%}\\[0pt]
\end{tabular}
\end{center}

An unrecognized specifier is left visible (\texttt{\%Q} stays \texttt{\%Q}).

\begin{itemize}
\item \texttt{\%t}, \texttt{\%T}: an active timestamp, date only or date and time. \texttt{\%u}, \texttt{\%U}: the same, inactive (\texttt{[...]} instead of \texttt{<...>}).
\item \texttt{\%i} is content shared in from another app (see \hyperref[sec:orge0ecaf9]{Capturing from other apps}), and \texttt{\%a} is a link to the shared page, \texttt{[[url][title]]}. Text before \texttt{\%i} on its line is added before every later line, so a \texttt{\%i} that follows indentation keeps every line indented. A shared line that would start a heading or a \texttt{\#+} line gets a leading comma, so shared text can never change the document's structure, and only web and mail addresses become links. Both are empty when nothing was shared.
\item \texttt{\%x} is the clipboard. It is read only for a template that uses it, the browser may ask permission (iOS shows a Paste confirmation), and a refusal leaves it empty while the capture still goes ahead.
\item \texttt{\%\textasciicircum{}\{Prompt\}} asks a question on the in-app form. \texttt{\%\textasciicircum{}\{Prompt|default\}} pre-fills an answer, and \texttt{\%\textasciicircum{}\{Prompt|default|choice1|choice2\}} also lists the choices as a hint under the field (there is no dropdown). Several prompts share one form, and cancelling discards the whole capture.
\item \texttt{\%\textasciicircum{}t}, \texttt{\%\textasciicircum{}T}, \texttt{\%\textasciicircum{}u}, \texttt{\%\textasciicircum{}U} are interactive timestamps: active date (\texttt{<2026-09-22 Tue>}), active date and time, inactive date (\texttt{[2026-09-22 Tue]}), inactive date and time. Active timestamps show in the \hyperref[sec:org105cb3a]{Agenda}; inactive ones don't. Each is a text field with a \textbf{📅 Pick date} (or \textbf{📅 Pick date \& time}) button. In the popup, \textbf{OK} writes the timestamp into the field, where it stays editable, and \textbf{Cancel} leaves the field as it was. The date \& time popup also has an \textbf{End time}, which follows the start by 30 minutes until you change it (\textbf{No end} removes it) and writes a range such as \texttt{<2026-09-22 Tue 09:00-10:30>}: an appointment with a start and an end, exported with both. They count as prompts, in reading order, for \texttt{\%\textbackslash{}N} and the live preview.
\item \texttt{\%\textbackslash{}N} (\texttt{\%\textbackslash{}1}, \texttt{\%\textbackslash{}2}, \ldots{}) re-inserts the answer already given to the Nth \texttt{\%\textasciicircum{}\{Prompt\}} or \texttt{\%?} (1-based) instead of asking again. Out of range gives empty text. In the JSON text the backslash must be doubled: \texttt{\%\textbackslash{}\textbackslash{}1}.
\item \texttt{@\#} is the table-formula row number, for \texttt{table-line} captures only (empty otherwise). It takes an offset, \texttt{@\# + n} or \texttt{@\# - n}, with n a non-negative integer.
\item \texttt{\%?} is org's cursor marker, treated as one more prompt: labeled "Text", asked on the form, and focused with the keyboard up when it is the first or only field.
\end{itemize}

\textbf{Two worked examples:}

\begin{verbatim}
{
  "key": "t",
  "description": "Tracking",
  "type": "table-line",
  "olp": ["🛍️ Bag - lists and notes", "❤️ Health", "🏩 Tracking", "%<%Y-%m>"],
  "preText": "#+PLOT: title:\"Tracking\" ind:4 deps:1 with:points ticks:6",
  "postText": "#+TBLFM: $4=round(trunc($3/100) + 100*($3/100 - trunc($3/100)) / 60)",
  "template": "| @# | %<%Y-%m-%d> | %<%H%M> ||",
  "prependHeading": true,
  "file": "github:home.org"
}
\end{verbatim}

\textbf{Tracking} logs a value with its date and time into a table, and plots it:

\begin{itemize}
\item \texttt{olp} is four segments deep. The last, \texttt{"\%<\%Y-\%m>"}, creates the current month's heading the first time anything is captured that month.
\item \texttt{preText} sets up the chart once, when the month's table is created: \texttt{ind:4} and \texttt{deps:1} plot the 4th column against the 1st (the row number) as points.
\item \texttt{postText} sets up the formula once. It converts a raw \texttt{HHMM} value (\texttt{\$3}, such as \texttt{1830}) to the nearest whole hour in \texttt{\$4}: \texttt{trunc(\$3/100)} is the hour, the rest turns the minutes into a fraction, and \texttt{round()} rounds the sum, so 1830 gives 19, not 18.5.
\item \texttt{template} is the one row: \texttt{@\#}, a date, an \texttt{HHMM} time, and a trailing \texttt{||} that leaves the 4th cell for the formula.
\item \texttt{prependHeading: true} keeps the newest month's heading on top, while its rows append in order.
\item \texttt{file: "github:home.org"} captures into that file on GitHub, whatever is open.
\end{itemize}

\begin{verbatim}
{
  "key": "v",
  "description": "vCard Contact",
  "type": "plain",
  "olp": ["New"],
  "file": "github:contacts.org",
  "omitEmptyEntries": true,
  "template": "* %^{First Name} %^{Last Name} :%^{Category/Tag}:\nBEGIN:VCARD\nVERSION:3.0\nFN:%\\1 %\\2\nN:%\\2;%\\1;;;\nEMAIL:%^{Email}\nTEL;TYPE=voice:%^{Phone Number}\nTEL;TYPE=cell:%^{Mobile Number}\nADR;TYPE=home:;;%^{Street};%^{City};%^{State};%^{Postal Code};%^{Country}\nBDAY:%^{Birthday (YYYY-MM-DD)}\nNOTE:%^{Note}\nORG:%^{Company/Organization}\nTITLE:%^{Job Title}\nEND:VCARD\n\n%?"
}
\end{verbatim}

\textbf{vCard Contact} builds a whole heading with a \texttt{plain} capture: a title with an inline tag prompt, and a raw multi-line body. \texttt{FN:\%\textbackslash{}1 \%\textbackslash{}2} and \texttt{N:\%\textbackslash{}2;\%\textbackslash{}1;;;} reuse the two name answers through backreferences instead of asking twice. The body has the shape the vCard export produces for its \textbf{new buffer} destination (see \hyperref[sec:org053e4f4]{Contacts (.vcf) export}), which Import reads back (see \hyperref[sec:orgc97ce31]{Import}). \texttt{omitEmptyEntries: true} drops every blank field, and the trailing \texttt{\%?} leaves the cursor below the vCard for a free-form note.

\subsubsection{Capturing from other apps}
\label{sec:orge0ecaf9}
On Android, the installed app is a share target: choose \textbf{org-pwa} in another app's Share sheet and the shared title, text and link open Capture. Pick a template, and its \texttt{\%i} and \texttt{\%a} take the content. Long-press the app icon for a \textbf{Capture} shortcut.

The built-in \textbf{Web page or text} template is made for this. Share a web page and it files a heading under \textbf{Inbox} that links to the page, titled as the page is; leave its title prompt blank. Share some text from a document and it files the text under a heading with the title you type. Either way the second prompt takes your own note, and \texttt{:CREATED:} records when. A part that isn't there (no text, no note) leaves a blank line. If you already have templates of your own, add it in \textbf{Settings → Capture Templates}:

\begin{verbatim}
{ "key": "w", "description": "Web page or text", "type": "plain", "olp": ["Inbox"],
  "template": "* %a%^{Title (leave blank for a web page)}\n:PROPERTIES:\n:CREATED: %U\n:END:\n  %i\n  %?", "emptyLines": 1 }
\end{verbatim}

Any link or bookmark can do the same. \texttt{index.html?capture=KEY} runs the template with that key, a bare \texttt{?capture} shows the list, and \texttt{\&text=}, \texttt{\&title=} and \texttt{\&url=} carry content. The address is cleaned afterwards, so reloading never captures twice, and it works offline.

Example: \texttt{index.html?capture=n\&text=Buy\%20milk} with this template files the text under \textbf{Inbox}:

\begin{verbatim}
{ "key": "n", "description": "Shared note", "type": "plain", "olp": ["Inbox"], "template": "* Note\n  %i\n  %a" }
\end{verbatim}

A template that files into another file (\texttt{file:}) needs a document from the same backend open; the app restores your tabs when it starts.

\textbf{iPhone}: iOS lists neither installed web apps in the Share sheet nor opens them from a link, so none of this reaches it. Copy the text instead (a Shortcut can do it), open the app and run a template that contains \texttt{\%x}.

\noindent\rule{\textwidth}{0.5pt}

\subsection{Extras menu}
\label{sec:org08ac3df}
A floating ☰ button at bottom-right opens a configurable popup of quick actions, a vertical list next to the button and extending to its left. It is this app's own feature, set with \texttt{org-xx-extra-menu} (Global or Local Variables; see \hyperref[sec:org75c5aa1]{Local Variables}). The button is hidden whenever the variable is unset or resolves to nothing selectable. The sibling variable \texttt{org-xx-menu-aliases} relabels or hides buttons in the existing File, More, Export and View menus; see \hyperref[sec:orgfca7ef3]{Menu customization}.

\textbf{Syntax}: a sequence of double-quoted \texttt{"SPEC;LABEL"} entries, spread over several lines with a trailing \texttt{\textbackslash{}} (see \hyperref[sec:org4f4f89f]{Configuration} for line continuation). Settings → Quick Settings shows and edits it in this one-entry-per-line form (see \hyperref[sec:orgbb54b74]{Settings}):

\begin{verbatim}
org-xx-extra-menu: "j;📔 Journal" \
                   "c;✅ Checklist" \
                   "m;👥 Meeting" \
                   "-----" \
                   "t;⭐ Tracking" \
                   "["🛍️ Bag - lists and notes", "❤️ Health", "🏩 Tracking", "%<%Y-%m>"];👁 Current" \
                   "-----" \
                   "'calendar;📅 Calendar"
\end{verbatim}

\textbf{LABEL} is the text shown, exactly as written (icon and all). \textbf{SPEC} is one of three shapes:

\begin{itemize}
\item A bare \textbf{capture-template key} (\texttt{t}, \texttt{j}, \texttt{c}, \texttt{m} above) runs that template, as if picked from the Capture menu (see \hyperref[sec:org83b5238]{Capture Templates}). A key with no template shows a status message.
\item A bracketed \textbf{OLP path}: an array of double-quoted heading titles, as in a capture template's \texttt{olp:}, each segment expanded for \texttt{\%<FORMAT>} when the item is tapped (\texttt{"\%<\%Y-\%m>"} gives the current year and month). It navigates to that heading. Unlike a capture template, it never creates a missing heading, because browsing a menu shouldn't write to the file: a path that doesn't fully exist creates nothing and says so.
\item A \textbf{quoted function reference} (\texttt{'calendar} above) runs the command palette command with that real Emacs or Org function name, as if chosen from the palette. It acts on the heading whose action menu is open, else the keyboard-focused one. A command that can't run says why in the status line (\texttt{Clock out: no clock is running.}), and no confirmation is asked, since choosing the item is the confirmation. The names are the ones the palette shows and searches (\texttt{org-clock-out}, \texttt{org-clock-in}, \texttt{org-refile}, \texttt{org-todo}, \texttt{org-agenda}, \texttt{save-buffer}, \texttt{isearch-forward}, \texttt{calendar}, \texttt{org-html-export-to-html} and the rest), so type a function name into the palette to find the command, and quote the same name here.

\textbf{Exceptions}: \textbf{New document} and \textbf{Undo history} have no Emacs function to name. Emacs has no dedicated new-buffer command (\texttt{switch-to-buffer} with an unused name is the usual way, and \texttt{scratch-buffer} only returns to the one \texttt{*scratch*} buffer), and its undo history is a variable, \texttt{buffer-undo-list}, not a command. A name that no command has shows \texttt{'name is not a recognized function.} when tapped; the entry stays in the menu, so a typo is visible.

\textbf{Recalculate all tables} (\texttt{org-table-recalculate-buffer-tables}) recalculates every \texttt{\#+TBLFM:} formula in the document as one undo step, where org's \texttt{C-c C-c} recalculates one table (see \hyperref[sec:org1458240]{Table formulas}). \textbf{Export as an Org buffer} is \texttt{org-org-export-as-org} (see \hyperref[sec:org053e4f4]{Export}).

\textbf{The calendar} (\texttt{'calendar}, the palette's \textbf{Calendar}) is named after Emacs's command but is this app's own single-month overview, with year and month navigation and a Today jump. Today is bold with a two-colour halo ring that stays visible on any background. Each day is colour-coded by what is on it, from the same agenda pipeline the Agenda uses: green for a contact's birthday (\texttt{org-contacts-anniversaries}, see \hyperref[sec:org105cb3a]{Agenda}), blue for any other event (SCHEDULED, DEADLINE, a plain timestamp or a diary-sexp entry), and burnt orange when both. Ambient daily diary entries (sunrise, sunset, weather) and logbook entries don't count, since they would colour nearly every day. Tapping a date closes the calendar and opens Agenda's Day view on that date, whatever is open or connected. The displayed month is remembered for the rest of the session.
\end{itemize}

Five hyphens (\texttt{"-{}-{}-{}-{}-"}) on their own are a visual separator rather than an item, for grouping related actions.

A malformed entry (bad JSON in an OLP array, a function symbol that isn't a valid name, a spec matching none of the three shapes, a missing label) is skipped without breaking the rest of the menu.

\noindent\rule{\textwidth}{0.5pt}

\subsection{Menu customization}
\label{sec:orgfca7ef3}
\texttt{org-xx-menu-aliases} (a Global or Local Variable; this app's own, not an org variable) relabels or hides buttons across the File, More, Export, Import and View menus at once. Its sibling \texttt{org-xx-extra-menu} adds a new quick-access menu instead (see \hyperref[sec:org08ac3df]{Extras menu}).

\textbf{Syntax}: double-quoted \texttt{"menu:Label;alias"} entries, space-separated and spread over several lines with a trailing \texttt{\textbackslash{}}. \texttt{menu} is \texttt{file}, \texttt{more}, \texttt{export}, \texttt{import} or \texttt{view}. Settings → Quick Settings shows and edits it one entry per line (see \hyperref[sec:orgbb54b74]{Settings}):

\begin{verbatim}
org-xx-menu-aliases: "more:Commands;💲 Commands" \
                     "more:Export;📤 Export" \
                     "more:Import;📥 Import" \
                     "more:Settings;⚙️ Settings"
\end{verbatim}

Export and Import are steps within the More menu and have their own namespaces, \texttt{export:} and \texttt{import:}. A form like \texttt{"more:export:HTML;📄"} isn't valid: the colon splits only once, so it is read as the label \texttt{export:HTML} in the \texttt{more} namespace, which matches nothing and has no effect, without an error.

\texttt{Label} is one of the menu's built-in button labels (the table below). A label that matches no button, or a \texttt{menu} other than the five, is skipped without an error. A label never mentioned keeps its default text: this is an override list, not a full redefinition.

\textbf{\texttt{alias}}, after the semicolon:

\begin{itemize}
\item \textbf{Non-empty}: replaces the button's visible text. The button does the same thing, and the original label becomes its accessible name (for screen readers and tooltips).
\item \textbf{Empty} (\texttt{"menu:Label;"}): the button is omitted entirely. \texttt{"export:Markdown;"} hides the Markdown button from the Export menu. \textbf{\texttt{more:Settings} is the exception}: it can never be hidden, because GitHub and WebDAV setup and Global Variables are reachable only from there, so an empty alias keeps the label \texttt{Settings}.
\end{itemize}

\textbf{Order}: buttons can be reordered, but only once every one of a menu's labels is listed (the table has the complete sets). A partial list leaves that menu's order untouched, so setting one alias never reorders the rest. Once all are listed, the buttons appear in the order the entries are written. An empty alias counts as listed but also hides the button (except \texttt{more:Settings}), so to keep a button visible with its default text, repeat that text as the alias: \texttt{"more:Settings;Settings"}. For example, to show the More menu as Settings, Commands, Import, then Export, list all four:

\begin{verbatim}
org-xx-menu-aliases: "more:Settings;Settings" "more:Commands;Commands" \
                     "more:Import;Import" "more:Export;Export"
\end{verbatim}

\textbf{Popup menus}: File, View, More and Extras open as small bordered popups next to their trigger button, below it if there is room and above it otherwise (the Extras button is near the bottom, so its popup opens upward). Each is sized to its widest row, with one full-width row per button, and on a mouse a row highlights as the pointer passes over it. Tapping outside closes a menu, and selecting an action closes it at once, with the result reported on the status line. Every popup (File, View, More, Export, the Open and Save As backend picker, and Extras) shares one size setting, Settings → Appearance → \textbf{Menu Size} (Regular or Small). Export is a step within the More popup, and a \textbf{‹ Back} button returns to the More list without closing anything.

The real labels of each menu, as used in \texttt{"menu:Label;alias"}:

\begin{center}
\begin{tabular}{ll}
\texttt{menu} & Labels\\[0pt]
\hline
\texttt{file} & \texttt{Export}, \texttt{New}, \texttt{Open}, \texttt{Save}, \texttt{Save As}, \texttt{Save Some Buffers}\\[0pt]
\texttt{more} & \texttt{Commands}, \texttt{Export}, \texttt{Import}, \texttt{Settings}\\[0pt]
\texttt{export} & \texttt{ASCII}, \texttt{As-org}, \texttt{Calendar (.ics)}, \texttt{Contacts (.vcf)}, \texttt{HTML}, \texttt{Markdown}, \texttt{ODT}\\[0pt]
\texttt{import} & \texttt{Contacts (.vcf)}, \texttt{iCalendar (.ics)}\\[0pt]
\texttt{view} & \texttt{Agenda}, \texttt{Org}, \texttt{Text}, \texttt{TODO}\\[0pt]
\end{tabular}
\end{center}

\texttt{Save Some Buffers} can be given an alias or hidden like the others, but it always comes last: it isn't part of a reordering list, so a File menu order you set earlier keeps working.

\noindent\rule{\textwidth}{0.5pt}

\subsection{File management}
\label{sec:org5584f85}
\textbf{File menu}: New, Open, Save, Save As, Save Some Buffers. A 💾 button beside the File menu icon does what \textbf{File → Save} does, in one tap. It is disabled under the same conditions as Save: no document open, or a new unsaved document with nowhere to save to yet.

\begin{itemize}
\item \textbf{New} creates a document in memory, with no file picker, name or backend to choose. You can edit it at once, and its edits are cached locally like any open document's, so a crash or an accidentally closed tab can be resumed with the usual "resume unsaved local version" prompt. Nothing is written to a backend until you choose to keep it. The tab bar and modeline show \texttt{*scratch*}, Emacs's name for a fresh unnamed buffer, in the same red-with-a-dot style as unsaved edits, until the first heading gets a title, which then becomes the label. \textbf{Save} stays disabled, and \textbf{Save As} picks a destination and makes it a persisted file. It opens in its own tab (see \hyperref[sec:orgda9f5ef]{Tabs}), and is handy as a scratch pad for trying a capture template, testing a table formula, or drafting something that may be thrown away.
\item \textbf{Open} and \textbf{Save As} go to one screen with a \textbf{GitHub / Local file / WebDAV} picker (\textbf{Import file} / \textbf{Download} on iOS, where a file is read once rather than linked). It opens on the type you used last, and switching types never closes it. \textbf{Local file} is one button into your device's file picker.
\item \textbf{GitHub and WebDAV show a real folder listing.} Tap a folder (📁) to enter it. Files shown are \texttt{.org} files (📄) and \texttt{\_archive} files (🗄️, org's archive-file naming; see \hyperref[sec:orgf0de015]{Archiving}), and other files are hidden. An \textbf{↑ Up} button appears below the root. GitHub uses its Contents API, and WebDAV uses \texttt{PROPFIND}, which Nextcloud, ownCloud and Apache \texttt{mod\_dav} support. The \textbf{Enter path…} field above the list opens a file by its path (press Enter), which also works when listing fails (a server without good \texttt{PROPFIND}, a token without list permission on a folder). A backend you haven't set up says so and links to its settings.
\item \textbf{Save} writes back to the backend the file came from, without asking again.
\item \textbf{Save Some Buffers} (\texttt{save-some-buffers}, \texttt{C-x s}) goes through every modified tab that has a file and asks about each in the minibuffer, one key: \texttt{y} or Space saves it, \texttt{n} or Backspace skips it, \texttt{!} saves it and the rest without asking, \texttt{.} saves it and stops, \texttt{q} (or Escape) stops, \texttt{d} shows what changed since it was last saved. Each save is the same \textbf{Save} as the File menu's (so an outside change is settled the usual way), and you return to the tab you started on. \textbf{File → Save Some Buffers} does not ask: it saves them all. It is dimmed when nothing is modified, and skips a new unsaved document, which needs \textbf{Save As}.
\item \textbf{Revert Buffer} (palette, \texttt{revert-buffer}) reads the open file again and drops the buffer's unsaved changes, after asking if there are any. \textbf{Undo} brings them back.
\item \textbf{Capture, Refile, Archive and Replace edit buffers, not files.} An entry sent to another file goes into that file's tab, which opens in the background if it has none, and shows the modified dot until you save it. Nothing is written behind your back, a file's tab keeps whatever you had typed in it, and the Agenda reads the buffer, so a new entry shows before it is saved. Like any edit, a modified buffer is also kept on the device, so closing the app doesn't lose it.
\item \textbf{Line endings} are kept. A Windows-style (CRLF) file is edited as ordinary text and saved back with CRLF, and a Unix-style one stays Unix-style, so a save changes only the lines you edited.
\item \textbf{Save As} picks a new destination, possibly on a different backend, or the first destination of a new unsaved document. Browse to a folder, type a name and press Save; if that file already exists you are asked before it is replaced. If someone else saves that file after you agreed, you are asked again rather than their version being overwritten.
\item If a file has unsaved local edits when you open it, you are asked whether to \textbf{resume} them or \textbf{discard} them and load the current version; either way the file opens. The prompt shows when the edit was made, and discarding clears it. \textbf{Settings → Pending Local Changes} lists every file with unsynced edits, with a Discard button for each and \textbf{Discard all} when there is more than one.
\item \textbf{Conflicts}: if the stored version changed since the last sync, Save merges the two where it safely can and asks only about the parts that clash (see \hyperref[sec:org57d263c]{Offline behavior and sync}).
\item \textbf{Recently opened}, below a double separator, lists files you have opened or saved as, never the file that is currently open. It persists across sessions, is sorted alphabetically and capped at five. Long-press an entry to remove just that one, or long-press the label to clear the whole list. Tapping an entry opens it directly, with no path or picker prompt, except for a file brought in with \textbf{Import} (the read-once fallback for browsers without File System Access, like iOS Safari), which can't be reopened without selecting it again.
\end{itemize}

\noindent\rule{\textwidth}{0.5pt}

\subsection{Tabs}
\label{sec:orgda9f5ef}
\textbf{Several documents can be open at once.} Opening a file (New, Open, a Recently opened entry, a \texttt{file:}, \texttt{github:} or \texttt{webdav:} link, or a result from Agenda or Search) adds it as a tab rather than replacing what is open, and all of these use the same open logic. Opening a file that is already open in another tab switches to that tab. Each tab keeps its own content, cursor and keyboard focus, mid-edit fields, scroll position, view mode and undo/redo history (which persists across a reopen; see \hyperref[sec:org9a920ce]{Undo/redo and History}).

The tab bar is always visible, even with one document. Each tab shows a short label (the filename, or the first heading's title for an unsaved document; several unsaved drafts can be open at once, each with its own local cache) and a solid dot when it has unsaved changes, the same red "modified" convention used elsewhere. Tap a tab to switch to it. The small × closes it, after asking first if it has unsaved changes (\texttt{"X" has unsaved changes. Close it anyway?}), as Emacs's \texttt{kill-buffer} does; a clean tab closes at once. Closing the last tab opens a fresh empty unsaved document in its place, so there is always a tab to work in.

The full set of open tabs persists across closing and reopening the app, each restored as it was left, with the previously active one active again.

\textbf{Help} has its own tab (see \hyperref[sec:orgb4d67ea]{Help}), always last and styled distinctly (green background, gold underline when active). Switching to a document tab leaves it open but inactive, and unlike document tabs it doesn't survive closing and reopening the app.

\noindent\rule{\textwidth}{0.5pt}

\subsection{Citations}
\label{sec:org261744f}
\textbf{Org-cite citations} work as in Emacs's default (\texttt{basic}) processor. Write \texttt{[cite:@knuth84]}, \texttt{[cite/t:@knuth84 p. 5]} or \texttt{[cite:see @a; @b, ch. 2]} and point at your references with \texttt{\#+BIBLIOGRAPHY: refs.bib} (BibTeX or CSL-JSON, read from the same place as the document, one line per file). In the outline a citation stays as typed, dimmed, and tapping a key shows its entry; a key with no entry is dotted underlined.

\begin{itemize}
\item Styles: default (author, year), \texttt{/t} (text), \texttt{/a} (author), \texttt{/na} (no author), \texttt{/n} (nothing printed), \texttt{/b} (numeric), with \texttt{/c} and \texttt{/b} variants for capitals and bare form.
\item \texttt{\#+PRINT\_BIBLIOGRAPHY:} marks where the reference list goes in HTML, Markdown, ASCII and ODT exports. \texttt{\#+CITE\_EXPORT: basic plain} (or \texttt{numeric}) picks the list style.
\item Settings (see \hyperref[sec:org75c5aa1]{Local Variables}): \texttt{org-cite-global-bibliography} (files used by every document), \texttt{org-cite-basic-sorting-field} and \texttt{org-cite-basic-author-year-separator}.
\end{itemize}

\noindent\rule{\textwidth}{0.5pt}

\subsection{Export}
\label{sec:org053e4f4}
\textbf{More → Export} gets a document out as ASCII, Markdown, HTML, ODT, As-org, a calendar file (.ics) or contacts (.vcf). File keywords such as \texttt{\#+TITLE} and \texttt{\#+OPTIONS:} match case-insensitively, as in org.

\textbf{Scope}: \textbf{Choose a heading…} exports just that heading and its descendants, and every heading is listed, even folded ones. The chosen heading becomes the top level of the output. \textbf{This file} exports everything. Calendar and Contacts also offer \textbf{This file + Agenda Files} or \textbf{This file + Contacts Files} once those are configured, fetched fresh at export time. In Tree style, Contacts lists only contacts (headings with \texttt{:KIND:}).

\subsubsection{Document settings}
\label{sec:orgbf809e0}

\texttt{\#+TITLE}, \texttt{\#+AUTHOR} and \texttt{\#+DATE} appear in a whole-document export in ASCII, HTML and ODT: a visible title with the author and date beneath it (in HTML, the date and author close the page, as in Emacs), and the format's own metadata where it has some. Markdown has none of them, as in org's \texttt{ox-md}. A subtree export starts at its heading instead.

\texttt{\#+OPTIONS:} takes these keys, chained on one line or split across several (a later value wins):

\begin{center}
\begin{tabular}{ll}
Key & Effect\\[0pt]
\hline
\texttt{toc:} & An integer depth limit, or \texttt{nil} for no table of contents. Independent of \texttt{num:}\\[0pt]
\texttt{num:} & Heading numbers on or off. In Markdown it switches the table of contents between an ordered and an unordered list\\[0pt]
\texttt{\textasciicircum{}:} & Sub/superscript strictness, as \texttt{org-use-sub-superscripts} (see \hyperref[sec:org75c5aa1]{Local Variables}); overrides it for this export\\[0pt]
\texttt{author:}, \texttt{date:} & Show each field or not\\[0pt]
\texttt{creator:} & A "Generated by org-pwa" notice, off by default\\[0pt]
\end{tabular}
\end{center}

HTML export also takes these, with Org's defaults:

\begin{center}
\begin{tabular}{ll}
Key & Effect\\[0pt]
\hline
\texttt{H:} & How many heading levels are headings (3); deeper ones become list items\\[0pt]
\texttt{num:N} & Number only the first N levels\\[0pt]
\texttt{tags:}, \texttt{todo:}, \texttt{pri:} & Show tags (\texttt{not-in-toc} leaves them out of the contents), TODO keywords, priorities (off)\\[0pt]
\texttt{p:} & Show SCHEDULED, DEADLINE and CLOSED (off)\\[0pt]
\texttt{\textbackslash{}n:t} & Every line break in a paragraph is kept\\[0pt]
\texttt{-:}, \texttt{*:}, \texttt{e:}, \texttt{f:}, \texttt{stat:} & Special strings (\texttt{-{}-{}}, \texttt{...}), emphasis, entities, footnotes and \texttt{[1/3]} cookies\\[0pt]
\texttt{d:} & Drawers: \texttt{t}, \texttt{nil}, or a list such as \texttt{("NOTES")} (all but LOGBOOK)\\[0pt]
\texttt{arch:}, \texttt{tasks:} & Archived subtrees (title only), TODO headings (\texttt{nil}, \texttt{todo}, \texttt{done})\\[0pt]
\texttt{<:}, \texttt{title:}, \texttt{email:} & Timestamps, the title, the author's address\\[0pt]
\end{tabular}
\end{center}

\texttt{\#+EXCLUDE\_TAGS:} (default \texttt{noexport}), \texttt{\#+SELECT\_TAGS:}, COMMENT headings, \texttt{\#+SUBTITLE:}, \texttt{\#+DESCRIPTION:}, \texttt{\#+KEYWORDS:}, \texttt{\#+LANGUAGE:} and \texttt{\#+MACRO:} are honoured. Other formats ignore these HTML-only keys; Markdown supports only \texttt{toc:} and \texttt{num:}.

\texttt{\#+INCLUDE: "path" [block-type] [language] [:lines "N-M"] [:minlevel N]} splices another document in at export time, before the first heading or in a heading's body. In a body it nests as children of that heading, one level deeper by default; at document level it keeps its own levels. The path is read from an explicit \texttt{github:}, \texttt{webdav:} or \texttt{local:} prefix, or else relative to the open file.

\begin{itemize}
\item With no block type, the file is parsed as org and its headings are spliced in (its own \texttt{\#+TITLE} and similar keywords are not). \texttt{:minlevel N} shifts its top level to N.
\item With a block type (\texttt{src}, \texttt{example}, \texttt{quote} or \texttt{export}, optionally with a language or backend, as in \texttt{\#+INCLUDE: "path" export html}), the raw text is wrapped in that block and placed in the body.
\item \texttt{:lines "N-M"}, \texttt{"-M"} or \texttt{"N-"} filters the text first. The end line is exclusive, as in org, despite the manual's wording.
\end{itemize}

A file that can't be found is skipped silently.

\texttt{\#+HTML\_HEAD:} inserts raw HTML before \texttt{</head>}, in HTML export only and for subtree exports too; several lines are joined in order. \texttt{\#+BEGIN\_EXPORT <backend>} \ldots{} \texttt{\#+END\_EXPORT} passes its content through unescaped when the backend matches the current format (\texttt{html}, \texttt{ascii}, \texttt{md} or \texttt{markdown}, \texttt{odt}, where the content must already be ODF XML) and is omitted otherwise.

\textbf{Header and footer}: \texttt{html-preamble:} and \texttt{html-postamble:} in \texttt{\#+OPTIONS:} take a quoted text, as in Emacs, placed above and below the page (\texttt{nil} drops the postamble). In it \texttt{\%t} is the title, \texttt{\%s} the subtitle, \texttt{\%a} the author, \texttt{\%d} the date, \texttt{\%e} the email, \texttt{\%C} or \texttt{\%T} when it was exported, and \texttt{\%\%} a percent sign; \texttt{\textbackslash{}n} starts a new line and HTML is allowed.

\begin{verbatim}
#+OPTIONS: html-postamble:"Prepared by %a, %d"
\end{verbatim}

For text on every printed page, such as a name and a page number, use \texttt{page-header:} and \texttt{page-footer:}. Parts are separated by \texttt{|}: three are left, centre and right, two are left and right, one is the centre. \texttt{\%p} is the page number, \texttt{\%P} the number of pages, and the codes above work too.

\begin{verbatim}
#+OPTIONS: page-footer:"%a | | Page %p of %P"
\end{verbatim}

They take effect when you print or use \textbf{Print → Save as PDF} in Chrome (also on a ChromeOS laptop); Firefox and Safari ignore them. Pages get a 2 cm margin to hold them. For your own CSS, such as \texttt{@page \{ @bottom-center \{ content: counter(page); \} \}}, use \texttt{\#+HTML\_HEAD:}.

\subsubsection{ASCII}
\label{sec:orgacd5367}

Plain text, org's \texttt{ox-ascii} for the core cases. Level-1 headings are underlined with \texttt{=} and deeper ones indented. Paragraphs wrap at \texttt{org-ascii-text-width} (default 72; see Local Variables), never mid-word. A link becomes \texttt{description (target)}. Bold, italic, underline and strikethrough are stripped, and \texttt{code} and \texttt{verbatim} become \texttt{`text'}. Footnote references become \texttt{[1]}, with their definitions under "Footnotes".

\textbf{Characters}: entities such as \texttt{\textbackslash{}alpha}, \texttt{\textbackslash{}rarr}, \texttt{\textbackslash{}copy} and \texttt{\textbackslash{}nbsp} are replaced by the chosen character set, as in ox-ascii. Set \texttt{org-ascii-charset} in Local Variables to \texttt{ascii} (the default), \texttt{latin1} or \texttt{utf-8}. UTF-8 also turns \texttt{-{}-{}}, \texttt{-{}-{}-} and \texttt{...} into dashes and an ellipsis, and draws underlines, bullets, checkboxes (☑ ☐ ☒), horizontal rules and block frames with box and symbol characters. Code and verbatim text are left alone.

\textbf{Blocks}: example and source blocks and lines starting with a colon sit in a frame (drawn with \texttt{,-{}-{}-{}-{}}, \texttt{|} and \texttt{`-{}-{}-{}-{}}, or with box-drawing characters in UTF-8). QUOTE and VERSE blocks are set in from the margin (a quote's text is also narrower); a verse keeps its line breaks. CENTER blocks centre each line. Other special blocks, including justify-left and justify-right, are set as ordinary text, and COMMENT blocks are omitted. \texttt{\#+ATTR\_ASCII:} lines are read and ignored: \texttt{:width} changes nothing in Emacs 29.3's output either. Headings are numbered (1, 1.1, 2) and, with more than one heading, a "Table of Contents" comes first; \texttt{toc:} and \texttt{num:} control both.

Tables are typeset: each column is padded to its widest cell, and mostly-numeric columns are right-aligned. A row of \texttt{<N>} cells, like \texttt{| <30> | <10> | <24> |} (usually above the header), fixes each column to N characters; wider content wraps within the row, breaking at hyphens when needed. That cookie row never appears in any format's output, and only ASCII uses the widths.

Not implemented: Latin-1 and UTF-8 substitutions, \texttt{:width} rules and justification blocks (see \hyperref[sec:org2d2d41c]{Differences from Emacs org-mode}).

\subsubsection{Markdown}
\label{sec:org1f56ceb}

GFM. Checkboxes become task-list items (\texttt{- [ ]}, \texttt{- [x]}), tables become pipe tables, code fences keep their language, and footnotes become \texttt{[\textasciicircum{}label]}. TODO keyword, priority and tags stay as text on the heading (\texttt{\# TODO [\#A] Buy milk}). Underline and sub/superscript use raw \texttt{<u>}, \texttt{<sub>} and \texttt{<sup>}.

Internal links (\texttt{[[\#custom-id]]}, \texttt{[[*Heading]]}, fuzzy titles) become GitHub-style \texttt{\#slug} anchors: lowercase, punctuation stripped, spaces to hyphens, duplicates suffixed \texttt{-1}, \texttt{-2}. A link outside the exported scope, or one that doesn't resolve, stays as its literal text.

\subsubsection{HTML}
\label{sec:org2247fc2}

A complete standalone page laid out like Emacs's \texttt{ox-html}: the title as \texttt{h1}, top-level headings as \texttt{h2} in numbered outline containers, \texttt{org-ul}, \texttt{org-ol} and \texttt{org-dl} lists whose items keep their paragraphs, sub-lists and blocks, tables with a header row (when a rule follows it), column alignment and caption, figures, and footnotes numbered in order of reference. Entities (\texttt{\textbackslash{}alpha}, \texttt{\textbackslash{}rarr}), timestamps, inline source, macros, \texttt{@@html:...@@} snippets, special blocks (\texttt{\#+BEGIN\_NOTE} becomes \texttt{<div class="note">}), verse and center blocks, \texttt{\#+TOC: headlines}, \texttt{tables} and \texttt{listings}, and \texttt{\#+ATTR\_HTML:} on tables and images all work. Source blocks are coloured for common languages. Pictures the app can read (GitHub, WebDAV, attachments) are embedded, so the file stands alone.

The CSS is print-aware (headings, tables, figures and code stay whole across pages, and long code lines wrap), so the browser's \textbf{Print → Save as PDF} is the route to PDF.

\subsubsection{ODT}
\label{sec:orga95287b}

A standalone \texttt{.odt} that opens in LibreOffice, Word (with the ODF filter) and Google Docs. Headings are outline levels; bold, italic, underline, strikethrough and code use named character styles; checkboxes show as ☐ and ☑; footnotes are real, editable ODF footnotes; and internal links are bookmark references. Images are referenced by path, not embedded.

\subsubsection{As-org}
\label{sec:org38a8899}

Org's \texttt{org-org-export-as-org}: a fully processed copy of the document in a new unsaved buffer named \texttt{*Org ORG Export*}, leaving the original untouched. Whole document only. It:

\begin{itemize}
\item resolves \texttt{\#+INCLUDE:} lines;
\item expands \texttt{\{\{\{macro\}\}\}} references: \texttt{\#+MACRO: name text} with \texttt{\$1}, \texttt{\$2} arguments, written \texttt{\{\{\{name(arg1,arg2)\}\}\}} with \texttt{\textbackslash{},} for a literal comma, plus \texttt{\{\{\{title\}\}\}}, \texttt{\{\{\{author\}\}\}}, \texttt{\{\{\{email\}\}\}}, \texttt{\{\{\{date\}\}\}} and \texttt{\{\{\{date(FORMAT)\}\}\}} (FORMAT applies only when \texttt{\#+DATE:} is a single well-formed timestamp), \texttt{\{\{\{time(FORMAT)\}\}\}} and \texttt{\{\{\{keyword(NAME)\}\}\}};
\item drops \texttt{:noexport:} headings with their subtrees, and strips \texttt{\#+BEGIN\_COMMENT} blocks and comment lines (those starting with \texttt{\#} and a space);
\item shifts the top level to 1.
\end{itemize}

Not implemented: running source blocks, and macros whose text starts with \texttt{(eval}. Both are left as written (see Differences from Emacs org-mode). In an Extras entry it is \texttt{'org-org-export-as-org}.

\subsubsection{Calendar (.ics)}
\label{sec:orgf0428d8}

Every dated item as an iCalendar file for any calendar app: SCHEDULED, DEADLINE, plain active title timestamps and \hyperref[sec:org105cb3a]{org-contacts-anniversaries}. It follows the Agenda's exclusions (archived, commented and completed headings). Re-exporting gives the same event IDs, so a calendar app updates events instead of duplicating them.

\textbf{To}: the file is downloaded. \textbf{To: file} opens the share sheet instead, which on an iPhone has \textbf{Save to Files} (a downloaded .ics there goes straight to Calendar); \textbf{To: new buffer} leaves the text in a new buffer.

\textbf{Scope}: \textbf{This file}, \textbf{Choose a heading…} and \textbf{This file + Agenda Files} write every dated item, with no window. \textbf{Agenda (as displayed)}, the last row, writes what \hyperref[sec:org66ee6b6]{the calendar sync} sends: what View → Agenda shows, from a month back to six months ahead.

\begin{center}
\begin{tabular}{lll}
Org source & vCalendar field & Notes\\[0pt]
\hline
SCHEDULED & \texttt{DTSTART}, \texttt{SUMMARY} (\texttt{DESCRIPTION} "Scheduled") & \\[0pt]
DEADLINE & \texttt{DTSTART}, \texttt{SUMMARY} (\texttt{DESCRIPTION} "Deadline") & A warning suffix (\texttt{-3d}) also becomes a \texttt{VALARM} reminder\\[0pt]
Plain active title timestamp & \texttt{DTSTART}, \texttt{SUMMARY} & Only when the heading has no SCHEDULED or DEADLINE\\[0pt]
A repeater (\texttt{+1w}, \texttt{++3d}, \texttt{.+1m}) on any of these & \texttt{RRULE} & The calendar app expands the recurrence\\[0pt]
A time range (\texttt{09:00-10:15}) & \texttt{DTEND} & \\[0pt]
A heading's tags & \texttt{CATEGORIES} & \\[0pt]
\texttt{:LOCATION:} & \texttt{LOCATION} & \\[0pt]
The start of the body & \texttt{DESCRIPTION} & At most 400 characters\\[0pt]
org-contacts-anniversaries & \texttt{DTSTART} with a yearly \texttt{RRULE}, \texttt{SUMMARY} (\texttt{Birthday: Name}) & One yearly event cannot carry an age, so the date is in the description\\[0pt]
\end{tabular}
\end{center}

\subsubsection{Contacts (.vcf)}
\label{sec:org72b2ec3}

vCard 3.0 for any contacts app, modeled on \href{https://github.com/pinoaffe/org-vcard}{org-vcard} (compatible with org-contacts.el). \href{https://www.rfc-editor.org/rfc/rfc6350}{RFC 6350} requires only \texttt{FN}, so no email address is needed. There are two styles, and \textbf{Tree} is the default:

\begin{itemize}
\item \textbf{Flat}: every heading with at least one recognized contact property (\texttt{:EMAIL:}, \texttt{:PHONE:}, \texttt{:CELL:}, \texttt{:WORK\_PHONE:}, \texttt{:ADDRESS:}, \texttt{:ADR:}, \texttt{:NICKNAME:}, \texttt{:NOTE:}, \texttt{:ANNIVERSARY:} or the birthday property). \texttt{:IGNORE: t} excludes one. A heading with none is skipped.
\item \textbf{Tree}: a contact heading (\texttt{:KIND: individual}, \texttt{:FIELDTYPE: name}) with each field as a descendant heading carrying its own \texttt{:FIELDTYPE:}, at any depth. A contact with no fields still exports a minimal vCard (\texttt{FN}, \texttt{N}).
\end{itemize}

The \textbf{To: new buffer} checkbox sends the vCard text to a new unsaved document (a \texttt{vCard(s)} heading, titled \texttt{*scratch*}) instead of downloading a file; Import reads that shape back.

There are no name, tag or property filters. \textbf{This file} respects an active narrow restriction (\hyperref[sec:org77ba459]{subtree narrowing} or \hyperref[sec:orgf0231a2]{Search's Narrow}), and the screen says so; narrow first, then export. A narrow applies only to the open document, never to Contacts Files.

\begin{center}
\begin{tabular}{lll}
Org source (flat) & vCard field & Notes\\[0pt]
\hline
Heading title & \texttt{FN} & \\[0pt]
\texttt{:N:} & \texttt{N} & Used as given; if absent, split from the title (last word as family name)\\[0pt]
\texttt{:EMAIL:} & \texttt{EMAIL} & \\[0pt]
\texttt{:EMAIL\_LABEL:} & \texttt{LABEL="..."} on the \texttt{EMAIL} line & When present\\[0pt]
\texttt{:PHONE:}, \texttt{:CELL:} & \texttt{TEL} & \\[0pt]
\texttt{:PHONE\_LABEL:} & \texttt{LABEL="..."} on the \texttt{TEL} line & \\[0pt]
\texttt{:WORK\_PHONE:} & \texttt{TEL;TYPE=WORK} & \\[0pt]
\texttt{:ADDRESS:}, \texttt{:ADR:} & \texttt{ADR} & Placed in the street component, since an address is one free-text string here\\[0pt]
\texttt{:ADDRESS\_LABEL:} & \texttt{LABEL="..."} on the \texttt{ADR} line & \\[0pt]
\texttt{:ADDRESS\_RAW:} & The structured \texttt{;;street;city;region;postal;country} part of \texttt{ADR} & The seven components, kept verbatim and used when present\\[0pt]
\texttt{:NICKNAME:} & \texttt{NICKNAME} & \\[0pt]
\texttt{:ORG:} & \texttt{ORG} & Several \texttt{;}-separated parts are joined into one string\\[0pt]
\texttt{:JOB\_TITLE:} & \texttt{TITLE} & \\[0pt]
\texttt{:URL:} & \texttt{URL} & \\[0pt]
\texttt{:URL\_LABEL:} & \texttt{LABEL="..."} on the \texttt{URL} line & \\[0pt]
\texttt{:PHOTO:} & \texttt{PHOTO} & An \texttt{http} or \texttt{https} URL as is; an embedded photo as a \texttt{data:} URI (to display photos, see \hyperref[sec:org75c5aa1]{\texttt{org-xx-startup-with-show-photos}})\\[0pt]
\texttt{:NOTE:} & \texttt{NOTE} & \\[0pt]
A heading's tags & \texttt{CATEGORIES} & \texttt{:CATEGORIES\_RAW:} is preferred when present\\[0pt]
\texttt{:CATEGORIES\_RAW:} & \texttt{CATEGORIES} & The original comma-separated text; tags can't hold spaces, so "Other Account" would otherwise become \texttt{Other\_Account}\\[0pt]
\texttt{org-contacts-birthday-property} (default \texttt{BIRTHDAY}) & \texttt{BDAY} & Accepts a bare \texttt{YYYY-MM-DD} with no description\\[0pt]
\texttt{:ANNIVERSARY:} & \texttt{ANNIVERSARY} (\texttt{YYYYMMDD}) & A vCard 4.0 property, written in the 3.0 card as a phone's own export does\\[0pt]
\texttt{:SPOUSE:} & \texttt{RELATED;TYPE=spouse;VALUE=text} & Likewise\\[0pt]
\end{tabular}
\end{center}

\begin{center}
\begin{tabular}{lll}
Tree style \texttt{:FIELDTYPE:} & vCard field & Notes\\[0pt]
\hline
\texttt{name} & The contact itself & Needs \texttt{:KIND: individual} on the same heading\\[0pt]
\texttt{n} & \texttt{N} & Used as given; else split from the title\\[0pt]
\texttt{cell} & \texttt{TEL;TYPE=CELL} & A \texttt{:LABEL:} property becomes a \texttt{LABEL="..."} parameter, here and in the rows below\\[0pt]
\texttt{phone}, \texttt{phone-work}, \texttt{phone-home} & \texttt{TEL}, \texttt{TEL;TYPE=WORK}, \texttt{TEL;TYPE=HOME} & \\[0pt]
\texttt{email}, \texttt{email-work}, \texttt{email-home} & \texttt{EMAIL}, \texttt{EMAIL;TYPE=WORK}, \texttt{EMAIL;TYPE=HOME} & \\[0pt]
\texttt{address} & \texttt{ADR} & Street component. A \texttt{:RAW:} property holds the seven preserved components, and a \texttt{\#+BEGIN\_VERSE} or \texttt{\#+BEGIN\_QUOTE} block in the body becomes the \texttt{LABEL}\\[0pt]
\texttt{address-work}, \texttt{address-home} & \texttt{ADR;TYPE=WORK}, \texttt{ADR;TYPE=HOME} & \\[0pt]
\texttt{nickname} & \texttt{NICKNAME} & \\[0pt]
\texttt{org} & \texttt{ORG} & \\[0pt]
\texttt{job-title} & \texttt{TITLE} & \\[0pt]
\texttt{url} & \texttt{URL} & \\[0pt]
\texttt{photo} & \texttt{PHOTO} & A URL is the heading's title. An embedded photo is titled "Photo", with its \texttt{data:} URI in a \texttt{:DATA:} property\\[0pt]
\texttt{note} & \texttt{NOTE} & Read from a \texttt{\#+BEGIN\_VERSE} or \texttt{\#+BEGIN\_QUOTE} block in the body, not the title\\[0pt]
\texttt{birthday} & \texttt{BDAY} & A bare \texttt{YYYY-MM-DD} is accepted\\[0pt]
\texttt{anniversary} & \texttt{ANNIVERSARY} & \\[0pt]
\texttt{spouse} & \texttt{RELATED;TYPE=spouse} & \\[0pt]
A contact's tags & \texttt{CATEGORIES} & \texttt{:CATEGORIES\_RAW:} on the contact heading is preferred\\[0pt]
\end{tabular}
\end{center}

Both styles cover the same set of fields, which is smaller than org-vcard's full mapping. \texttt{:PREFERRED:} is accepted but doesn't yet affect field order or add a \texttt{PREF} parameter.

\subsubsection{What is dropped}
\label{sec:org3fd2b56}

Property drawers and \texttt{ARCHIVE\_*} properties (see \hyperref[sec:orgf0de015]{Archiving}) are not exported. SCHEDULED and DEADLINE show as a small italic line under the heading in Markdown, HTML and ODT; Calendar is built from them.

\noindent\rule{\textwidth}{0.5pt}

\subsection{Import}
\label{sec:orgc97ce31}
\textbf{More → Import} lists the formats it reads, as Export does: \textbf{Contacts (.vcf)} and \textbf{iCalendar (.ics)}. Each opens its own panel. The palette commands are \texttt{org-vcard-import} and \texttt{icalendar-import-file}.

\textbf{Contacts (.vcf)} is the reverse of \hyperref[sec:org053e4f4]{Contacts (.vcf) export}. A \texttt{.vcf} file can hold one contact or many, and each becomes a new heading appended to the end of the open document. Nothing already in the file is touched.

\textbf{Style}: \textbf{Flat} or \textbf{Tree}, the two structures Export supports, independent of the style you last exported. Tree is the default, because Flat keeps only one value per field type (see below). Exporting a contact, importing that file and exporting again gives an equivalent vCard.

\textbf{Sources}: \textbf{Choose vCard file…}, \textbf{Choose a heading…} and \textbf{This file}. A document can hold headings whose body is vCard text, such as Export's \textbf{To: new buffer} output or text you paste in. \textbf{Choose a heading…} scans that heading and its descendants for bodies starting with \texttt{BEGIN:VCARD}, and \textbf{This file} scans the whole document. Both say plainly when there is nothing to import ("No vCard data found in this file"). The \textbf{To: new buffer} checkbox sends the imported contacts to a new unsaved \texttt{*scratch*} document instead of appending to the open one.

\textbf{What is read}: \texttt{FN} is the only required property, as in \href{https://www.rfc-editor.org/rfc/rfc6350}{RFC 6350}; a vCard block without one is skipped. \texttt{N}, \texttt{EMAIL}, \texttt{TEL}, \texttt{ADR}, \texttt{URL}, \texttt{NICKNAME}, \texttt{ORG}, \texttt{TITLE}, \texttt{PHOTO}, \texttt{NOTE}, \texttt{BDAY} and \texttt{CATEGORIES} are mapped when present, the reverse of Export's tables.

\begin{itemize}
\item \texttt{N} is kept exactly as given, never recomputed from \texttt{FN} (which would mangle a business name or a single name).
\item \texttt{ADR}'s components (PO Box, street, city, region, postal code, country) are joined into one comma-separated string.
\item A value with embedded newlines (a multi-line \texttt{NOTE} or \texttt{ADR}) becomes one line, with a space-semicolon-space separator, wherever it lands in a single-line Flat property.
\item \texttt{CATEGORIES} become org tags on the heading, and export writes them back as \texttt{CATEGORIES:}.
\item Folded (continued) lines are unfolded, and escapes (\texttt{\textbackslash{},}, \texttt{\textbackslash{};}, \texttt{\textbackslash{}\textbackslash{}}, \texttt{\textbackslash{}n}) are reversed, the inverse of Export's escaping.
\end{itemize}

\textbf{Apple and Google \texttt{itemN.} grouping}: these contacts tie a label to a property with a group prefix, such as \texttt{item1.TEL} with \texttt{item1.X-ABLabel} (standard vCard 3.0). This is always recognized, because otherwise the prefixed properties would be dropped. The label becomes \texttt{:EMAIL\_LABEL:}, \texttt{:PHONE\_LABEL:} or \texttt{:URL\_LABEL:} (Flat), or a \texttt{:LABEL:} property on the field heading (Tree), like \texttt{:ADDRESS\_LABEL:} for addresses. A label you assigned ("Landline") is kept as is. Apple's placeholder form (\texttt{"\_\$!<HomePage>!\$\_"}) is reduced to its inner name ("HomePage"), and an empty label or Apple's "no label" placeholder (\texttt{"\_\$!!\$\_"}) gives no label. \texttt{URL} can have several values. Apple's \texttt{X-ABADR} country-code companion is ignored.

\textbf{Clean data (Google, Apple)} is a checkbox, on by default, for non-standard quirks of those exports:

\begin{enumerate}
\item Google appends an eighth component to \texttt{ADR}: the address again, as a multi-line label. On, it is kept alongside the seven structured components, so a later export can rebuild the original. Off, it is ignored, as RFC 6350's seven-component \texttt{ADR} requires. The standard \texttt{ADR;LABEL="..."} parameter is recognized either way.
\item Google escapes colons as \texttt{\textbackslash{}:}, which isn't a real vCard escape. On, it is turned back into \texttt{:} everywhere; off, it is left as written.
\item Apple's \texttt{X-ABLabel} placeholders are interpreted as described above; off, they are stored as written.
\end{enumerate}

\textbf{Flat style keeps one value per type}: one \texttt{:EMAIL:}, \texttt{:PHONE:} or \texttt{:WORK\_PHONE:}, \texttt{:URL:} and \texttt{:ADDRESS:}. When a vCard has several of one type, the first is kept and the rest are dropped, so use Tree style to keep them all (org-vcard lists multiple \texttt{EMAIL} as an open item too). Labels survive as \texttt{:ADDRESS\_LABEL:}, \texttt{:EMAIL\_LABEL:}, \texttt{:PHONE\_LABEL:} and \texttt{:URL\_LABEL:}.

\textbf{Tree style has no such limit}: each value becomes its own descendant heading. The \texttt{address-work} and \texttt{address-home} fieldtypes (like \texttt{email-work} and \texttt{phone-home} in Export's table) round-trip a \texttt{TYPE=WORK} or \texttt{TYPE=HOME} address. A field whose value contains a newline (a note, always, or a Google address whose label spans lines) gets a fixed title ("Note", "Home Address") with its text as a \texttt{\#+BEGIN\_VERSE} block in the body. Verse keeps every line break, so separate facts (old phone numbers, street and city) don't blur together, and a body can hold newlines that Flat must flatten. Export reads \texttt{\#+BEGIN\_VERSE} or \texttt{\#+BEGIN\_QUOTE} and falls back to the heading's title when there is neither. A short single-line label is a plain \texttt{:LABEL:} property on the field heading.

\textbf{Warnings}: properties the parser doesn't know are collected, listed once by name in the status message after the import ("Imported 3 contacts from vCard. 2 unrecognized properties were skipped: X-SOCIALPROFILE, X-YAHOO."). \texttt{VERSION}, \texttt{PRODID}, \texttt{REV}, \texttt{UID} and \texttt{X-ABADR} are known and deliberately unmapped, so they aren't reported. A base64 \texttt{PHOTO} gets its own count in the same message, since it is kept as a \texttt{data:} URI (see Export's \texttt{:PHOTO:} row).

\texttt{ANNIVERSARY} and a \texttt{RELATED;TYPE=spouse} name become \texttt{:ANNIVERSARY:} and \texttt{:SPOUSE:}; other relations are reported as skipped. A \texttt{.vcf} file with no valid (\texttt{FN}) contacts gives a clear status message. A malformed line inside an otherwise valid contact is skipped without aborting the import.

\textbf{iCalendar (.ics)} is the reverse of \hyperref[sec:org053e4f4]{Calendar (.ics) export}, and reads a calendar exported from Google, Apple, Outlook or Nextcloud. Each event becomes a new heading appended to the end of the open document, and each task a TODO heading. Nothing already in the file is touched. The sources are the same as for Contacts: \textbf{Choose iCalendar file…}, \textbf{Choose a heading…} and \textbf{This file} (which look for bodies starting with \texttt{BEGIN:VCALENDAR}), and the \textbf{To: new buffer} checkbox, which has its own setting. There is no style or clean-up option.

\begin{itemize}
\item An event's time is a timestamp in the heading's title, where the agenda looks for one: \texttt{* Dentist <2026-10-05 Mon 09:00-10:30>}. An all-day event has a date only. A multi-day event is a range, \texttt{<2026-10-05 Mon>-{}-{}<2026-10-07 Wed>}, which the agenda shows on its first and last day.
\item Times are shown in this device's time zone. A time in UTC or in a named zone (\texttt{TZID}) is converted, as Emacs's \texttt{icalendar-import-file} does. A time with no zone is kept as written, and so is one in a zone name the browser doesn't know (Outlook's own names, for example); the status line says so.
\item \texttt{LOCATION}, \texttt{URL} and \texttt{UID} become properties, \texttt{CATEGORIES} become tags, \texttt{DESCRIPTION} becomes the body, and a cancelled event keeps \texttt{:STATUS: CANCELLED}.
\item A recurrence becomes a repeater when one says the same thing (\texttt{FREQ} with \texttt{INTERVAL}, such as \texttt{+1w}). A rule a repeater cannot express (an end date or \texttt{COUNT}, several weekdays, "the second Tuesday") is imported once, at its first date, with the rule in an \texttt{:RRULE:} property, and the status line counts them. Nothing becomes a repeater that would run on the wrong days or for ever.
\item A task (\texttt{VTODO}) becomes a \texttt{TODO} heading, or \texttt{DONE} with \texttt{CLOSED} once completed. \texttt{DUE} is its \texttt{DEADLINE}, \texttt{DTSTART} its \texttt{SCHEDULED}, and \texttt{PRIORITY} becomes \texttt{[\#A]}, \texttt{[\#B]} or \texttt{[\#C]}.
\item Attendees, organizer, attachments and exceptions (\texttt{EXDATE}) are not carried over, and the status line names them ("Imported 2 events and 1 task from iCalendar. 1 unrecognized property was skipped: ATTENDEE."). Alarms are ignored.
\end{itemize}

\noindent\rule{\textwidth}{0.5pt}

\subsection{Help}
\label{sec:orgb4d67ea}
Open Help from the command palette's \textbf{Help} (or god-mode's \texttt{h SPC m}, see \hyperref[sec:org996c4cc]{Keybindings}). It is this file, \texttt{README.org}, opened like any other document and read-only (see \hyperref[sec:org75c5aa1]{Local Variables}). It shows as its own tab (see \hyperref[sec:orgda9f5ef]{Tabs}), always last and visually distinct, and opening Help again switches to it if it is already open. It opens with \texttt{\#+STARTUP: show2levels}, as any document with that line does.

\noindent\rule{\textwidth}{0.5pt}

\subsection{Settings}
\label{sec:orgbb54b74}
Reached via \textbf{More (⋮) → Settings}, which replaces the outline until you navigate away, like switching views. Settings opens on a list of six categories, each a page of small cards: \textbf{Appearance}, \textbf{Editing and startup}, \textbf{Agenda and weather}, \textbf{Sync and storage}, \textbf{Export} and \textbf{Advanced}. At the foot of the list, \textbf{Check for updates} is centred on a line of its own, with the version centred beneath it; the result of the check (or “Out of date — reload”) appears on that line in bold red. Tap one to open it and \textbf{‹ Settings} to come back; a number on a category counts what you have set there: changed Quick Settings, connected services, ticked export items. \textbf{Search settings} finds a card from any category. In a window wide enough (about 720 px), the categories become a sidebar with the page beside it. Settings reopens on the page you left.

\textbf{Order of settings} (Advanced) moves categories, and the cards inside each, up or down with ▲ ▼; \textbf{Reset to default order} puts them back. The order is kept on the device and travels with \textbf{Export Settings}.

The cards, by what they hold:

\begin{itemize}
\item \textbf{Quick Settings}: labeled controls for the Global and Local Variables the app supports (see \hyperref[sec:org75c5aa1]{Local Variables}), grouped by area, with an \textbf{Advanced} group for those with more complex syntax (refile targets, agenda files, the \texttt{org-xx-*} menu variables; see \hyperref[sec:org08ac3df]{Extras menu} and \hyperref[sec:orgfca7ef3]{Menu customization}). It shares storage with the Global Variables text box, so each reflects changes made in the other.
\item \textbf{Global Variables} and \textbf{Capture Templates}: raw \texttt{name: value} lines and raw JSON, each edited in a popup.
\item \textbf{GitHub} and \textbf{WebDAV}: connection details for each backend. Use a fine-grained GitHub token scoped to just that repository (Contents read and write only), and for WebDAV an app-specific password if your server offers one.
\item \textbf{Calendar (CalDAV) / Contacts (CardDAV)}: the calendar your agenda is mirrored to and the address book your contacts are, right after WebDAV (see \hyperref[sec:org66ee6b6]{Calendar and contacts sync}).
\item \textbf{Agenda notifications}: reminders before appointments, just before GitHub (see \hyperref[sec:org0153874]{Agenda notifications}).
\item \textbf{Appearance}: theme (System, Light or Dark), fonts, independent font sizes for main text, tables and menus, reading width, custom colors for Light and Dark, and \textbf{Paragraph Spacing}. Paragraph Spacing is two independent values laid out like the font sizes: one for paragraphs and one, \textbf{Tables}, for tables and other secondary blocks (source and quote blocks, rules). Each is 0 to 32px with a 10px default. Where a paragraph meets a table or block the larger value applies, and headings and list items space themselves.
\item \textbf{Backup}: \textbf{Export Settings} bundles everything on this page into one downloadable JSON file, and \textbf{Import Settings} merges a previously exported file back in, leaving anything it doesn't mention untouched.
\item \textbf{Updates}: shows the version of the code that is running and checks for a new one. If a newer version is already installed but this page was opened before it took over, a banner and a note here say so; \textbf{Reload} brings the page up to date. The installed Android app has no button, since a newer version is installed instead.
\item \textbf{Export} (General, ASCII, HTML): one card per export type. \textbf{General} sets a \textbf{Page header} and \textbf{Page footer} for printed pages (\texttt{Draft | | \%d} is left, centre, right). \textbf{HTML} sets a \textbf{Preamble}, \textbf{Postamble} and \textbf{Head}. Each has a tick box and is off until ticked; a ticked item is used for every export whose file does not set its own, and the file always wins. A \textbf{Format codes} list in each card explains \texttt{\%t}, \texttt{\%a}, \texttt{\%d}, \texttt{\%p}, \texttt{\%P} and the rest.
\end{itemize}

A file's own \texttt{\#+STARTUP:} or Local Variables block still overrides an app-wide default set here.

\noindent\rule{\textwidth}{0.5pt}

\subsection{Offline behavior and sync}
\label{sec:org57d263c}
Edits are cached locally at once; writing to disk, GitHub or WebDAV happens only on an explicit Save.

\textbf{Conflict detection}: Save compares the file's current state with what the app last knew, so one version never silently overwrites the other. The app also checks for external changes when a file is opened, when the tab regains focus, and when you switch into a tab with real backing storage. A banner then appears with \textbf{Merge}, \textbf{Reload} and \textbf{Dismiss}. For a local file this background check never asks the browser for permission: it runs only if access is already granted, so the banner may not appear until you have saved or reloaded the file in the session. Save, Reload and Merge still ask when they need to.

\textbf{Resolving a conflict}: the app keeps the full text of the file as of the last successful open or save, the common ancestor a three-way merge needs. When Save finds the file changed elsewhere under your unsaved edits, it compares the ancestor, your version and the other one line by line:

\begin{itemize}
\item \textbf{The other version has no real change, or matches yours}: Save just proceeds. A server that changes a file's ETag without changing its content isn't a conflict.
\item \textbf{Your changes and theirs are in different places}: they are combined automatically and Save completes with a status message (\texttt{Saved — merged with the changes made elsewhere}). The merge is its own undo step (\textbf{Merged with changes from elsewhere}), so Undo returns to your text from just before it.
\item \textbf{Both sides changed the same lines}, or adjacent lines with no unchanged line between them (the same place, as in git and GNU diff3): a dialog lists each such region with your version and the other one. Choose \textbf{Yours}, \textbf{Other} or \textbf{Both} (yours, then theirs) for each, or \textbf{All yours} or \textbf{All other} for every region at once. \textbf{Apply} stays disabled until every region has a decision, and everything outside the regions is already combined. \textbf{Cancel}, Escape or tapping outside changes nothing: the other version is untouched, your edits stay unsaved, and the next Save raises the same conflict.
\item \textbf{No common ancestor is known} (a file last synced by a version that didn't keep one, until its next open or save), or the two versions are too different to compare by lines: every difference, or the whole file as one choice, is shown for you to decide.
\end{itemize}

\textbf{The banner's Merge} does the same without waiting for Save. It combines your unsaved edits with the version now stored, asks about clashing regions in the same dialog, and applies the result as an ordinary, undoable, still-unsaved edit. The other version becomes the new baseline, so the Save that follows is a plain write. With no unsaved changes Merge just reloads, and a read-only buffer can't be merged into (make it writable first).

\textbf{What this isn't}: the merge works on lines, not org structure. A subtree moved on one side and edited on the other is an ordinary clash, and two sides that each add the same new heading in different places keep both. A clean automatic merge is still worth a glance; that is what the undo step is for.

\noindent\rule{\textwidth}{0.5pt}

\subsection{Platform support}
\label{sec:orgc3c9c80}
Local file access ("Local", a live writable handle) requires the File System Access API -- Chrome/Edge, desktop or Android. Not supported on iOS, a WebKit limitation.

On unsupported platforms, \textbf{Import} replaces "Local": pick a file, edit it, and Save downloads the new version for you to move into place yourself. GitHub and WebDAV work the same everywhere, including iOS.

\textbf{Sharing into the app}: Android's Share sheet and the icon's \textbf{Capture} shortcut work for the installed app. iOS has no equivalent, so use the clipboard (\texttt{\%x}); see \hyperref[sec:orge0ecaf9]{Capturing from other apps}.

\subsubsection{Android app}
\label{sec:org3d41625}
The Android app is the same app in a native wrapper, with a few things done the Android way.

\begin{itemize}
\item \textbf{Local files}: \textbf{File → Open → Local file} uses Android's file picker, and the app keeps access to the file, so Save writes straight back and the file reopens from the recent list. \textbf{Save As → Local file} names a new one. The picker starts in \texttt{Documents/org-pwa}, in the shared Documents folder where a file manager finds it, and you can choose any other folder. A file is known by its name, so two files with the same name from different folders replace each other.
\item \textbf{Agenda files}: \texttt{local:name.org} names a file in the org-pwa folder (below), like \texttt{github:journal.org;local:home.org}. It loads silently and read-only; a file you opened with \textbf{Open → Local file} works too.
\item \textbf{Attaching a file}: \textbf{Attach} opens Android's chooser with the camera apps beside your files, as the browser's does. Pick \textbf{Take a photo} or \textbf{Record a video}, and the result is attached.
\item \textbf{Saving a file out}: exports and settings backups go straight into the \textbf{Downloads} folder, and the status line says so.
\item \textbf{The org-pwa folder}: the first time something needs it (a local file's first attachment, or a \texttt{local:} name that isn't found), the app asks you to choose a folder, starting at \texttt{Documents/org-pwa}. Attachments go under \texttt{data/} inside it, as on GitHub and WebDAV, and \texttt{local:} files are found in it by name, anywhere a \texttt{github:} or \texttt{webdav:} file can be named: \texttt{org-agenda-files}, \texttt{org-contacts-files}, \texttt{org-refile-targets}, a capture template's \texttt{file}, a link or \texttt{\#+INCLUDE}. \texttt{local:contacts.org} is \texttt{contacts.org} in that folder, and \texttt{local:areas/home.org} is \texttt{home.org} in its \texttt{areas} folder. A file you opened with the file picker is found by its name first. \textbf{Choose the org-pwa folder} changes it.
\item \textbf{Reminders}: agenda notifications (see \hyperref[sec:org0153874]{Agenda notifications}) are handed to Android, so they arrive with the app closed. Allow notifications when asked; without the \textbf{Alarms \& reminders} setting they may be a few minutes late.
\item \textbf{Long-press the app icon}: your capture templates are listed, in the order they are in \textbf{Settings → Capture Templates}, as many as the launcher shows. Tapping one opens that template. Drag one onto the home screen for an icon, or run \textbf{Add a capture icon to the home screen}.
\item \textbf{Sharing in}: share a link or text from any app to \textbf{org-pwa}; Capture opens with it (see \hyperref[sec:orge0ecaf9]{Capturing from other apps}).
\item \textbf{Separate from the browser}: the app keeps its own settings, tokens and documents, so enter GitHub, WebDAV, calendar and contacts settings again. A WebDAV, CalDAV or CardDAV server must allow the origin \texttt{https://localhost} (CORS).
\item \textbf{Updates}: the app is updated by installing a newer APK over the old one. \textbf{Settings → Updates} shows which version you have.
\end{itemize}

\textbf{Screen size}: mobile-first below 900px wide. Above that, Settings and Help open in a resizable side panel next to the outline instead of replacing it, and the app uses the full window width. Reading width is unlimited by default; Settings → Appearance has an opt-in narrower column.

No auto-capitalization anywhere in this app.

\noindent\rule{\textwidth}{0.5pt}

\subsection{Command palette}
\label{sec:orgc687081}
A searchable list of what the app can do, the touch counterpart of Emacs's \texttt{M-x}. Open it from \textbf{More → Commands} (first by default) or, from a keyboard, in \hyperref[sec:orgb59a778]{god-mode} with \texttt{g x} (which is \texttt{M-x}) or \texttt{h SPC a} (\texttt{C-h a}).

\textbf{Searching}: type a few letters or words. Each word must match the command's name, the real Org or Emacs command it corresponds to, or one of a few extra keywords (so \texttt{deadline} finds \textbf{Edit details}). Case doesn't matter. Real names work as you'd type them in \texttt{M-x}: \texttt{org-clock-in} finds \textbf{Clock in} but not \textbf{Clock out}, because a hyphenated word is matched as that exact phrase, while \texttt{clock in} typed with a space is two words that can match anywhere. A category name (\textbf{Clocking}, \textbf{Heading}, \textbf{Export}) also matches at the start of its words. There is no fuzzy or skip-letter matching.

\textbf{Order}: commands that can run now come first, then the rest, dimmed, each saying why it can't (\textbf{nothing is narrowed}, \textbf{no clock is running}, \textbf{the buffer is read-only}, \textbf{tap a heading first}). A command is never mysteriously absent, and choosing a dimmed one repeats the reason as a status message. With nothing typed, your most recently used commands come first (marked \texttt{↺}), remembered across sessions. ↑ and ↓ move the highlight, Enter or a tap runs it, and Escape or a tap outside closes it.

\textbf{What it acts on}: heading commands work on the heading whose action menu is open (tap a heading's title first), else the keyboard-focused heading. A command that needs a heading and can't find one says so, except \textbf{Add heading after} and \textbf{Add TODO heading after}, which also work with nothing selected (see the table in \hyperref[sec:org996c4cc]{Keybindings} for where the new heading lands).

\textbf{Names}: a command's real Emacs or Org function name is shown beside it and is searchable. The same name, quoted, names the command in an Extras-menu entry (see \hyperref[sec:org08ac3df]{Extras menu}), so the palette, the god-mode chords and Extras share one vocabulary. The commands with no real name are this app's own: \textbf{New document}, \textbf{Undo history}, \textbf{Sync agenda to calendar}, \textbf{Rebuild calendar from the agenda}, \textbf{Show display measurements}, \textbf{Add a capture icon to the home screen} and \textbf{Choose the org-pwa folder}. A command with a \hyperref[sec:orgb59a778]{god-mode} key runs the same code as that key and shows it.

\textbf{What's in it}, by group:

\begin{center}
\begin{tabular}{ll}
Group & Commands\\[0pt]
\hline
\textbf{Heading} & cycle TODO state, set priority, set effort estimate, edit details (schedule, deadline, timestamps, tags, properties), add heading or TODO heading after, promote, demote, move up, move down, cut, paste, archive, unarchive, refile, attachments, narrow, widen\\[0pt]
\textbf{Clocking} & clock in, clock out, cancel clock, continue last clock, go to clocked task, go to a recently clocked task\\[0pt]
\textbf{Document} & save, save as, open, new, undo, redo, undo history, toggle read-only, capture, recalculate all tables\\[0pt]
\textbf{View} & agenda, TODO list, outline view, text view, cycle visibility of the whole document, calendar, search, sync agenda to calendar, rebuild calendar, sync contacts to the address book, rebuild contacts\\[0pt]
\textbf{Export} & the Export menu; this file as HTML, Markdown, ASCII, ODT or an Org buffer\\[0pt]
\textbf{Import} & Contacts (.vcf), \texttt{org-vcard-import}; iCalendar (.ics), \texttt{icalendar-import-file}\\[0pt]
\textbf{App} & settings, help, show display measurements, add a capture icon to the home screen and choose the org-pwa folder (Android app only)\\[0pt]
\end{tabular}
\end{center}

\textbf{Also in it}: your own capture templates and \hyperref[sec:org08ac3df]{Extras menu} entries, one command each. \textbf{Capture: <description>} (such as \textbf{Capture: Meeting}) runs that template directly, and \textbf{Extras: <label>} runs that Extras entry as tapping it would. Both are read fresh each time the palette opens, so an edited template list or a document's own Extras menu (set in its Local Variables) is current. The plain \textbf{Capture} command still opens the template picker.

\textbf{Commands}, like the other More entries, can be renamed, hidden or reordered with \hyperref[sec:orgfca7ef3]{menu customization} (\texttt{org-xx-menu-aliases}). It is the ordinary \texttt{more:} label \texttt{Commands}, which renders first by default.


\subsection{Security and privacy}
\label{sec:org74c50ec}
There is no server of ours: the app talks only to the places you set up, and nothing is sent anywhere else.

\textbf{Where things are kept.} Your files stay where they are (a file on the device, GitHub or WebDAV). The app keeps its settings in the browser's own storage on this device, including the GitHub token, the WebDAV, calendar and contacts passwords, and the init script. They are stored as plain text, as in any browser app, so use a token limited to one repository and an app-specific password, and sign out of a shared device by clearing its site data. Nothing but what you set up is synced, exported or read from a file.

\textbf{What can run.} Only the app's own code runs in the page. The page's Content-Security-Policy refuses inline and injected scripts. Code from a document or the init script runs in a \hyperref[sec:org34d11ec]{sandbox} with no access to the page, your files or your tokens. It reaches a server only if you allow it.

\subsection{Keybindings}
\label{sec:org996c4cc}
A practical subset for a keyboard or mouse. On a phone or tablet it applies only to a paired Bluetooth keyboard.

These are \textbf{not org's own bindings}, on purpose. Org leans on the \texttt{C-c} prefix, and \texttt{C-c}, \texttt{C-v}, \texttt{C-a} and \texttt{C-f} are reserved by browsers for copy, paste, select-all and find, so a page can't reliably take them; org's \texttt{M-←} and \texttt{M-→} collide with the browser's back and forward. Where org's key is safe it is used as is; otherwise a safe key stands in. god-mode (below) reaches org's fuller set.

\begin{center}
\begin{tabular}{lll}
Key & Action & Matches real org?\\[0pt]
\hline
\texttt{↓}, \texttt{↑} & Move keyboard focus to the next or previous visible line: any row (heading, paragraph, table, list item, block). Moving down into a folded heading reveals its first body row. A status message reports the first or last line & No equivalent\\[0pt]
\texttt{S-↓}, \texttt{S-↑}, \texttt{S-←}, \texttt{S-→} & Only with focus on a table: move cell by cell (the header row counts; a rule between rows doesn't). Outside a table they aren't intercepted. After you edit a cell by tap, \texttt{Enter} commits and \texttt{Escape} cancels (the \textbf{⎌} button does the same by tap), and the cell stays highlighted for more Shift+arrow navigation. Clicking anything that isn't a table cell clears it & No equivalent\\[0pt]
\texttt{i} & Edit the current line (heading title, paragraph, table cell or list item), cursor at the start & Like vim's \texttt{i}\\[0pt]
\texttt{j}, \texttt{k} & Move focus to the next or previous visible heading, faster than \texttt{↓} and \texttt{↑}. A status message reports the first or last heading & Vim-style\\[0pt]
\texttt{Tab} & Cycle the focused heading's fold state & Yes, \texttt{org-cycle}\\[0pt]
\texttt{Enter} & Open or close the focused heading's action row & No equivalent\\[0pt]
\texttt{/} & Focus the search box & A web convention\\[0pt]
\texttt{Escape} & The first that applies: cancel a god-mode sequence, exit god-mode, clear keyboard focus, enter god-mode & A web convention, extended for god-mode\\[0pt]
\end{tabular}
\end{center}

Keyboard focus, a thin highlight around the focused row, is separate from tapping a heading open. It starts unset; the first arrow, \texttt{j} or \texttt{k} picks the first or last visible heading, and shortcuts that need the focused heading do nothing until one exists. \texttt{Tab} follows the same archived-heading rule as the swipe gesture (see \hyperref[sec:org75c5aa1]{Local Variables}). Shortcuts never fire while you are typing in a text field, and Ctrl and Cmd combinations are never intercepted, so the browser's copy, paste and find stay safe.

Cycling TODO state, moving a heading, promoting and demoting, and inserting a heading are available by tap, in god-mode (\texttt{c t}, \texttt{g↑} and \texttt{g↓}, \texttt{g←} and \texttt{g→}, \texttt{gRET}), and from the command palette (\textbf{Add heading after}).

\subsubsection{god-mode}
\label{sec:orgb59a778}
An alternate input mode that reaches org's fuller keybinding set (\texttt{C-c C-t}, \texttt{M-RET} and the rest) without the browser ever seeing a held Ctrl or Alt. \texttt{Escape} turns it on (see the priority order above), or on a phone the \textbf{[g]} button does (\hyperref[sec:org7b13dfa]{below}). While it is on, every key is part of a sequence instead of a plain shortcut, and the minibuffer (see \hyperref[sec:org6cade45]{Modeline and minibuffer}) echoes the sequence.

God-mode starts on the heading you last used, or else on the row at the top of the screen. Running a block with \texttt{C-c C-c} keeps the cursor on it.

When a sequence completes, the minibuffer names what ran, using the command palette's own label: \texttt{God-mode · C-c C-x C-j → Go to clocked task}. A chord with no palette command (\texttt{Tab}, the arrows, \texttt{C-f}) shows just the chord. The name stays until the next keystroke. Entering god-mode focuses the first visible heading if none is focused, since its actions need one.

While a sequence is in progress, a card under the top bar lists the keys that would continue it and what each does, as Emacs's which-key does. Each entry is a button: tap it to send that key, so a whole chord can be completed by tapping, with no device keyboard. A key that leads on to more is marked \texttt{…}. \texttt{Escape} cancels the sequence and the card goes away.

Opening an overlay panel (the Extras menu, Settings, Docs, Capture, Search, the File, View and More menus, the History panel or the Calendar) exits god-mode, and the plain shortcuts don't act while a panel is open; \texttt{Escape} closes the open panel first. Help opened with \texttt{h SPC m} is the exception (see below).

The translation rules follow real god-mode.el (github.com/emacsorphanage/god-mode), plus the \texttt{c} rule, which is this app's own:

\begin{itemize}
\item A plain key means Control: \texttt{t} is \texttt{C-t}.
\item \texttt{g} prefixes the next key with Meta (\texttt{g↑} is \texttt{M-↑}), so a bare \texttt{C-g} can't be typed. \texttt{G} (Shift+g) prefixes Control+Meta (\texttt{Gn} is \texttt{C-M-n}).
\item \texttt{SPC} toggles literal mode: following keys are typed with no modifier until \texttt{SPC} is pressed again. \texttt{c SPC g} is \texttt{C-c g}, not \texttt{C-c C-g}.
\item A named key (arrow, Tab, Return) typed bare is never controlled, so ordinary movement keeps working.
\item \texttt{c} as the first key gives \texttt{C-c} and starts an org-style chain: further letters are controlled too (\texttt{c t} is \texttt{C-c C-t}, \texttt{c c} is \texttt{C-c C-c}) while punctuation stays literal (\texttt{C-c .}, \texttt{C-c |}). One \texttt{c} is enough for any \texttt{C-c C-<letter>} command.
\end{itemize}

\begin{center}
\begin{tabular}{lll}
god-mode sequence & Emacs chord & Action\\[0pt]
\hline
\texttt{↓}, \texttt{↑} & \texttt{<down>}, \texttt{<up>} & Next or previous visible line; inside a table, the next or previous row\\[0pt]
\texttt{←}, \texttt{→} & \texttt{<left>}, \texttt{<right>} & Collapse or expand the heading; inside a table, the previous or next column\\[0pt]
\texttt{S-↓}, \texttt{S-↑}, \texttt{S-←}, \texttt{S-→} (the floating \textbf{S} then an arrow, or Shift+arrow) & \texttt{S-<down>}, \texttt{S-<up>}, \texttt{S-<left>}, \texttt{S-<right>} & On a focused heading: move it down or up among its siblings, promote it (left) or demote it (right), like the heading's action buttons. In a table they move between cells. With nothing focused, \texttt{S-↓} and \texttt{S-↑} go to the next or previous heading. On a laptop, \textasciitilde{}Alt\textasciitilde{}+arrow moves, promotes and demotes the heading from anywhere, even in a table\\[0pt]
\texttt{S-←}, \texttt{S-→} & \texttt{S-<left>}, \texttt{S-<right>} & Table cell to the left or right\\[0pt]
\texttt{i} & (not a chord) & Edit the current line, cursor at the start (see \texttt{a} and \texttt{e})\\[0pt]
\texttt{a} & \texttt{C-a} & Before \texttt{i}: put the cursor at the start of the line\\[0pt]
\texttt{e} & \texttt{C-e} & Before \texttt{i}: put the cursor at the end of the line\\[0pt]
\texttt{Tab} & \texttt{Tab} & Cycle the focused heading's fold state\\[0pt]
\texttt{S-Tab} & \texttt{S-Tab} & Cycle the whole document (overview, fully shown)\\[0pt]
\texttt{←}, \texttt{→} & \texttt{<left>}, \texttt{<right>} & Collapse the focused heading fully, or expand it one level\\[0pt]
\texttt{g↑}, \texttt{g↓} & \texttt{M-↑}, \texttt{M-↓} & Move the focused heading up or down among its siblings\\[0pt]
\texttt{g←}, \texttt{g→} & \texttt{M-←}, \texttt{M-→} & Promote or demote the focused heading with its subtree\\[0pt]
\texttt{gS-←}, \texttt{gS-→} & \texttt{M-S-←}, \texttt{M-S-→} & The same (there is no separate whole-subtree action)\\[0pt]
\texttt{c f}, \texttt{c b} & \texttt{C-c C-f}, \texttt{C-c C-b} & Focus the next or previous heading at the same level; a status message says when there is none\\[0pt]
\texttt{c u} & \texttt{C-c C-u} & Focus the parent heading (a status message at the top level)\\[0pt]
\texttt{c n}, \texttt{c p} & \texttt{C-c C-n}, \texttt{C-c C-p} & Focus the next or previous visible heading, at any level\\[0pt]
\texttt{c SPC l} & \texttt{C-c l} & Copy an \texttt{[[id:…]]} link to the focused heading (\texttt{org-store-link})\\[0pt]
\texttt{c x w} & \texttt{C-c C-x C-w} & Cut the focused heading and its subtree to the clipboard (\texttt{org-cut-subtree})\\[0pt]
\texttt{c x y} & \texttt{C-c C-x C-y} & Paste a subtree after the focused heading at its level, shifting the pasted levels to match (\texttt{org-paste-subtree})\\[0pt]
\texttt{x SPC n s} & \texttt{C-x n s} & \hyperref[sec:org77ba459]{Narrow} to the focused heading's subtree\\[0pt]
\texttt{x SPC n w} & \texttt{C-x n w} & \hyperref[sec:org77ba459]{Widen}: remove any narrowing\\[0pt]
\texttt{x SPC x r} & \texttt{C-x x r} & \textbf{Rename buffer…} (\texttt{rename-buffer}): ask for a new name in the minibuffer. An empty name or one another buffer has is refused; the file is not touched\\[0pt]
\texttt{gRET} & \texttt{M-RET} & Insert a heading after the focused one, as its sibling. With nothing focused, an empty document gets a heading at the top and a document with headings gets one at the bottom\\[0pt]
\texttt{gS-RET} & \texttt{M-S-RET} & The same, with a TODO keyword\\[0pt]
\texttt{c t} & \texttt{C-c C-t} & Cycle the focused heading's TODO state; with no state and more than one \texttt{\#+TODO:} workflow, open a state picker\\[0pt]
\texttt{c ,} & \texttt{C-c ,} & Cycle priority (none, A, B, C)\\[0pt]
\texttt{c v} & \texttt{C-c C-v} & Switch to the TODO view\\[0pt]
\texttt{c .}, \texttt{c !}, \texttt{c d}, \texttt{c s} & \texttt{C-c .}, \texttt{C-c !}, \texttt{C-c C-d}, \texttt{C-c C-s} & Open the details editor (SCHEDULED, DEADLINE, timestamps)\\[0pt]
\texttt{c w} & \texttt{C-c C-w} & Refile the focused heading (\texttt{org-refile}, see \hyperref[sec:org6433005]{Refile})\\[0pt]
\texttt{c a} & \texttt{C-c C-a} & Open the attachment choices (\texttt{org-attach}, see \hyperref[sec:org7b27204]{Attachments})\\[0pt]
\texttt{c SPC a} & \texttt{C-c a} & Open the Agenda view\\[0pt]
\texttt{c SPC c} & \texttt{C-c c} & Open Capture, which exits god-mode (see below)\\[0pt]
\texttt{h SPC m} & \texttt{C-h m} & Open Help (\texttt{describe-mode}) without exiting god-mode\\[0pt]
\texttt{s} & \texttt{C-s} & Open search\\[0pt]
\texttt{x q} & \texttt{C-x C-q} & Toggle \texttt{buffer-read-only} (see \hyperref[sec:org75c5aa1]{Local Variables})\\[0pt]
\texttt{x s} & \texttt{C-x C-s} & Save (\texttt{save-buffer})\\[0pt]
\texttt{x SPC s} & \texttt{C-x s} & Offer to save every modified tab that has a file, one key each (\texttt{save-some-buffers})\\[0pt]
\texttt{x w} & \texttt{C-x C-w} & Save under a new name or place (\texttt{write-file}, Save As)\\[0pt]
\texttt{x f} & \texttt{C-x C-f} & Open a file (\texttt{find-file})\\[0pt]
\texttt{c x SPC e} & \texttt{C-c C-x e} & Set the focused heading's effort estimate (\texttt{org-set-effort}, see \hyperref[sec:orgb9e46b3]{Clocking})\\[0pt]
\texttt{c x i} & \texttt{C-c C-x C-i} & Clock in (\texttt{org-clock-in})\\[0pt]
\texttt{c x o} & \texttt{C-c C-x C-o} & Clock out (\texttt{org-clock-out})\\[0pt]
\texttt{c x q} & \texttt{C-c C-x C-q} & Cancel the running clock without recording its time (\texttt{org-clock-cancel})\\[0pt]
\texttt{c x x} & \texttt{C-c C-x C-x} & Resume the most recently clocked heading (\texttt{org-clock-in-last})\\[0pt]
\texttt{c x j} & \texttt{C-c C-x C-j} & Jump to the clocked-in heading, in whichever tab it runs (\texttt{org-clock-goto})\\[0pt]
\texttt{u c x j} & \texttt{C-u C-c C-x C-j} & Choose from recently clocked tasks and jump to it (\texttt{org-clock-select-task})\\[0pt]
\texttt{c e} & \texttt{C-c C-e} & Open the export dispatcher (\texttt{org-export-dispatch}, see \hyperref[sec:org053e4f4]{Export})\\[0pt]
\texttt{f} & \texttt{C-f} & Arm the next \texttt{/} to redo instead of undo\\[0pt]
\texttt{/} & \texttt{C-/} & Undo, or redo if armed\\[0pt]
\texttt{g x} & \texttt{M-x} & Open the \hyperref[sec:orgc687081]{command palette}\\[0pt]
\texttt{h SPC a} & \texttt{C-h a} & The same (\texttt{apropos-command}: find the commands that match)\\[0pt]
\end{tabular}
\end{center}

Some valid sequences have no equivalent yet and say so instead of doing nothing: \texttt{c c} (toggling a checkbox or running a code block needs finer focus than a heading), \texttt{c |} and the table row and column inserts (use the +row and +col controls under any table), and \texttt{c l} and \texttt{c o} (there is no link insertion or cursor-based link opening).

\textbf{Pickers}: \texttt{c SPC c}, and \texttt{c t} on a heading with no TODO state and several workflows, exit god-mode and show a picker. Its options are tappable buttons, each also labelled with a one-character hotkey you can type at once. A hotkey is shown and accepted only when no other option shares it; those options are tap-only. Opened any other way (a tap, the ⋮ Capture button, a TODO-less heading's badge), the pickers show no hotkeys. \texttt{Escape} cancels either.

\textbf{Help}: \texttt{h SPC m} doesn't exit god-mode. While Help is the only thing open, god-mode keeps working underneath it, and \texttt{Escape} cancels a sequence and then exits god-mode instead of closing Help. Opening any other panel, or leaving god-mode, restores the ordinary rule.

\textbf{Redo}: Emacs has no redo command; it undoes the undo, after something breaks the chain (god-mode users conventionally press \texttt{C-f}). Here \texttt{f} just arms the next \texttt{/} to redo. \texttt{/} keeps redoing while armed, any other key disarms it, and with nothing left to redo \texttt{/} falls back to undo. The palette's \textbf{Undo} and \textbf{Redo} are unaffected.

\textbf{\texttt{i}, \texttt{a} and \texttt{e}} aren't chords; they are loosely modeled on vim and act only as the first key of a fresh sequence (\texttt{g i} is still \texttt{M-i} and \texttt{c i} still \texttt{C-c C-i}). \texttt{a} and \texttt{e} set where the cursor lands the next time \texttt{i} opens the editor, and say so in the minibuffer (\texttt{C-a → edit at start (press i)}). Once editing, every key is ordinary text-field input.

\noindent\rule{\textwidth}{0.5pt}

\subsubsection{God-mode without a hardware keyboard}
\label{sec:org7b13dfa}
A phone has no \texttt{Escape} key, so a floating \textbf{[g]} button, the third floating button, stacked above Extras or in its place when there is no Extras menu, does three things together: it enters god-mode, shows a floating keyboard and brings up the device's keyboard. Tapping it again ends all three, and it is drawn solid while the keyboard is showing. It rides above the mode line and the device keyboard, and nothing already on screen moves for it. Settings → Quick Settings → \textbf{Keyboard} has a \textbf{Show the floating [g] god-mode button} checkbox (\texttt{org-xx-god-mode-button}, on by default) for devices with a real \texttt{Escape} key. A document's Local Variables line overrides it, and turning it off while the keyboard is showing ends god-mode.

The floating keyboard has the keys a phone keyboard can't type, in one row: \textbf{⌨}, \texttt{Tab}, \texttt{←}, \texttt{→}, \texttt{↑}, \texttt{↓}, \texttt{S} and \texttt{g}. Its thick accent-coloured border makes it easy to spot, and the keys are 34px wide and 44px tall so the row fits a 360px screen.

\begin{itemize}
\item \textbf{S} shifts the next key (\texttt{Tab} becomes \texttt{S-Tab}, \texttt{g} becomes \texttt{G}, the Control+Meta prefix). After an arrow it stays on, so \texttt{S} then \texttt{↑} \texttt{↑} \texttt{↑} moves the focused heading up three places; tap \textbf{S} again to turn it off. It also shifts the next key typed on the phone keyboard, so \texttt{gS-RET} (\texttt{M-S-RET}, a TODO heading) is: tap \texttt{g}, tap \textbf{S}, tap Return. A letter typed after \textbf{S} becomes its capital.
\item \textbf{⌨} hides or shows the device keyboard and changes nothing else. It is lit while the keyboard is up. After a popup closes, the keyboard doesn't return by itself; tap ⌨.
\item The move area (⠿) drags the panel, and its position is remembered. Tap it without dragging to minimize the panel to its handle, and tap the handle to restore it.
\end{itemize}

\textbf{Placement}: the panel sits above the mode line and the floating buttons, so the device keyboard never hides it and it never covers the mode line or \textbf{[g]}. By default it sits just above the floating buttons. If the keyboard comes up under it, it is pushed up and stays there when the keyboard goes away; only dragging lowers it.

\textbf{Lifetime}: god-mode and the panel end through \textbf{[g]}, \texttt{Escape} or another panel opening, never because the device keyboard went away. Hide the keyboard however you like and god-mode, the panel and its keys carry on, leaving the whole screen to read the outline. Entering god-mode with a real \texttt{Escape} never shows the panel.

If a letter doesn't register on your keyboard app, the command palette reaches the same commands by name.


\subsection{Differences from Emacs org-mode}
\label{sec:org2d2d41c}
Read this if you know org-mode well and want to know where org-pwa is a subset, a simplification or just different. It lists only the deltas: a topic not mentioned here matches org's behavior.

\textbf{Interaction model} (see \hyperref[sec:org996c4cc]{Keybindings}): touch-first, not keyboard-driven. The \hyperref[sec:orgc687081]{command palette} is \texttt{M-x} only in spirit: a fixed, curated list that you can't add to, not every Emacs command. god-mode's table doesn't cover table cells, per-list-item checkboxes or link insertion (see \hyperref[sec:orgb59a778]{god-mode}).

\textbf{Folding} (see \hyperref[sec:org22c076e]{Editing your outline}): no equivalent of Emacs's richer visibility state machine, but three flags per heading, on top of which \texttt{overview}, \texttt{content}, \texttt{showall}, \texttt{showeverything} and \texttt{show2levels} to \texttt{show5levels} are implemented.

\textbf{Priority, tags, properties} (see \hyperref[sec:org22c076e]{Editing your outline}): tags and property keys complete from what is used in the open document and loaded agenda files (tags also from a \texttt{\#+TAGS:} line). Priority is a fixed A/B/C/None row with nothing to complete, and nothing is treated specially beyond \texttt{CUSTOM\_ID} and the \texttt{ARCHIVE\_*} set.

\textbf{Inheritance} (see \hyperref[sec:org75c5aa1]{Local Variables}): property inheritance is a simple boolean. Tag inheritance takes \texttt{t}, \texttt{nil}, a list or a regexp, and \texttt{org-tags-exclude-from-inheritance}. Both apply to search, to the Agenda and TODO tag filters, and to the dimmed inherited tags shown after a heading's own tags (switched by \texttt{org-xx-show-inherited-tags} in the outline and \texttt{org-agenda-show-inherited-tags} in the Agenda).

\textbf{Tables} (see \hyperref[sec:org1458240]{Table formulas}): \texttt{\#+TBLFM:} covers arithmetic, conditionals, named constants, scalar and aggregate functions, date and time arithmetic and formatting, and unit conversion. Complex numbers and basic vectors and matrices are covered; not reproduced: the rest of Emacs Calc (symbolic algebra), named-field references, and recalculation on every keystroke.

\textbf{Footnotes} (see \hyperref[sec:orga27d7d4]{Footnote commands}): in the outline a footnote goes in through the Footnote button of the paragraph editor, and there is no \texttt{org-ref} integration.

\textbf{Clocking} (see \hyperref[sec:orgb9e46b3]{Clocking}): no custom \texttt{org-duration-units}, no \texttt{Effort\_ALL+}, and no clock-in reminder beyond the modeline's over-estimate marker.

\textbf{Agenda} (see \hyperref[sec:org105cb3a]{Agenda}): effort sorting, filtering and the per-day limit are session settings in the ⏱ panel, not \texttt{org-agenda-sorting-strategy} or \texttt{org-agenda-max-effort}. Date stamps in titles are hidden unless \texttt{org-xx-agenda-hide-title-stamps} is \texttt{nil} (Emacs strips only a time of day that the row already shows, \texttt{org-agenda-remove-times-when-in-prefix}). The tag, state and category chips are not Emacs's \texttt{/}, \texttt{<} and \texttt{=} filters: they match a heading's own tags only. There is no negated (\texttt{-<}) filter, no effort column format (\texttt{\%e}) and no effort sort in the TODO view. The sexp evaluator supports a closed, documented set of functions (including \texttt{org-class}, \texttt{org-date} and \texttt{and}, \texttt{or}, \texttt{not}), not a Lisp interpreter: no \texttt{let}, \texttt{date} variable or \texttt{calendar-*} helpers, no \texttt{diary-remind}, and no holiday-skipping in \texttt{org-class}. \texttt{diary-anniversary}, \texttt{diary-cyclic}, \texttt{diary-block} and \texttt{diary-date} are left out on purpose, since their \texttt{org-*} twins take unambiguous year-month-day dates (see \hyperref[sec:org36f3daa]{General sexp timestamps}). A \texttt{local:} agenda file is read-only, and only a file already opened on this device can be named.

\textbf{Babel}: \texttt{js} blocks run, with \texttt{:results}, \texttt{:var}, \texttt{:timeout}, \texttt{:wrap}, \texttt{:net}, \texttt{:location}, \texttt{:cache}, \texttt{:noweb} and tangling (see \hyperref[sec:org41ed666]{Source blocks} and \hyperref[sec:org229bca3]{Literate programming}). These do not exist: blocks in any other language, \texttt{:session}, running one block from inside another (\texttt{<{}<{}name()>{}>{}}), and results that refresh when a file is opened.

\textbf{Links and structure}: \texttt{[[id:...]]} links are followed through the open tabs and Agenda Files, then the files opened here before, from a table kept in this browser rather than \texttt{org-id-locations}. \texttt{<{}<{}targets>{}>{}} and radio targets work in the outline and in every export. \texttt{\#+SETUPFILE:} is read from the same place as the document, or from an http(s) URL, and kept for offline use. \textbf{Citations} (see \hyperref[sec:org261744f]{Citations}): only the \texttt{basic} processor. \texttt{\#+CITE\_EXPORT:} naming \texttt{csl}, \texttt{biblatex} or \texttt{natbib} is exported as \texttt{basic}, \texttt{note} styles print as \texttt{text}, and \texttt{org-cite-export-processors} is not a setting (use \texttt{\#+CITE\_EXPORT:}).

\textbf{Undo/redo} (see \hyperref[sec:org9a920ce]{Undo/redo and History}): no undo tree, and fold state isn't preserved across an undo or redo jump (the file's \texttt{\#+STARTUP} visibility is reapplied). There is no per-keystroke granularity; the browser's field-level undo covers that, much as Emacs batches consecutive typing.

\textbf{Search} (see \hyperref[sec:orgf0231a2]{Searching}): Search and Replace both reach every file in \texttt{org-agenda-files}, as org's \texttt{org-search-view} and \texttt{tags-query-replace} do. Both subtree narrowing (see \hyperref[sec:org77ba459]{Narrowing}) and Search's Narrow survive a reload.

\textbf{Export} (see \hyperref[sec:org053e4f4]{Export}): no \LaTeX{} export. ASCII export matches \texttt{ox-ascii} for characters, entities, blocks and lists, but lays out headings, tables and links its own way (numbered headings with indented bodies, tables with \texttt{|} borders, links as \texttt{description (target)}, no \texttt{Notes} list of links). Justify-left and justify-right blocks are set as ordinary text, which is what Emacs 29.3 produces too. HTML export differs from \texttt{ox-html} in a few deliberate ways: its own colouring instead of \texttt{htmlize}, KaTeX instead of MathJax, readable element ids, real checkboxes, \texttt{::} inside code or verbatim text is not a description-list separator, and no \texttt{\#+HTML\_DOCTYPE:}.

\textbf{Refile} (see \hyperref[sec:org6433005]{Refile}): you can't create a new destination heading by typing a name that doesn't exist yet.

\textbf{Attachments} (see \hyperref[sec:org7b27204]{Attachments}): the same folder model as \texttt{org-attach} (an \texttt{ID} or \texttt{DIR} folder and the \texttt{:ATTACH:} tag) and the dispatcher keys \texttt{a} \texttt{b} \texttt{n} \texttt{o} \texttt{f} \texttt{d} \texttt{D} \texttt{s} \texttt{S} \texttt{z}, plus \texttt{r} to record audio. The Folder panel stands in for \texttt{org-attach-dired}: select files, then open, save or delete them. Not built: copy, move and link (\texttt{c}, \texttt{m}, \texttt{l}), attaching from a Dired buffer, and \texttt{DIR} paths outside the storage. Needs GitHub, WebDAV, or the org-pwa folder in the Android app.

\textbf{Conflict resolution} (see \hyperref[sec:org57d263c]{Offline behavior and sync}): a line-based three-way merge with a region-by-region resolver. It merges lines, not org structure, so a subtree moved on one side and edited on the other is an ordinary clash, and a region has no side-by-side or word-level diff.

\noindent\rule{\textwidth}{0.5pt}

\subsection{Known limitations / not built yet}
\label{sec:org52c5ce1}
Gaps against org-mode itself are in \hyperref[sec:org2d2d41c]{Differences from Emacs org-mode}. This section covers app-specific and platform limits:

\begin{itemize}
\item Local and relative images show a placeholder on local files (File System Access) and iOS imports, the same picker-permission limit as archiving and capture to a file. GitHub and WebDAV documents are unaffected.
\item No File System Access on iOS, a platform limit rather than a bug (see \hyperref[sec:orgc3c9c80]{Platform support}).
\item A read-only buffer accepts no edit anywhere in the document, though Save and Save As still work, as in Emacs (see \hyperref[sec:org996c4cc]{Keybindings}).
\end{itemize}

\noindent\rule{\textwidth}{0.5pt}

\subsection{Development}
\label{sec:org45678a6}
A static site: \texttt{index.html} and ES modules, with no build step to run or deploy it. Serve it any way you like, including \texttt{python3 -m http.server} or GitHub Pages.

The one exception is \texttt{src/vendor/mathjs/mathjs-custom.min.js}, a tree-shaken custom mathjs build (Unit, BigNumber and Fraction, used by \hyperref[sec:org1458240]{Table formulas}'s \texttt{uconv()}, \texttt{F} and \texttt{pN}). It is regenerated occasionally with \texttt{tools/build-mathjs-vendor.sh}, a dev-only step; the committed output is what ships, as with the vendored KaTeX files.

\textbf{Tests}: the engine (\texttt{src/}) and the browser adapters (\texttt{src-browser/}) use Node's built-in test runner, with no external dependencies:

\begin{verbatim}
npm test
\end{verbatim}

The tests run with a frozen clock, so they give the same result on any day. \texttt{npm run test:zones} runs them in five time zones.

Behaviour that Emacs defines is held to Emacs itself. \texttt{tools/gen-repeat-oracle.mjs} and \texttt{tools/gen-agenda-oracle.mjs} ask Emacs 29.3 (with \texttt{faketime}) what completing a repeating task does and which days the agenda lists, and store the answers under \texttt{test/fixtures/}; the tests compare the app with them.

\texttt{npm run lint} (ESLint) and \texttt{npm run typecheck} (\texttt{tsc} on files that begin with \texttt{// @ts-check}) need \texttt{npm install -{}-{}no-save eslint @eslint/js globals typescript}. \texttt{node tools/release.mjs N} checks, bumps the cache name and writes \texttt{org-pwa-vN.zip}; \texttt{.github/workflows/ci.yml} runs the same checks and the smoke test on every push.

\texttt{node tools/roundtrip-corpus.mjs DIR} checks that every \texttt{.org} file in a directory comes back byte for byte after an open and save with no edit.

One test checks that every module the app imports is in the service worker's precache list, because nothing is cached at runtime and a missing module would stop the app starting offline.

\textbf{Code layout}: \texttt{app.js} is the entry point (about 1,400 lines): it imports everything, sets up the shared state, wires event listeners and runs \texttt{bootstrap}. Features live in \texttt{src-browser/}, one module per area. All runtime state is one object, \texttt{S}, in \texttt{app-state.js}, and the page's DOM elements are looked up once in \texttt{dom.js}. The split is by file rather than decoupled (modules call each other, so the import graph has cycles), and browser modules need a DOM, so Node can't import them. \texttt{test/module-structure.test.js} and \texttt{tools/check-scope.mjs}, a scope analysis that needs the acorn and eslint-scope parsers installed once (\texttt{npm install -{}-{}prefix /tmp/acorn acorn eslint-scope}), keep the structure from decaying.

Everything the app asks of the device (files on the device, saving a file out, the clipboard, the service worker) goes through \texttt{src-browser/platform.js}, which has browser defaults. A native shell supplies its own by setting \texttt{window.orgPwaPlatform} before the app loads, and hands a share to \texttt{window.orgPwaLaunch()}. The browser tests use a fake platform the same way.

\textbf{Browser smoke test}: it covers what unit tests can't reach, since \texttt{app.js} is UI wiring:

\begin{verbatim}
node tools/smoke.mjs             # everything
node tools/smoke.mjs conflict    # only checks whose name contains "conflict"
SMOKE_SLICE=0:27 node tools/smoke.mjs    # only the first 27 checks; 27: runs the rest
\end{verbatim}

It drives the real app in headless Chromium through flows that have broken before, against a built-in static server and a small in-memory WebDAV server. It is dev-only, adds nothing to \texttt{package.json}, needs Playwright and a Chromium once (\texttt{npm install -g playwright \&\& npx playwright install chromium}, or set \texttt{CHROMIUM\_PATH}), and exits non-zero on any failure.

\texttt{test/testing.org} holds an older record of the testing approach and is no longer maintained.
\end{document}
